\documentclass[12pt, a4paper]{article}

% Packages
\usepackage[utf8]{inputenc}
\usepackage[T1]{fontenc}
\usepackage{amsmath, amssymb, amsthm}
\usepackage{mathtools}
\usepackage{geometry}
\usepackage{setspace}
\usepackage{booktabs}
\usepackage{array}
\usepackage{bm}
\usepackage{multirow}
\usepackage{caption}
\usepackage{graphicx}
\usepackage{tikz}
\usetikzlibrary{arrows.meta, positioning, calc}
\usepackage{hyperref}
\usepackage[round]{natbib}
\usepackage{xcolor}
\usepackage{longtable}
\usepackage{float}
\usepackage{tcolorbox}
\newcolumntype{P}[1]{>{\raggedright\arraybackslash}p{#1}}
\usepackage{pdflscape}
\usepackage{float}
\usepackage{booktabs}
\usepackage{graphicx}
\usepackage{threeparttable}

\geometry{margin=2.5cm}
\onehalfspacing

% Custom commands
\newcommand{\Sk}{S_k}
\newcommand{\Fk}{F_k}

\title{\textbf{The Pandemic Trilemma} \\[1em] 
\large Work in Progress. Please do not cite or circulate.}

% Der Trick: \setcounter{footnote}{1} schiebt den Zähler unsichtbar vor, 
% sodass \thanks direkt als zweites Symbol das Kreuz wählt.
\author{Aulis Pesenti\setcounter{footnote}{1}\thanks{University of Basel, Wirtschaftswissenschaftliche Fakultät, Peter Merian-Weg 6, 4002 Basel, Switzerland; (e-mail: aulis.pesenti@unibas.ch).} \\ University of Basel}

\date{\today}

\begin{document}
\maketitle

\begin{abstract}

\end{abstract}
% =============================================================
% Section 1 — Introduction
% =============================================================

\newpage

\section{Introduction}

Pandemics are not standard recessions. In a standard downturn, stabilization policy typically aims to support demand and income. In a pandemic, by contrast, governments must use containment to reduce disease pressure, even though this same instrument suppresses economic activity. Fiscal policy can compensate for the output loss, but only by creating public liabilities. The policy problem is therefore not a conventional stabilization problem, but a dynamic trade-off between mortality, output, and public debt. 

COVID-19 made this trade-off visible at global scale. What had appeared in late 2019 as a marginal news item turned, within weeks, into a worldwide shock. By March 2020, governments reacted with stay-at-home orders, business closures, restrictions on movement, and severe limits on international travel. Trade, mobility, and everyday economic activity were disrupted on a scale not seen since World War II \citep{OECD2020EconomicOutlook, WorldBank2020GEP}, with excess mortality rising sharply and output gaps reaching up to -25 percentage points in the most affected countries.

Since the situation was unprecedented, there was no common policy template for governments. Countries reacted with highly diverse combinations of containment measures and especially fiscal support. Understanding these choices requires a framework that captures the joint constraints between health, output, and public debt.

This paper formulates this problem as a trilemma. A government confronting a pandemic faces three simultaneous objectives: protecting public health by minimizing excess mortality, preserving economic stability by minimizing the deviation of output from potential, and maintaining fiscal sustainability by avoiding destabilizing increases in public debt. These objectives are mutually constrained by the structural architecture of the pandemic environment, forming a \textit{trilemma} in the economic sense of the Mundell-Fleming impossible trinity \citep{mundell1963}: no feasible policy of a benevolent social planner can simultaneously minimize mortality, output losses, and debt accumulation.

The impossibility arises from four key properties of the pandemic environment, which I formalize in Section~\ref{sec:trilemma}. First, once introduced exogenously, the viral prevalence of a lethal pathogen grows endogenously through human contact. Second, in the absence of an effective vaccine, the main available suppression instruments---lockdowns, social distancing mandates, mobility restrictions, and other non-pharmaceutical interventions (NPIs)---protect public health by reducing human interaction, but this same mechanism contracts economic activity. Third, fiscal instruments can compensate for the resulting output loss, but their effects are heterogeneous and they do so at the cost of public debt accumulation. Fourth, even without discretionary fiscal intervention, the output contraction itself erodes the fiscal position through automatic stabilizers: falling incomes reduce tax revenues and trigger transfer payments, generating debt independently of discretionary policy choices. Given this structure, the planner's decision is not whether to intervene fiscally, but how to compose the intervention across instruments with fundamentally different transmission channels. Demand injection---for instance, direct transfers and cash payments---supports aggregate expenditure during containment and generates a demand-side output effect. Capacity preservation---such as wage subsidies that maintain employment relationships, loans, and guarantees that prevent firm insolvency---limits the persistence of the contraction by preventing the destruction of organizational capital, supplier networks, and job-specific human capital. These instruments differ in their mechanisms, fiscal costs, and dynamic profiles. The optimal fiscal response therefore depends not only on the aggregate size of support, but also on its composition and timing. As the pandemic state evolves, this composition should adjust accordingly, turning the social planner's problem into an intertemporal dynamic optimization problem.

%%Noch genauer und anpassen
A broad literature has formalized important parts of this problem. On the macroeconomic modelling side, \citet{Atkeson2020} introduce the baseline SIR structure into a macroeconomic setting, \citet{alvarez2020} solve for the optimal lockdown trajectory of a planner trading off fatalities against output costs, and \citet{eichenbaum2021macro} endogenize the behavioral recession by embedding epidemiological dynamics into macroeconomic equilibrium. Related work studies targeted containment policies \citep{acemoglu2021} and endogenous labor-supply responses to infection risk \citep{Jones2021}. On the empirical side, a large literature estimates the effect of non-pharmaceutical interventions on economic activity across countries \citep{Ashraf2022, Konig2020, Cross2020, Gordon2021}, documenting the contractionary impact of containment but generally treating fiscal policy as absent, given, or secondary.

A second strand studies the fiscal response to COVID-19 more directly. \citet{chetty2020} show that U.S.\ direct transfers like demand injection had limited spending effects during containment; \citet{Autor2022} estimate the employment effects of the Paycheck Protection Program; \citet{Pappa2022} analyze the effects of on-budget fiscal measures on growth in Europe; and \citet{fariaecastro2021} show in a nonlinear DSGE model that the pandemic shock changes the ranking of fiscal multipliers, with liquidity assistance dominating transfers for employment stabilization. The most comprehensive cross-country evidence is provided by \citet{deb2021}, who find that what they call emergency-lifeline instruments outperform income-support measures during strict containment, while the ranking reverses upon reopening.


The closest contribution to my approach is \citet{KaplanMollViolante2020}, who integrate SIR dynamics into a heterogeneous-agent model with both lockdown and fiscal transfers. Their analysis constructs a pandemic possibility frontier and shows how fiscal transfers affect the distributional incidence of the health--economy trade-off. Yet fiscal policy enters as undifferentiated lump-sum transfers, while government debt appears as a residual accounting outcome rather than as a state variable in the planner's objective. As a result, the compositional and intertemporal fiscal margins remain outside of the optimization problem.

Two gaps emerge from this literature. First, existing work primarily studies a health--output dilemma, while fiscal sustainability is rarely treated as a co-equal objective in cross-country comparison. Second, fiscal policy is often represented as an aggregate instrument, although pandemic fiscal packages differed sharply in their transmission channels, timing, and budgetary pass-throughs.

Therefore, the central research question is: \textit{What is the model-implied optimal composition and timing of fiscal support between demand injection and capacity preservation over the course of a pandemic and what are the costs---in foregone output stabilization and excess debt---of deviating from this benchmark?}

To answer this question, I constructed a measure-level database of pandemic fiscal interventions: 2,301 individual measures across all 38 OECD economies, recorded at daily frequency with their announcement date, size, and primary transmission channel. This granularity allows me to distinguish demand injection, capacity preservation, and health-related spending at the country-quarter level---a decomposition that aggregate fiscal data cannot deliver.

The paper makes four contributions. First, it develops a theoretical framework that formalizes pandemic policy as a dynamic trilemma between excess mortality, the output gap, and public debt. Second, it introduces a transmission-mechanism-based classification of pandemic fiscal measures. Rather than treating fiscal support as a homogeneous aggregate, I distinguish between demand injection, capacity preservation, and health-related spending. Third, I use the measure-level database to parameterize the instrument-specific transition channels. Fourth, I use the estimated transition system to solve a finite-horizon planner problem. This converts the empirical transition parameters into a model-implied benchmark and allows me to compare observed country strategies to this benchmark.

Taken together, these contributions provide a retrospective assessment of the highly heterogeneous fiscal strategies adopted during COVID-19. The paper evaluates which types of support stabilized output, at what fiscal cost, and how close observed country strategies came to a model-implied benchmark. This assessment is relevant beyond COVID-19 because the pandemic differed fundamentally from standard demand-driven recessions. Understanding which instruments worked under these conditions can inform fiscal policy in future crises characterized by similar supply-side or mobility restrictions. It also matters for current debt debates, as rising inflation and borrowing costs after the main pandemic period made pandemic-era debt more costly to service and reduced fiscal space.

%\begin{figure}[H]
    %\centering
    %\includegraphics[width=\textwidth]{Files/text/main/outcomes_y_cumdebt_excess_aer.pdf}
    %\caption{Pandemic Dynamics: Infection Prevalence, Excess Mortality, 
    %and Containment}
    %\label{fig:planing_horizon}
%\end{figure}

The analysis yields four sets of findings. First, the trilemma is a structural property of the pandemic environment, not an assumption. Within the estimated transition system, public debt is dynamically unstable and cannot be stabilized by any available instrument, containing the epidemic requires stringency high enough to depress output, and no fiscal instrument is self-financing, every instrument leaves a positive net debt footprint. Each stabilizing action on one dimension therefore feeds the instability of another.

Second, every leg of the trilemma is active in the data. The panel estimates identify all three trade-off channels jointly and with the expected signs: containment reduces infection transmission but depresses output; the output losses feed into public debt; and every fiscal instrument stabilizes output only at a positive net debt cost. Beyond their existence, the transmission channels differ sharply, and one of them is state-dependent in a way that proves central: the output effect of demand injection declines with containment and turns negative once stringency exceeds a threshold that the average observed OECD containment path only barely undercut. Capacity-preserving support---above-the-line
programs and, above all, liquidity instruments such as loans and guarantees, operating through persistent stocks---retains its effectiveness under containment. Pandemic fiscal policy is thus not a question of aggregate
size: the same unit of support has fundamentally different output effects depending on the instrument and the containment state in which it is deployed.

Third, judged against the model-implied benchmark, the observed average OECD response was strictly Pareto-dominated: holding the other two dimensions at their observed levels, a feasible policy path---using no instrument, scale, or adjustment speed beyond what OECD governments demonstrably implemented---would have reduced the discounted output loss by 79 percent, excess mortality by 22 percent, or terminal debt by 6.3 percentage points of GDP. The source of the gains is composition and timing rather than volume: the dominating policies front-load fiscal support, shift it toward effective capacity preservation, retime containment into the acute waves, and abandon
demand injection entirely. 

%%%5NEU SCHREIBEN
[CROSS COUNTRY FINDINGS UND GENERELL]
%Fourth, the same pattern holds country by country: no OECD economy operated on its own frontier,  The countries closest to their own frontier---Mexico, Sweden, Portugal, Finland, France, and Norway---ran front-loaded, liquidity-oriented packages with little demand injection or, as Mexico, deployed so little that almost nothing was left to convert. The furthest---the United States, New Zealand, Japan, Australia, and Denmark---were far from it in one of two ways: large, back-loaded, demand-injection-heavy programs whose scale bought little at the margin, or well-composed packages, as in Denmark and Japan, whose realized debt paths left far more attainable at the same terminal debt. The mechanism is visible in the realized data: the output gain per point of fiscally induced debt correlates at \(-0.89\) with the demand-injection share, and the four countries with a negative realized return are precisely those with the highest such shares. Italy is the instructive exception: hit first and hardest, it tightened containment faster than the common policy menu admits, so its observed strategy cannot be replicated---let alone dominated---within it; under a relaxed ramping limit, Italy would rank among the five countries closest to their frontier, consistent with its early, liquidity-led capacity-preservation response.

\noindent The paper proceeds as follows. Section~\ref{sec:trilemma} establishes the formal framework of the pandemic trilemma, introduces the relevant state variables, controls, and transition dynamics, and demonstrates the structural uniqueness. Section~\ref{sec:data} introduces the measure-level database and the construction of the key variables. Section~\ref{sec:estimation} estimates the transition parameters, and Section~\ref{subsec:calibration} calibrates and validates the resulting system. Section~\ref{sec:planner} formulates the planner's problem, and
Section~\ref{sec:results} presents the normative results: the planner benchmark, the constrained policy frontier, and the cross-country assessment. Section~\ref{sec:conclusion} concludes.

\newpage

\section{The Pandemic Trilemma}
\label{sec:trilemma}

%This section formalizes the trilemma by identifying four constitutive properties of the pandemic environment that made it unique, develops the theoretical structure with reference to the literature, and introduces the variables and parameters that enter the framework in Section~\ref{sec:planner}.

%% ============================================================
\subsection{Constitutive Properties}
\label{subsec:properties}
%% ============================================================


%% --- Property 1 ---
\subsubsection*{Property 1: Exogenous Shock with Endogenous Dynamics}


A pandemic begins with the introduction of a novel pathogen into the population. I denote the pandemic pressure at time $k$ by $\theta_k$, which I refer to as the viral prevalence. The initial outbreak constitutes an exogenous shock---it originates outside the economic system and is not the result of endogenous market dynamics or policy failure. I represent this shock by a stochastic disturbance $\varepsilon_k$.

The defining characteristic of a pandemic shock is its endogenous growth dynamic: in the absence of intervention, each infected individual transmits the virus, causing the number of infections to grow over time. I capture this by a persistence parameter $\rho_\theta > 1$ (similar to the logic of the reproduction rate during the pandemic), implying that viral prevalence grows exponentially if left unchecked \citep{Atkeson2020, alvarez2021, eichenbaum2021macro}.\footnote{The linear specification abstracts from the susceptible-depletion mechanism of the full SIR framework, where the effective reproduction rate declines endogenously as the susceptible population shrinks. This linearization is appropriate for the model horizon (2020--2021)} It is the lethality of the pathogen that transforms this explosive dynamic into a policy imperative: a non-lethal pathogen with $\rho_\theta > 1$ would spread but impose no mortality cost of sufficient magnitude to warrant contractionary public health interventions \citep{Jones2021, Barro2020}, thereby establishing the necessity of active policy intervention at the cost of economic contraction (Property~2).

This property distinguishes the pandemic shock qualitatively from an impulse shock---where the disturbance dissipates naturally over time ($\rho < 1$) and the system returns to equilibrium without active intervention \citep{jorda2022longer}. These potential consequences are captured by the state variable: \textit{excess mortality} ($d_k$), defined as the deviation of the periodic death rate from its historical seasonal average encompassing both direct COVID-19 deaths and indirect deaths attributable to avoidance of care and healthcare system overload \citep{Msemburi2023WHO, owid_em}. I denote the parameter governing the transmission from infections to deaths by $\delta_\theta$. Crucially, the biological lag between infection and death implies that mortality in period $k+1$ is already determined by the current viral prevalence $\theta_k$ and thus predetermined from the planner's perspective. Consequently, the planner's intervention $S_k$---which suppresses transmission to $\theta_{k+1}$---only impacts mortality starting from period $k+2$. This relationship is given by:

\begin{equation}
    d_{k+2} = \delta_\theta \, \theta_{k+1}, \qquad \text{with} \quad \theta_{k+1} = \rho_\theta \theta_k+\varepsilon_k
    \label{eq:mortality_intro}
\end{equation}

The interaction between the lethality parameter $\delta_\theta$ and the transmission potential $\rho_\theta$ plays a decisive role. If $\delta_\theta = 0$---that is, if the virus were non-lethal---or if the pathogen were not transmissible ($\rho_\theta = 0$), there would be no reason for the planner to suppress $\theta_{k+1}$. This logic is formally derived in Appendix \ref{app:Theta import}, which proves that the necessity of a trade-off-inducing policy ($S_k > 0$) depends strictly on the joint condition of positive lethality ($\delta_\theta > 0$) and active transmission ($\rho_\theta > 0$). 

%\medskip
%\noindent\textit{Summary.} Property~1 introduces the epidemiological state: the infection rate $\theta_k$ (self-reinforcing with $\rho_\theta > 1$), the mortality outcome $d_k$ (driven by $\delta_\theta$), and the exogenous shock $\varepsilon_k$. Together, they establish that the planner must act to suppress viral prevalence, as the system is unstable without intervention.


%% --- Property 2 ---
\subsubsection*{Property 2: Suppression Instrument}
Given the necessity of intervention established by Property~1, the planner's primary tools to suppress viral transmission during a pandemic are non-pharmaceutical interventions (NPIs): measures including business closures, restrictions on gatherings, school closings, restrictions on internal movement as well as international travel controls and stay-at-home orders. I denote the intensity of these measures at time $k$ by $\Sk$, where $\Sk = 0$ represents the absence of restrictions and $\Sk = 100$ represents a complete lockdown in every aspect. 

As NPIs suppress viral prevalence by reducing human contact \citep{Flaxman2020}, I denote the effectiveness of this suppression by the parameter $\phi_S > 0$. Crucially, NPIs do not remove a fixed number of infections from the population; rather, they reduce the rate at which existing infections generate new ones. Each unit of suppression $S_k$ reduces the transmission potential $\rho_\theta$ by the fraction $\phi_S$, yielding an effective reproduction rate $\rho_\theta^{\text{eff}}(S_k) = \rho_\theta(1 - \phi_S S_k)$. Combining this with the mortality transmission from Property~1, the full two-period channel from policy action to health outcome is:
%
\begin{equation}
    d_{k+2} = \delta_\theta \cdot \underbrace{\bigl[\rho_\theta (1 - \phi_S S_k) \, \theta_k + \varepsilon_k\bigr]}_{\theta_{k+1}} 
    \label{eq:excess dynamics}
\end{equation}
%
This equation encodes two features that are central to the planner's problem. First, the suppression effect is multiplicative in the current level of prevalence: at high $\theta_k$, each unit of lockdown prevents more infections---and therefore more deaths---than at low $\theta_k$. The marginal health benefit of suppression, $\partial d_{k+2} / \partial S_k = -\delta_\theta \rho_\theta \phi_S \theta_k$, is state-dependent and vanishes as the pandemic recedes (see Appendix \ref{app:Theta import}). Second, the planner does not control the infection level or the death rate directly but only the growth rate of the epidemic. When $\phi_S S_k$ is sufficient to push $\rho_\theta^{\text{eff}}$ below unity, infections decline; otherwise, they continue to grow despite the intervention. This threshold property---$\rho_\theta^{\text{eff}} < 1$ requires $S_k > \frac{1 - 1/\rho_\theta}{\phi_S}$---implies a minimum lockdown intensity below which suppression is ineffective at reversing the epidemic trajectory.

The central feature of Property~2 is that the channel through which NPIs suppress viral transmission simultaneously reduces the volume of economic activity. Containment measures close businesses, idle labor, and curtail the flow of goods, services, and persons that constitutes aggregate economic output. This duality is inherent to the nature of any instrument that operates through contact reduction.

To formalize this, I introduce the next state variable \textit{output gap} ($y_k$): the percentage deviation of real GDP from its potential trend, where $y_k = 0$ represents normal economic conditions and $y_k < 0$ represents a recession. The contractionary effect of NPIs on output is governed by the parameter $\alpha_S > 0$. Every unit of suppression $\Sk > 0$ therefore simultaneously implies some values of:
%
\begin{alignat}{2}
&\text{Epidemiological benefit:} &\qquad \Delta\theta &= -\phi_S \, \Sk < 0 \label{eq:S_health}\\[4pt]
&\text{Economic cost:} &\qquad \Delta y &= -\alpha_S \, \Sk < 0 \label{eq:S_econ}
\end{alignat}
%

The suppression instrument cannot control the epidemiological shock without simultaneously affecting the economic state. Beyond the direct contractionary effect of mandated containment, the pandemic generates an additional output cost through the behavioral response. When individuals observe rising death tolls, they voluntarily reduce economic activity independently of any government-mandated restrictions \citep{GOOLSBEE2021} making it a do or die strategy for the social planner. I capture this channel through the parameter $\beta_d < 0$, which governs the output response to the mortality flow ($d_k$) and is commonly referred to as the \textit{fear term}. 

A natural objection to Property~2 is that alternative instruments, such as testing, might control the pandemic without suppressing economic activity, thereby dissolving the contractionary link. However, testing capacity constitutes a complement to, rather than a substitute for, NPIs \citep{repec:nbr}. Critically, testing does not alter the biological parameters driving the pandemic: it reduces neither the inherent transmissibility ($\rho_\theta$) nor the lethality ($\delta_\theta$) of the virus. Instead, it operates exclusively through the suppression channel $S_k$ by identifying infected agents. Ideally, this improves the ratio of epidemiological benefit ($\phi_S$) to economic cost ($\alpha_S$) by replacing broad lockdowns with specific quarantines. However, testing cannot eliminate the contractionary effect ($\alpha_S > 0$) since even targeted isolation disrupts economic transactions. Moreover, the net effect on the efficiency ratio is ambiguous as higher detection rates may intensify responses.
%Thus, testing modulates the parameter ratio but does not resolve the structurally contractionary nature of suppression.

Vaccination, by contrast, is the instrument that dissolves the trilemma by invalidating Property~2. It operates through both channels, reducing viral transmission (lowering $\rho_\theta$) and reducing the lethality of breakthrough infections (lowering $\delta_\theta$), which, as shown above, is the critical mechanism. Crucially, it achieves it without operating through the suppression channel $S_k$---and thus, no restriction of economic activity \citep{Giordano2021V}, thereby eliminating the contractionary link. However, effective vaccines were not widely available until mid-2021, and the critical vaccination thresholds required to sustainably decouple mortality from viral prevalence were not reached in most OECD economies until the end of 2021 \citep{OECD2021}. Consequently, for the main phase of the pandemic (Q1 2020--Q4 2021)---the period during which critical policy decisions were taken \citep{Hale2021}---NPIs remained the only available suppression instrument, and were thus inherently contractionary.

Figure~\ref{fig:planing_horizon} illustrates on the OECD average this mechanism. During Waves~1--3, infection prevalence and excess mortality move in parallel, and governments maintain elevated containment. With the onset of Wave~4 (Omicron, late 2021), the relationship breaks: infection prevalence reaches its pandemic peak, yet excess mortality declines sharply---reflecting the reduction in $\delta_\theta$ achieved by widespread vaccination. Governments respond by withdrawing containment to near-zero levels despite record case counts. This decoupling is precisely the condition $\delta_\theta \to 0$ formalized above: once the fatality rate  falls sufficiently, the rational to act vanishe, and the trilemma ceases to bind. This observation defines the planner's horizon: the trilemma is operative from Q1\,2020 through Q4\,2021. Beyond the terminal date, the planner deals exclusively with the fiscal legacy of the crisis.


\begin{figure}[H]
    \centering
    \includegraphics[width=0.9\textwidth]{fig_theta_excess_stringency.pdf}
    \caption{Pandemic Dynamics: Infection Prevalence, Excess Mortality, 
    and Containment}
    \label{fig:planing_horizon}
\end{figure}

%\medskip
%\noindent\textit{Summary.} Property~2 introduces the suppression instrument $S_k$, its epidemiological effectiveness $\phi_S$, its direct economic cost $\alpha_S$, the behavioral cost $\beta_d$ driven by observed mortality, and the output gap $y_k$. It establishes that pandemic suppression is inherently contractionary: the planner cannot protect health without damaging the economy unless an effective vaccine is available.

%================================================================
% Property 3: Fiscal Compensation Instruments
% Revised: 20.05.2026
% ================================================================

\subsubsection*{Property 3: Fiscal Compensation Instruments}

Given that pandemic suppression is inherently contractionary (Property~2), the social planner requires compensating instruments to stabilize the economic state. The pandemic damages the economy through two main channels. Containment and behavioral retrenchment directly reduce the level of output by restricting economic activity and suppressing expenditure. Simultaneously, the prolonged disruption erodes the economy's capacity to recover---through firm failures, dissolution of employment relationships or disruption of supply chains---prolonging the contraction beyond the duration of the shock itself. To address these two channels, the social planner has to target these channels separately.

The first category, I frame it as demand injection ($F^{DI}_k$). Those comprises measures that directly increase the flow of disposable income in the economy (level): stimulus checks (such as the U.S.\ Economic Impact Payments under the CARES Act), direct transfers to households (such as the Australian JobSeeker supplement or the Japanese universal cash transfer of \textyen100,000 per resident), and tax reductions that raise disposable income (such as the German temporary VAT reduction from 19\% to 16\% in H2~2020). These instruments transmit through the standard Keynesian multiplier channel and are widely used, shifting the output gap by $\alpha_{DI} > 0$ per unit deployed. \footnote{Pre-pandemic estimates place expenditure multipliers in a range of approximately 0.5 to 1.5 \citep{ramey2018, blanchard2002}.}

However, applying these instrument to a pandemic environment is questionable on both theoretical and practical grounds. The pandemic  recession is not a demand shock but a supply disruption induced by containment \citep{guerrieri2022}: businesses are closed by government order, not by insufficient demand. Injecting purchasing power into an economy where the supply of goods and services is physically restricted cannot translate into expenditure at the usual rate \citep{chetty2020, coibion2020}. This raises the question whether $\alpha_{DI}$ under pandemic conditions retains a meaningful magnitude, is merely delayed, or vanishes entirely.

The second category, capacity preservation ($F^{\text{CP}}_k$), comprises measures that maintain the economy's productive structure rather than directly raising contemporaneous demand: short-time work schemes, such as the Swiss \textit{Kurzarbeitsentsch\"{a}digung} or the UK Coronavirus Job Retention Scheme; payroll subsidies; direct grants; and government-backed liquidity support, such as the French \textit{Pr\^{e}ts Garantis par l'\'{E}tat} or the German KfW-Sonderprogramm. When containment closes firms, CP does not immediately restore output; instead, it prevents temporary disruptions from becoming persistent losses by preserving firm-worker matches, supplier relationships, and organizational capital---the amplification mechanism emphasized by \citet{guerrieri2022}. This transmission justifies a lag: CP affects output mainly through the subsequent recovery rather than at implementation.

Capacity preservation is therefore decomposed into above-the-line and below-the-line instruments. Above-the-line measures, such as short-time work schemes, payroll subsidies, and direct grants, operate as fiscal flows that replace lost income or revenue. Below-the-line measures, such as loans and guarantees, operate through liquidity provision; their relevant effect is the cumulative stock of available support.

Combining the pandemic shocks from Property~2 with the fiscal instruments above, the output evolves according to:

%Above und Below einführen und den Lag
\begin{equation}\label{eq:output}
\begin{aligned}
      y_{k+1} &= \rho_y\, y_k + \alpha_S\, S_k
          + \alpha_{\text{above}}\, F^{\text{above}}_{k-2}
          + \alpha_{\text{below}}\, \bar F^{\text{below}}_k
          + \alpha_{\text{DI}}\, F^{\text{DI}}_{k-1}
          + \alpha_{S,\text{DI}}\, S_k F^{\text{DI}}_{k-1}
          - \beta_d\, d_k 
\end{aligned}
\end{equation}

\noindent where $\bar F^{\text{below}}_k = \sum_{s \leq k}(\tilde F^{\text{loans}}_s + \tilde F^{\text{guar}}_s)$ denotes the cumulative below-the-line stock at quarter $k$.

\noindent Equation~\eqref{eq:output} captures several distinct mechanisms. 

First, output dynamics are governed by the autoregressive coefficient $\rho_y$, which captures persistence in the output gap. Previously implemented stringency is therefore absorbed into this term through the lagged state $y_k$. Containment policy $S_k$ enters additively with a contemporaneous level effect $\alpha_S < 0$.

Second, capacity-preservation measures enter through two distinct channels. Above-the-line measures enter with a two-period lag through $F^{\text{above}}_{k-2}$ and are governed by the coefficient $\alpha_{\text{above}}$. These measures substitute income or revenue streams, help maintain income, and reduce the persistence of the output shock. Below-the-line measures enter through the cumulative stock variable $\bar F^{\text{below}}_k$, with coefficient $\alpha_{\text{below}}$. This reflects the fact that loans and guarantees operate less through contemporaneous expenditure flows than through the availability of liquidity. The relevant mechanism is therefore the accumulated amount of loans and guarantees available up to quarter $k$, which helps reduce liquidity shortages and thereby limits structural damage.

Third, demand-injection measures enter with a one-period lag through $F^{\text{DI}}_{k-1}$ and operate through the usual multiplier effect, governed by $\alpha_{\text{DI}}$. Their effectiveness, however, depends on the current restrictiveness of the economy. This state dependency is captured by the interaction term $S_k F^{\text{DI}}_{k-1}$, with coefficient $\alpha_{S,\text{DI}}$.

Finally, mortality losses enter through $d_k$ with coefficient $-\beta_d$, capturing the adverse effect of excess deaths on economic activity.

The preceding decomposition might suggest a straightforward resolution of the output--health trade-off: deploy fiscal support to offset the lockdown-induced contraction, then exit. However, this reasoning abstracts from the budgetary consequences of the fiscal responses. Each instrument that stabilizes the output gap carries fiscal consequences for the government's balance sheet---and the magnitude of this fiscal cost is heterogeneous across categories and across several dimensions not yet established in the literature. To formalize this constraint, I introduce the third state variable: the \textit{pandemic debt gap} (\(b_k\)), defined as the deviation of real gross public debt from its country-specific pre-pandemic level, normalized by 2019 GDP.\footnote{I normalize by 2019 GDP rather than contemporaneous GDP to avoid mechanical denominator effects from the recession. The state \(b_k\) therefore captures changes in the real stock of government liabilities relative to the country's own pre-pandemic debt position, rather than cross-country differences in inherited debt levels.}


\begin{equation}\label{eq:debt}
\begin{aligned}
b_{k+1} &= (1+r)\, b_k - \gamma_y\, y_k
        + \kappa_{\text{above}}\, F^{\text{above}}_k
        + \kappa_{\text{loans}}\, F^{\text{loans}}_k 
        +\kappa_{\text{guar}}\, F^{\text{guar}}_k
        + \kappa_{\text{DI}}\, F^{\text{DI}}_{k-1}
        +\kappa_H F^H_{k-1}
\end{aligned}
\end{equation}

\noindent where \(r\) is the quarterly sovereign interest rate, \(\gamma_y>0\) captures the automatic stabilizer channel, and the \(\kappa\) coefficients are instrument-specific budgetary pass-throughs into the pandemic debt gap. Since \(y_k<0\) represents a negative output gap, the term \(-\gamma_y y_k\) increases \(b_{k+1}\) through lower revenues and higher automatic transfers.

The split of capacity-preservation measures in the debt equation reflects their different fiscal footprints. Above-the-line measures, such as short-time work schemes, payroll subsidies, and direct grants, generate immediate budgetary outlays, although these may partly operate through existing fiscal structures. Loans also create immediate fiscal exposure because the state extends new credit and must finance the lending operation, even though it receives a financial claim on the borrower in return. Guarantees, by contrast, do not generate an immediate disbursement but create an immediate contingent liability whose realized fiscal cost depends on the default rate among supported borrowers \citep{imf2022fiscalmonitor}. The empirical magnitudes of $(\kappa_{\text{above}}, \kappa_{\text{loans}}, \kappa_{\text{guar}})$ are estimated separately in Section~\ref{sec:results_empirical}; their heterogeneity is central to the planner's problem, because it implies that the fiscal cost of output stabilization depends not only on the volume but on the composition of the fiscal package.

Demand injection and health spending enter with a one-period lag because their budgetary impact may arise only after transfers are paid, tax relief affects revenues, procurement is settled, or reimbursement claims are recorded. For demand injection, this timing also aligns the debt cost with the period in which the measure becomes economically effective. Since the marginal output effect in equation~\eqref{eq:output} is $\alpha_{\text{DI}}+\alpha_{S,\text{DI}}S_k$ with $\alpha_{S,\text{DI}}<0$, the output return to demand injection falls with containment intensity, while its debt cost remains positive. Demand injection is therefore less attractive when containment is high. Health spending, $F^H_{k-1}$, follows the same lagged fiscal-pass-through logic, but enters only as an exogenous debt component because it is not part of the planner's optimized fiscal choice set.

The optimal fiscal mix therefore depends jointly on three margins: the prevailing level of stringency, which determines the effectiveness of demand injection; the composition of the capacity-preservation package, which determines the fiscal cost of liquidity support; and the cumulative fiscal legacy of past decisions, which affects the debt state $b_k$. The problem is therefore genuinely dynamic: the planner does not simply choose the volume of fiscal support, but must allocate support across instruments whose output effects and fiscal costs differ across states and over time.

%\vspace{5mm}
%\noindent\textit{Summary.} Property~3 introduces the fiscal compensation instruments---demand injection ($F^{\text{DI}}_k$), above-the-line capacity preservation ($F^{\text{above}}_k$), and below-the-line liquidity support ($\bar F^{\text{below}}_k$)---as well as the pandemic debt-gap state (\(b_k\)). Demand injection operates through a state-dependent multiplier, governed by $\alpha_{\text{DI}}$ and $\alpha_{S,\text{DI}}$, while capacity preservation stabilises output through delayed fiscal flows and cumulative liquidity provision, governed by $\alpha_{\text{above}}$ and $\alpha_{\text{below}}$. Since these instruments differ in both output effects and fiscal pass-throughs, the optimal fiscal response is inherently dynamic.


\subsubsection*{Property 4: Endogenous Fiscal Feedback via Automatic Stabilizers}

While Property~3 introduced the direct budgetary costs of discretionary fiscal instruments, Property~4 captures the endogenous fiscal feedback generated by the recession itself. When NPIs contract economic output (Property~2), the resulting negative output gap triggers automatic stabilizer effects. The parameter $\gamma_y > 0$ in the debt transition~\eqref{eq:debt} captures the aggregate sensitivity of the budget to the output gap.\footnote{For OECD economies, estimates of the budget semi-elasticity range from approximately 0.35 to 0.65, with a cross-country average around 0.5 \citep{repec:oecd}.}

Because this term operates symmetrically, it has a direct implication for instrument choice: a fiscal instrument whose output effect is large relative to its fiscal cost generates a double dividend---it stabilizes output directly and attenuates the automatic fiscal deterioration indirectly, as the improved $y_k$ feeds back through $\gamma_y$ to slow debt accumulation.

The central contribution of this property is to close the trilemma by establishing that fiscal abstinence is not a feasible strategy \citep{Fornaro2020}. Setting all discretionary instruments to zero ($F_k^{DI} = F_k^{CP} = 0$) in~\eqref{eq:debt} yields:
%
\begin{equation}
    b_{k+1} = (1+r) \, b_k - \gamma_y \, y_k
    \label{eq:debt_no_fiscal}
\end{equation}
%
Since \(y_k<0\) during containment, \(-\gamma_y y_k>0\): the pandemic debt gap increases even without any discretionary fiscal expenditure. A lockdown without accompanying fiscal measures still erodes the fiscal position---the trilemma cannot be circumvented by fiscal abstinence. Combined with the necessity of suppression established in Property~1, this means that the planner's decision is not whether to deploy fiscal support but which composition to deploy and when---returning the problem to the instrument choice formalized in Property~3.

%\vspace{5mm}
%\noindent\textit{Summary.} Property 4 introduces the budget semi-elasticity $\gamma_y$ and establishes that fiscal abstinence is an illusion: suppression passively erodes the fiscal position even absent discretionary intervention. The planner's decision is therefore not whether to intervene fiscally, but how to compose the intervention to minimize the pandemic debt footprint.

\vspace{5mm}
% ============================================================
%\subsection{Summary}
%\label{sec:Sumary}
% ============================================================
\noindent The four properties jointly characterize and formalized the structural environment of the pandemic trilemma that define this unique situation government faced.\footnote{An application of the four constitutive properties to four alternative crisis types: armed conflict, natural disasters, financial crises, and climate change is showed in Appendix~\ref{sec:comparison}.} The epidemiological block establishes the need for suppression; the suppression instrument improves the health state only by reducing economic activity; fiscal instruments can compensate the output loss but differ in their transmission channels and fiscal costs; and the debt state makes the problem dynamically constrained through both discretionary spending and automatic stabilizers. The planner therefore faces a joint transition problem over output, the pandemic debt gap, viral prevalence, and excess mortality. Therefore, the resulting transition system is:

\begin{align}
  y_{k+1} &= \rho_y\, y_k + \alpha_S\, S_k
          + \alpha_{\text{above}}\, F^{\text{above}}_{k-2}
          + \alpha_{\text{below}}\, \bar F^{\text{below}}_k \notag  \\
          &\quad + \alpha_{\text{DI}}\, F^{\text{DI}}_{k-1}
          + \alpha_{S,\text{DI}}\, S_k F^{\text{DI}}_{k-1}
          - \beta_d\, d_k
          \label{eq:trans_y} \\[4pt]
  b_{k+1} &= (1+r)\, b_k - \gamma_y\, y_k
          + \kappa_{\text{above}}\, F^{\text{above}}_k
          + \kappa_{\text{loans}}\, F^{\text{loans}}_k \notag \\
          &\quad + \kappa_{\text{guar}}\, F^{\text{guar}}_k
          + \kappa_{\text{DI}}\, F^{\text{DI}}_{k-1} +\kappa_{\text{H}}F^H_{k-1}
          \label{eq:trans_b} \\[4pt]
  \theta_{k+1} &= \rho_\theta \left(1 - \phi_S\, S_k\right)\theta_k
          + \varepsilon_k
          \label{eq:trans_theta} \\[4pt]
  d_{k+1} &= \delta_\theta\, \theta_k
          \label{eq:trans_d}
\end{align}

% =====================================================================
% Justification of the linear reduced-form approach
% =====================================================================
\noindent 

The transition system is deliberately reduced-form. Each channel has a microfounded counterpart in the pandemic macro literature, but the paper does not claim to identify policy-invariant deep parameters. Its purpose is to discipline a local retrospective benchmark within the realized OECD policy space. The reduced-form specification allows the transition parameters to be estimated directly and makes it possible to assess whether the state-transition structure implied by the theoretical framework is supported by the data.

The underlying mechanisms are well established in the pandemic-macro literature, although these mechanisms are typically studied separately. \citet{guerrieri2022} and \citet{bayer2023c} show that a sectoral supply shock can generate an endogenous demand shortfall and that standard transfers lose effectiveness while supply constraints bind, corresponding to the state-dependent demand-injection multiplier and the persistence channel of capacity preservation. \citet{BaqaeeFarhi2022} show that COVID-19 combined supply and demand shocks across sectors and that these shocks propagated, reinforcing the need to distinguish fiscal instruments by transmission channel rather than by aggregate size alone. \citet{fariaecastro2021} finds in a nonlinear DSGE model calibrated to the U.S. labor market that liquidity assistance to firms is more effective than household transfers consistent with the predicted ranking of capacity preservation over demand injection under high stringency. Similarly, \citet{CespedesChangVelasco2022} emphasize the role of firm-level financial frictions and productivity losses, providing a microfounded rationale for liquidity-preserving interventions. \citet{eichenbaum2021macro} derive the behavioral recession from optimizing agents who reduce consumption and labor supply to lower infection risk, the microfounded counterpart of the fear term. Finally, \citet{KaplanMollViolante2020} show in a heterogeneous-agent SIR economy that containment policy determines the health--output trade-off, while fiscal transfers mainly shape the distributional and aggregate economic consequences of that trade-off. Taken together, these contributions mirror the separation in Eqs.~(\ref{eq:trans_y})--(\ref{eq:trans_d})  The reduced-form system therefore does not invent its transmission channels. Rather, it combines mechanisms derived separately in the existing literature into a three-dimensional state-transition framework with explicit trade-offs what makes the problem tractable and allow me to analyze which strategies stabilized output most effectively relative to their fiscal cost and health environment.


%%Maybe overdo afiter conclusion

\section{Empirical Parametrization}

The purpose of this section is to bring the transitions system to the data. The resulting estimates are not interpreted as policy-invariant causal effects, but as empirical transition parameters within the realized pandemic policy space. They provide the numerical inputs required to solve the planner problem and to evaluate observed country strategies against the resulting benchmark. Table~\ref{tab:planner_objects} summarizes the target variables, state variables, and transition-system parameters that need to be measured or parameterized.

\begin{table}[H]
  \centering
  \scriptsize
  \renewcommand{\arraystretch}{1.15}
  \begin{tabular}{lp{9.5cm}l}
  \toprule
  \textbf{Symbol} & \textbf{Description} & \textbf{Block} \\
  \midrule

  \multicolumn{3}{l}{\textit{A. Planner objective and target variables}} \\
\midrule

$y_k$ 
& Output gap; measured as the percentage deviation of real GDP from its pre-pandemic potential trend 
& Target \\

$d_k$ 
& Excess mortality; measured as the deviation of deaths from the historical seasonal baseline 
& Target \\

$b_k$ 
& Pandemic debt-gap state; measured as the deviation of real gross public debt from the country-specific pre-pandemic level, normalized by 2019 GDP
& Target \\

$\theta_k$ 
& Viral prevalence; epidemiological state that affects welfare only through future excess mortality 
& State \\

  \midrule
  \multicolumn{3}{l}{\textit{B. Transition-system parameters}} \\
  \midrule

  $\rho_y$                 
  & Persistence of the output gap                                          
  & Output \\

  $\alpha_S$               
  & Contemporaneous effect of stringency $S_k$ on output                          
  & Output \\

  $\alpha_{\text{above}}$  
  & Lagged effect of above-the-line fiscal spending on output                    
  & Output \\

  $\alpha_{\text{below}}$  
  & Effect of the cumulative below-the-line liquidity stock on output  
  & Output \\

  $\alpha_{\text{DI}}$     
  & Lagged baseline effect of demand injection on output                              
  & Output \\

  $\alpha_{S,\text{DI}}$   
  & State-dependent interaction between stringency and demand injection         
  & Output \\

  $\beta_d$                
  & Effect of observed excess mortality on output through the behavioral response            
  & Output \\

  \midrule

  $r$                      
  & Quarterly sovereign interest rate                                  
  & Debt \\

  $\gamma_y$               
  & Budget semi-elasticity with respect to the output gap                    
  & Debt \\

  $\kappa_{\text{above}}$  
  & Fiscal pass-through of above-the-line spending into public debt                  
  & Debt \\

  $\kappa_{\text{loans}}$  
  & Fiscal pass-through of loan flows into public debt                        
  & Debt \\

  $\kappa_{\text{guar}}$   
  & Fiscal pass-through of guarantee flows into public debt                   
  & Debt \\

  $\kappa_{\text{DI}}$     
  & Fiscal pass-through of demand injection into public debt                        
  & Debt \\

  \midrule

  $\rho_\theta$            
  & Persistence of viral prevalence in the absence of containment                      
  & Infection \\

  $\phi_S$              
  & Effectiveness of stringency in reducing viral transmission                                       
  & Infection \\

  $\delta_\theta$          
  & Transmission from viral prevalence to excess mortality                                     
  & Mortality \\

  \bottomrule
  \end{tabular}

 \caption{Target variables and transition-system parameters}
\label{tab:planner_objects}
\end{table}



\subsection{Data and Measurement}
\label{sec:data}

I use a quarterly panel dataset covering all 38 OECD member states. The underlying data span Q1~2015 through Q4~2024, while the main estimation sample focuses on the pandemic period during which the trilemma was binding. The dataset combines macroeconomic outcomes, containment policy, epidemiological variables, and a newly constructed database of fiscal interventions. 


% --- TABLE 1: Variables and Sources ---
\begin{table}[H]
\centering
\caption{Variables and Data Sources}
\label{tab:data_sources}
\scriptsize
\begin{tabular}{@{}lp{7cm}l@{}}
\toprule
\textbf{Variable} & \textbf{Description} & \textbf{Source} \\
\midrule
\multicolumn{3}{@{}l}{\textit{Dependent variables}} \\[3pt]
$y_{ik}$
  & Output gap: deviation of real GDP from HP-filtered
    trend (2015--2019, $\lambda = 1600$) (pp)
  & \citet{OECDQNA} \\[6pt]
$b_{ik}$
  & Pandemic debt gap: deviation of real gross public debt from the country-specific pre-pandemic debt level, normalized by 2019 GDP (pp)
  & \citet{OECDQNA} \\[6pt]
\midrule
\multicolumn{3}{@{}l}{\textit{Policy variables}} \\[3pt]
$S_{ik}$
  & Stringency Index, population-weighted (0--100)
  & Own calc.\ (Appendix); \citet{Hale2021} \\[6pt]
$F_{ik}^{CP}$
  & Capacity preservation (\% of 2019 GDP)
  & Own database (Data Appendix) \\[6pt]
$F_{ik}^{DI}$
  & Demand injection (\% of 2019 GDP)
  & Own database (Data Appendix) \\[6pt]
$F_{ik}^{H}$
  & Health spending (\% of 2019 GDP)
  & Own database (Data Appendix) \\[6pt]
\midrule
\multicolumn{3}{@{}l}{\textit{Epidemiological variables}} \\[3pt]
$d_{ik}$
  & Excess mortality (P-score)
  & \citet{Msemburi2023WHO}; \citet{owid_em} \\[6pt]
$\theta_{ik}$
  & Infection prevalence, estimated from excess
    mortality via country-specific IFRs
  & Own calc.\ (Appendix); \citet{Levin2020} \\[6pt]
$\text{Vax}_{ik}$
  & Cumulative vaccination doses per capita
  & \citet{Mathieu2021} \\[6pt]
\bottomrule
\end{tabular}

\smallskip
{\scriptsize \textit{Notes:} All absolute values are expressed in real values.}
\end{table}

% --- PROSE: follows table order top to bottom ---

\noindent The output gap $y_{ik}$ is measured as the percentage deviation of quarterly real GDP from its Hodrick--Prescott-filtered trend, estimated exclusively on pre-pandemic data (2015--2019, $\lambda = 1600$) and projected forward, yielding a country-specific counterfactual that is comparable across countries in percentage-point terms.\footnote{Applying the HP filter to the full sample including 2020--2021 would pull the estimated trend downward, mechanically reducing the measured output gap. The pre-pandemic-only estimation ensures that counterfactual potential output reflects the economy's trajectory absent the pandemic shock.} The debt variable \(b_{ik}\) is constructed as the deviation of real gross public debt from its country-specific pre-pandemic level, normalized by 2019 GDP. Thus, \(b_{ik}=0\) does not represent the absence of public debt, but the absence of additional debt accumulation relative to the country's pre-pandemic debt position. Specifically, I take the real gross public debt stock in 2019Q4 as the baseline and express subsequent deviations from this level in percentage points of 2019 GDP. This measurement isolates pandemic-related debt dynamics from cross-country differences in initial debt levels. For the measurement of the containment I use a modified version of the widely applied Oxford Stringency Index \citep{Hale2021}, corrected for the share of the population
effectively affected by each measure: many restrictions applied heterogeneously---to the unvaccinated, to specific age groups, or exempting the recently recovered---so that the headline index overstates de facto stringency in periods of targeted policies. Weighting the index components
by the affected population share reduces the composition and aggregation vulnerability inherent in any policy index. Appendix~\ref{app:robustness_output} assesses sensitivity to this choice against alternative containment measures. 

The three core fiscal aggregates $F_{ik}^{CP}$, $F_{ik}^{DI}$, and $F_{ik}^{H}$ are constructed from a novel micro-level database of 2,301 individual policy measures across all 38 OECD member states, building on and substantially extending an initial draft by \citet{lacey2021}. The database records each measure at daily frequency with its announcement date, description, primary source, and fiscal cost (in local currency and as \% of 2019 GDP). The classification distinguishes instruments by their primary economic transmission mechanism required by the theoretical framework in Section~\ref{sec:trilemma}. Table~\ref{tab:fiscal_classification} details the composition of each channel while the full classification protocol and a high-frequency descriptive analysis for all 38 OECD countries are documented in the Data Appendix (Appendix~\ref{app:data}).

% --- TABLE 2: Fiscal Classification  ---
\begin{table}[H]
\centering
\caption{Fiscal Instrument Classification}
\label{tab:fiscal_classification}
\scriptsize
\begin{tabular}{@{}llr@{}}
\toprule
\textbf{Channel} & \textbf{Instruments} & \textbf{Measures} \\
\midrule
\multicolumn{3}{@{}l}{\textit{Capacity Preservation ($F^{CP}$)}} \\[3pt]
\quad Above-the-line
  & Short-time work, payroll subsidies, direct grants
  & 948 \\[3pt]
\quad Below-the-line
  & Government loans, credit guarantees
  & 400 \\[6pt]
\multicolumn{3}{@{}l}{\textit{Demand Injection ($F^{DI}$)}} \\[3pt]
\quad Transfers
  & Cash transfers, unemployment top-ups, ad-hoc benefits
  & 346 \\[3pt]
\quad Tax measures
  & Income tax relief, VAT reductions, tax credits
  & 162 \\[3pt]
\quad Demand stimulus
  & Infrastructure, green investment, tourism support
  & 38 \\[6pt]
\multicolumn{3}{@{}l}{\textit{Health ($F^{H}$)}} \\[3pt]
\quad
  & Medical procurement, hospital capacity, vaccine rollout
  & 232 \\[6pt]
\midrule
\multicolumn{2}{@{}l}{\textit{Excluded: Tax deferrals (IMF Cat.\ 3)}}
  & 154 \\[3pt]
\multicolumn{2}{@{}l}{\textit{Excluded: Unclassified and excluded policy codes}}
  & 21 \\[3pt]
\midrule
\multicolumn{2}{@{}l}{\textbf{Total}}
  & \textbf{2,301} \\
\bottomrule
\end{tabular}

\smallskip
{\scriptsize \textit{Notes:} Classification by primary transmission mechanism. Tax deferrals
excluded as realized fiscal cost approaches zero. Full protocol in Data Appendix
(Appendix~\ref{app:data}).}
\end{table}

\noindent Excess mortality \(d_{ik}\)---the deviation of all-cause deaths from a projected seasonal baseline---is the preferred measure of pandemic health impact for cross-country comparison. Unlike reported COVID-19 deaths, which depend on testing capacity, health infrastructure, and cause-of-death attribution practices that varied widely across countries, excess mortality captures both direct viral fatalities and indirect deaths from healthcare-system overload, deferred care, and resource diversion \citep{Msemburi2023WHO}. Excess mortality also enters the output equation as the fear term \(d_{ik}\), capturing voluntary reductions in economic activity based on \citet{GOOLSBEE2021}. Reported case counts are an unreliable measure of true infection prevalence, as testing capacity and reporting standards varied systematically across countries and over time. The infection prevalence variable \(\theta_{ik}\) is therefore recovered indirectly. I first convert excess-mortality P-scores into excess-death counts using country-specific baseline mortality, and then infer infections by dividing excess deaths by country-specific age-adjusted infection fatality rates \citep{Levin2020}. Normalizing by population yields a cross-country comparable measure of pandemic pressure that is not contaminated by differential testing regimes.


\subsection{Descriptives}

Table~\ref{tab:descriptives} provides descriptive statistics for the key variables at the country-quarter level (Panel~A) and the country level (Panel~B). 

\begin{table}[H]
\centering
\caption{Summary Statistics}
\label{tab:descriptives}
\scriptsize
\begin{tabular*}{\textwidth}{@{\extracolsep{\fill}} l r r r r r r @{}}
\toprule
 & Mean (SD) & Median & Min & P25 & P75 & Max \\
\midrule
\multicolumn{7}{@{}l}{\textbf{Panel A: Quarterly observations ($N = 380$ country-quarters)}} \\[4pt]
\multicolumn{7}{@{}l}{\textit{\textbf{Outcome variables}}} \\[2pt]
Output gap, pp of potential GDP
  & $-4.36\;(4.77)$  & $-3.74$ & $-23.47$ & $-6.34$ & $-1.74$ & $11.13$ \\[2pt]
Change in pandemic debt gap, \(\Delta b_{ik}\), pp of 2019 GDP
  & $0.85\;(2.48)$   & $0.46$  & $-7.46$  & $-0.61$ & $2.00$  & $14.20$ \\[3pt]
\multicolumn{7}{@{}l}{\textit{\textbf{Policy instruments}}} \\[2pt]
Containment stringency index (0--100)
  & $40.19\;(18.70)$ & $39.80$ & $0.00$   & $28.24$ & $54.45$ & $85.66$ \\[2pt]
Capacity preservation $F^{CP}$, \% of 2019 GDP
  & $1.31\;(3.03)$   & $0.12$  & $0.00$   & $0.00$  & $1.09$  & $25.42$ \\[2pt]
Demand injection $F^{DI}$, \% of 2019 GDP
  & $0.26\;(0.71)$   & $0.00$  & $0.00$   & $0.00$  & $0.14$  & $8.32$  \\[2pt]
Health expenditures $F^{H}$, \% of 2019 GDP$^{\dagger}$
  & $0.12\;(0.38)$   & $0.00$  & $0.00$   & $0.00$  & $0.00$  & $4.25$  \\[3pt]
\multicolumn{7}{@{}l}{\textit{\textbf{Epidemiological controls}}} \\[2pt]
Excess mortality, P-score (\%)
  & $10.23\;(14.87)$  & $5.51$  & $-13.36$ & $0.57$  & $15.30$ & $87.21$ \\[2pt]
Infection prevalence, \% of population
  & $1.18\;(1.89)$    & $0.41$  & $0.00$   & $0.10$  & $1.43$  & $11.45$ \\[2pt]
Cum.\ vaccination rate, \% of pop.
  & $31.35\;(33.04)$  & $12.46$ & $0.00$   & $0.00$  & $66.52$ & $89.68$ \\[6pt]
\midrule
\multicolumn{7}{@{}l}{\textbf{Panel B: Cross-country totals ($N = 38$ countries)}} \\[4pt]
\multicolumn{7}{@{}l}{\textit{\textbf{Outcome variables}}} \\[2pt]
Avg.\ output gap, pp of potential GDP
  & $-4.36\;(3.05)$    & $-4.22$  & $-10.58$ & $-5.96$  & $-2.58$  & $5.99$   \\[2pt]
Cum.\ change in pandemic debt gap, pp of 2019 GDP
  & $8.55\;(5.44)$     & $8.47$   & $-2.21$  & $5.31$   & $12.31$  & $20.64$  \\[3pt]
\multicolumn{7}{@{}l}{\textit{\textbf{Policy instruments}}} \\[2pt]
Avg.\ containment stringency (0--100)
  & $40.19\;(6.22)$    & $40.25$  & $27.66$  & $36.51$  & $43.65$  & $56.09$  \\[2pt]
Cum.\ capacity preservation $F^{CP}$, \% of 2019 GDP
  & $13.11\;(6.99)$    & $11.02$  & $0.88$   & $8.01$   & $17.86$  & $29.79$  \\[2pt]
Cum.\ demand injection $F^{DI}$, \% of 2019 GDP
  & $2.61\;(2.54)$     & $1.89$   & $0.11$   & $0.95$   & $3.83$   & $11.89$  \\[2pt]
Cum.\ health expenditures $F^{H}$, \% of 2019 GDP$^{\dagger}$
  & $1.15\;(1.24)$     & $0.73$   & $0.00$   & $0.31$   & $1.70$   & $4.73$   \\[3pt]
\multicolumn{7}{@{}l}{\textit{\textbf{Epidemiological outcomes}}} \\[2pt]
Cum.\ excess mortality, P-score (\%)
  & $102.25\;(76.95)$  & $88.74$  & $-13.59$ & $49.54$  & $131.05$ & $355.59$ \\[2pt]
Cum.\ infection prevalence, \% of pop.
  & $11.82\;(5.39)$    & $10.88$  & $3.33$   & $8.13$   & $15.11$  & $31.28$  \\[2pt]
Cum.\ vaccination rate, \% of pop.
  & $73.31\;(9.47)$    & $74.50$  & $45.67$  & $67.49$  & $80.31$  & $89.68$  \\
\bottomrule
\end{tabular*}

\vspace{4pt}
\parbox{\textwidth}{\footnotesize \textit{Notes:} Standard deviations in parentheses. Panel~A reports the distribution of quarterly country-level observations; Panel~B reports the cross-country distribution of period averages (output gap, stringency) and cumulative totals (all other variables) over Q1\,2020--Q2\,2022. All variables defined in Section~\ref{sec:data}. $^{\dagger}$Health expenditures enter the debt equation only; excluded from the output equation.}
\end{table}

\noindent The average OECD country experienced an output gap of $-4.4$~pp over the sample period, with the deepest contraction reaching $-23.5$~pp in a single country-quarter and the worst country-level average equal to $-10.6$~pp. Cumulative pandemic debt accumulation averaged $8.6$~pp of 2019 GDP, ranging from $-2.2$~pp to $20.6$~pp, indicating substantial cross-country heterogeneity in fiscal outcomes. Containment stringency was relatively synchronized across countries: the cross-country standard deviation of period-average stringency is $6.2$, compared to a within-country quarterly standard deviation of $18.7$. Fiscal composition exhibits the opposite pattern. The quarterly distributions of both $F^{CP}$ and $F^{DI}$ are strongly right-skewed: the mean of $F^{CP}$ ($1.31$\% of GDP) exceeds the median ($0.12$\%) by an order of magnitude, with $F^{CP}=0$ in $37.6$\% of country-quarters, reflecting the concentrated deployment of capacity preservation during the acute containment phase of 2020. Cumulatively, CP ($13.1$\% of 2019 GDP) dominated DI ($2.6$\%) by a factor of five, but the cross-country coefficient of variation is substantially larger for DI, indicating that countries that deployed demand injection did so at very different scales.

The epidemiological variables also display wide cross-country dispersion. Cumulative excess mortality ranges from $-13.6$ P-score points---reflecting a net reduction in non-COVID mortality in some countries---to $355.6$ points, a cumulative mortality burden more than three times the projected baseline. Infection prevalence and vaccination rates follow similar patterns of cross-country divergence. This compositional heterogeneity, together with the weak association between fiscal support and containment intensity documented in Section~\ref{sec:estimation}, provides the empirical variation used to parameterize the transition system. 

\begin{figure}[H]
    \centering
    \includegraphics[width=0.6\textwidth]{fig_ts_key_variables_v2.pdf}
    \caption{OECD Pandemic Trajectories by Quarter}
    \label{fig:trajectories}
    \vspace{4pt}
    \parbox{\textwidth}{\footnotesize \textit{Notes:} Each panel plots the cross-country median (solid line) and interquartile range (P25--P75, shaded band) across 38 OECD countries per quarter, Q1\,2020--Q2\,2022.}
\end{figure}

\noindent Figure~\ref{fig:trajectories} shows the evolution of the key variables over time. The narrow interquartile range for containment stringency, compared to the much wider bands for capacity preservation and demand injection, illustrates that countries followed relatively similar containment trajectories but differed substantially in the composition and scale of fiscal support. For more dynamics of the variable see Appendix \ref{app:further_descriptives}


\section{Parameterization Strategy}
\label{sec:estimation}

% -----------------------------------------------------------------
%  3. ESTIMATING EQUATIONS
% -----------------------------------------------------------------
I parameterize the transition system equation by equation. The main text focuses on the output and debt transitions, where the paper's contribution lies. The health-side equations are estimated separately and reported in Appendix~\ref{app:health_estimations}.\footnote{\label{fn:weekly}These equations are estimated at weekly frequency, because quarterly estimation would average over time frames too coarse for the underlying epidemiological dynamics. The estimates align with the existing literature; for the baseline specification I therefore adopt the corresponding literature values for the infection parameters and calibrate the wave-specific mortality mapping from the literature alone. My own estimates serve as a robustness check (see \ref{tab:params}).} The output estimating equation maps directly onto Eq.~(\ref{eq:trans_y}):

%
\begin{equation}
\begin{aligned}
    y_{ik} &= \rho_y\, y_{i,k-1}
    + \alpha_S\, S_{ik}
    + \alpha_{\text{above}}\, F^{\text{above}}_{i,k-2}
    + \alpha_{\text{below}}\, K^{\text{below}}_{ik} \\
    &\quad
    + \alpha_{\text{DI}}\, F^{\text{DI}}_{i,k-1}
    + \alpha_{S,\text{DI}}\, S_{ik} F^{\text{DI}}_{i,k-1}
    + \beta_d\, d_{ik}
    + \gamma_i + \varepsilon_{ik}.
\end{aligned}
\label{eq:output_est}
\end{equation}
%
where $K^{\text{below}}{ik}=\sum{s\leq k}
\left(F^{\text{loans}}{is}+F^{\text{guar}}{is}\right)$ denotes the cumulative below-the-line liquidity stock, and $\gamma_i$ denotes country fixed effects. The baseline specification includes country fixed effects but omits unrestricted quarter fixed effects. This choice is deliberate. The objective is to parameterize the transition of output and debt over the pandemic trajectory where quarter fixed effects would absorb much of the common pandemic variation through which containment and fiscal deployment occurred. Since the model explicitly includes observed measures of the pandemic state and policy response, the baseline follows \citet{deb2021} in controlling for the shock directly rather than absorbing it. Sensitivity checks with alternative time controls are reported in Section~\ref{app:robustness_output}.


The autoregressive coefficient $\rho_y$ captures persistence in the output gap, while $\alpha_S$ measures the contemporaneous output effect of containment. The preferred timing specification follows the mechanisms derived in Section~\ref{sec:trilemma}: above-the-line capacity preservation enters as a two-quarter-lagged flow, $F^{\text{above}}_{i,k-2}$; below-the-line support enters through the cumulative liquidity stock $K^{\text{below}}_{i,k}$; and demand injection enters with a one-quarter lag, $F^{\text{DI}}_{i,k-1}$. The
interaction $S_{i,k} F^{\text{DI}}_{i,k-1}$ captures whether the demand-injection multiplier depends on the restrictiveness of containment.\footnote{Section~\ref{app:robustness_output} reports alternative lag
specifications.} A negative $\alpha_{S,\text{DI}}$ would indicate that demand injection becomes less effective under high containment. The fear channel remains part of the theoretical transition system as an exogenous control, but is omitted from the quarterly baseline specification because the excess-mortality proxy is weakly identified at the quarterly level and has no influence on the coefficients of interest.\textsuperscript{\ref{fn:weekly}}



For the debt transition, I adjust the estimating equation so that the mechanical debt-service component is netted out before estimation rather than estimated directly:
%
\begin{equation}
\begin{aligned}
    \underbrace{\Delta b_{ik} - r_t b_{i,k-1}}_{\equiv\, \Delta \tilde b_{ik}}
    &= \gamma_y, y_{i,k-1}
    + \kappa_{\text{above}}\, F_{ik}^{\text{above}}
    + \kappa_{\text{loans}}\, F_{ik}^{\text{loans}}
    + \kappa_{\text{guar}}\, F_{ik}^{\text{guar}} \\
    &\quad
    + \kappa_{\text{DI}}\, F_{ik-1}^{\text{DI}}
    + \kappa_H\, F_{ik-1}^{H}
    + \mu_i + u_{ik}.
\end{aligned}
\label{eq:debt_est}
\end{equation}
%

\noindent where $\Delta b_{ik}=b_{ik}-b_{i,k-1}$ is the quarterly change in the pandemic debt-gap state, measured in percentage points of 2019 GDP, and $r_t b_{i,k-1}$ approximates the mechanical debt-service component. The estimates are not sensitive to whether the mechanical debt-service component is explicitly netted out. This is consistent with the low borrowing-cost environment during the main pandemic period: the average cost of new borrowing remained low, increasing only from around 1\% in 2020 to 1.4\% in 2021 before rising after the pandemic period \citep{oecd2021borrowing,oecd2023borrowing}. I nevertheless use the adjusted specification as the baseline because it aligns with the debt-gap transition in Eq.~(\ref{eq:trans_b}).

The lagged output gap $y_{i,k-1}$ is used because the debt equation describes the transition from $k-1$ to $k$. It therefore treats output as a predetermined state variable observed at the beginning of the transition. Therefore, the coefficient $\gamma_y$ captures the automatic-stabilizer channel: a larger negative output gap increases debt accumulation through lower revenues and higher transfer payments.

The fiscal variables enter the debt equation according to their measurement and budgetary timing. Since all fiscal measures are assigned to their announcement date, the timing captures the pass-through from announced policy commitments to measured debt accumulation. Above-the-line capacity preservation enters contemporaneously because these measures were designed as immediate emergency support and typically created fiscal obligations within the same transition period. Loans and guarantees are also recorded as announced commitments, but are adjusted by take-up rates to approximate effective below-the-line fiscal exposure, following evidence that announced envelopes were only partially utilized during the pandemic \citep{esrb2021monitoring, europeanparliament2022guarantees}. Demand injection and health-related spending enter with a one-quarter lag. This accounts for delayed budgetary effects. The coefficients $\kappa_{\text{above}}$, $\kappa_{\text{loans}}$, $\kappa_{\text{guar}}$, $\kappa_{\text{DI}}$, and $\kappa_H$ therefore measure instrument-specific debt-cost parameters in the transition system. Country fixed effects $\mu_i$ absorb time-invariant differences in fiscal institutions, debt-management practices, and accounting conventions. Robustness checks with contemporaneous output, alternative take-up assumptions, contemporaneous demand injection and health spending, and alternative timing structures are reported in Section~\ref{app:robustness_debt}.


\subsection{Parameterization and Scope of Interpretation}
\label{subsec:scope}

Several features of the pandemic setting complicate causal interpretation. This is well recognized in the related literature, which explicitly aims to identify causal effects of pandemic policies. Although this paper does not pursue such a causal identification strategy, the interpretation of the estimated transition parameters still requires clarification. Containment and fiscal support were not randomly assigned: both may respond to pandemic severity, expected output losses, and country-specific policy capacity. Moreover, containment and fiscal support were implemented simultaneously, while treatment intensity was continuous, time-varying, and global.
The baseline specification therefore relies on a conditional sequential-exogeneity assumption appropriate for a dynamic panel setting: after controlling for country fixed effects, lagged economic states, observed pandemic conditions, and the timing structure implied by the transition system in Section~\ref{sec:trilemma}, the remaining innovations to output and debt are assumed not to be systematically correlated with the predetermined policy variables entering the transition equations. Country fixed effects absorb time-invariant differences in institutions, fiscal capacity, health systems, and policy preferences, while the lag structure reduces contemporaneous reverse causality from realized output to fiscal deployment. The weak correlation between containment intensity and fiscal deployment in the data \((r=0.04\)) further mitigates the concern that fiscal support was mechanically driven by stringency alone (see \ref{app:concerns} for the full discussion).

The lagged output gap follows directly from the state-transition structure of the model. Pandemic output losses were persistent, and omitting the lag would force the policy variables to absorb part of this persistence. Its inclusion turns the specification into a dynamic fixed-effects panel, so the estimates are interpreted under the conditional sequential-exogeneity assumption discussed above. The within estimate of $\rho_y$ may therefore be subject to dynamic-panel bias.\footnote{I interpret the persistence estimate as a lower bound. Bias-corrected dynamic-panel specifications show that the results are not sensitive to this concern; see Section~\ref{app:concerns}.}

All baseline specifications report cluster-robust HC1 standard errors clustered at the country level. Since the number of country clusters is moderate (\(N=38\)), I additionally report wild cluster bootstrap (\(p\))-values clustered at the country level as a robustness.

The relevant empirical variation comes primarily from within-country and within-quarter differences, e.g., in the timing, scale, and composition of fiscal responses across countries facing a broadly common pandemic trajectory. In particular, countries differed substantially in their reliance on different fiscal instruments as seen in Figure \ref{fig:trajectories}. This compositional heterogeneity is used to discipline the channel-specific transition parameters. The estimates should therefore not be read as randomized treatment effects, but as conditional transition estimates that quantify whether the signs, timing, and relative magnitudes implied by the theoretical framework in Section~\ref{sec:trilemma} are consistent with the data. This allows me to assess whether the mechanisms underlying the pandemic trilemma were empirically present and binding in the realized OECD policy space. Remaining concerns---including alternative time controls, vaccination dynamics, fiscal classification choices, dynamic-panel bias as well as a discussion about alternative estimation methods---are addressed in the Appendix ~\ref{app:identification} and ~\ref{app:robustness}.
%%Maybe hier noch weitere causal results hinein-> um zu stärken

\subsection{Results}
\label{sec:results_empirical}
\subsubsection{Output}
\label{subsub:results_output}

Table~\ref{tab:output_transition} reports the output-transition estimates, gradually moving from a parsimonious state/shock specification to the preferred baseline transition equation.

\begin{table}[H]
\centering
\caption{Output Transition Estimates}
\label{tab:output_transition}
\scriptsize
\renewcommand{\arraystretch}{1.2}
\begin{tabular}{@{} l c c c c @{}}
\toprule
 & (1) & (2) & (3) & (4) \\
 & State/shock & Aggregate fiscal & Decomposed fiscal & Full baseline \\
 & only & support & channels & (preferred) \\
\midrule
\multicolumn{5}{@{}l}{\textit{Persistence channel:}} \\[2pt]
$y_{i,k-1}$ \,($\rho_y$) & $0.154***$ & $0.064$ & $0.221***$ & $0.231***$ \\
 & (0.043) & (0.046) & (0.052) & (0.050) \\[4pt]
\multicolumn{5}{@{}l}{\textit{Containment channel:}} \\[2pt]
$S_{ik}$ \,($\alpha_S$) & $-0.101***$ & $-0.087***$ & $-0.101***$ & $-0.095***$ \\
 & (0.011) & (0.010) & (0.011) & (0.013) \\[4pt]
\multicolumn{5}{@{}l}{\textit{Aggregate fiscal channel:}} \\[2pt]
$F^{\text{total}}_{i,k-1}$ &  & $-41.315***$ &  &  \\
 &  & (9.556) &  &  \\[4pt]
\multicolumn{5}{@{}l}{\textit{Capacity-preservation channels:}} \\[2pt]
$F^{CP,\,\text{above}}_{i,k-2}$ \, \,[$\alpha_{\text{above}}$] &  &  & $0.588***$ & $0.544***$ \\
 &  &  & (0.202) & (0.189) \\[4pt]
$K^{CP,\,\text{below}}_{ik}$ \, \,[$\alpha_{\text{below}}$] &  &  & $0.304**$ & $0.261*$ \\
 &  &  & (0.142) & (0.145) \\[4pt]
\multicolumn{5}{@{}l}{\textit{Demand-injection channels:}} \\[2pt]
$F^{DI}_{i,k-1}$ \,($\alpha_{DI}$) &  &  & $-0.615$ & $1.470**$ \\
 &  &  & (0.419) & (0.569) \\[4pt]
$F^{DI}_{i,k-1} \times S_{ik}$ \,($\alpha_{S\cdot DI}$) &  &  &  & $-0.041***$ \\
 &  &  &  & (0.014) \\[4pt]
\midrule
Country FE & $\checkmark$ & $\checkmark$ & $\checkmark$ & $\checkmark$ \\
Observations & 418 & 418 & 418 & 418 \\
$R^2$ (adj.) & 0.565 & 0.638 & 0.584 & 0.590 \\
Within $R^2$ & 0.395 & 0.498 & 0.426 & 0.436 \\
\bottomrule
\end{tabular}
\vspace{4pt}
\parbox{\textwidth}{\scriptsize\textit{Notes:} Dependent variable: output gap $y_{ik}$, measured in percentage points of potential GDP. Sample: 38 OECD countries over the pandemic estimation window $t \in [4,14]$. All specifications include country fixed effects. Standard errors clustered at the country level with small-sample correction are reported in parentheses. Column~(1) includes only the lagged output gap and containment intensity. Column~(2) adds aggregate fiscal support. Columns~(3) and~(4) decompose fiscal support by transmission channel. Column~(4) is the preferred specification. The fear term is omitted from the baseline, as already stated in Section~\ref{sec:estimation} see Appendix \ref{app:robustness_output}. * $p<0.1$, ** $p<0.05$, *** $p<0.01$.}
\end{table} 


\noindent Containment remains stable, negative, and highly significant across all specifications. Output persistence is positive throughout, but becomes insignificant in the aggregate fiscal specification that is shown in Column~(2). This specification also shows why aggregate fiscal support is difficult to interpret: the coefficient has the wrong sign and is economically implausibly large. This suggests that an aggregate fiscal measure with a one-period lag mixes different instruments, timing structures, and potential reverse causality. The decomposition in column~(3) therefore provides a more meaningful specification: the capacity-preservation coefficients become positive and plausible, while demand injection remains insignificant without the interaction term.\footnote{This reflects aggregation, not irrelevance: the marginal effect $\alpha_{DI}+\alpha_{S\cdot DI}\,S$ is positive when the economy is open and vanishes at $S^{\star}\approx 36$. Since DI was deployed mostly during high-stringency quarters, the unconditional coefficient averages these regimes and is pulled below zero. A joint Wald test rejects $\alpha_{DI}=\alpha_{S\cdot DI}=0$ ($F=4.23$, $p=0.022$); given the high correlation between the two terms, only the marginal effect at a given $S$ is interpretable. See Appendix~\ref{app:identification}.}

In the preferred specification, the coefficients have the expected signs and are statistically significant at conventional levels. The persistence coefficient, at $0.231$, is clearly lower than what would typically be expected in a normal business-cycle environment. This is consistent with the pandemic setting, where output dynamics were driven less by standard internal propagation and more by containment, fiscal interventions, and the timing of the shock. Containment has a statistically significant negative level effect on output, in line with the theoretical framework and with the empirical literature on the economic costs of NPIs \citep{deb2021, Ashraf2022}.

Turning to the fiscal channels, above-the-line capacity preservation increases the output gap by $0.54$ percentage points per percentage point of 2019 GDP, while the below-the-line liquidity stock increases it by $0.26$ percentage points. These estimates should not be interpreted as standard fiscal multipliers. They
capture transition effects in a pandemic environment, where fiscal policy operates under binding health restrictions and supply constraints \citep{guerrieri2022, chetty2020}. Within this environment, the positive
coefficients indicate that capacity-preserving measures limited persistent damage to firms, employment relationships, and liquidity conditions, consistent with the role of emergency lifelines emphasized in the theoretical model and in the literature by \citet{deb2021, fariaecastro2020f}.

Demand injection is the channel closest to a standard fiscal multiplier. Its baseline coefficient is larger than one, but the negative interaction with containment implies strong state dependence. The marginal effect is given by $\alpha_{DI}+\alpha_{S,\text{DI}}S_{ik}$, which becomes zero at a stringency level of roughly $S^{*}\approx 36$. Above this threshold, demand injection no longer translates into higher output in the estimated transition equation. This is consistent with the underlying mechanism: demand-side measures rely on households and firms being able to spend the additional income. When containment restricts consumption opportunities and economic activity, the multiplier is attenuated or delayed \citep{chetty2020, guerrieri2022, coibion2020}. 

\subsubsection{Debt}
\label{subsub:results_debt}
The debt specifications in Table~\ref{tab:debt_transition} move from broad fiscal categories in column~(1), to separate headline amounts for loans and guarantees in column~(2), and finally to the preferred take-up-adjusted specification in column~(3).

 \begin{table}[H]
  \centering
  \caption{Debt Transition Estimates}
  \label{tab:debt_transition}
  \scriptsize
  \renewcommand{\arraystretch}{1.2}
  \begin{tabular}{@{} l c c c @{}}
  \toprule
   & (1) & (2) & (3) \\
   & Coarse & Headline & Preferred \\
   & categories & commitments & take-up \\
  \midrule
  \multicolumn{4}{@{}l}{\textit{Predetermined output state:}} \\[2pt]
  $y_{i,k-1}$ \,($-\gamma_y$) & $-0.120***$ & $-0.117***$ & $-0.117***$ \\
   & (0.039) & (0.039) & (0.039) \\[4pt]
  \multicolumn{4}{@{}l}{\textit{Above-the-line capacity preservation:}} \\[2pt]
  $F^{CP,\,\text{above}}_{ik}$ \,[$\kappa_{\text{above}}$] & $0.639***$ & $0.664***$ & $0.664***$ \\
   & (0.202) & (0.203) & (0.203) \\[4pt]
  \multicolumn{4}{@{}l}{\textit{Below-the-line liquidity support:}} \\[2pt]
  $F^{CP,\,\text{below}}_{ik}$ \, (aggregate) & $0.584***$ &  &  \\
   & (0.114) &  &  \\[4pt]
  $F^{CP,\,\text{loans}}_{ik}$ \, (headline) &  & $0.501***$ &  \\
   &  & (0.096) &  \\[4pt]
  $F^{CP,\,\text{guar}}_{ik}$ \, (headline) &  & $0.134*$ &  \\
   &  & (0.071) &  \\[4pt]
   $F^{CP,\,\text{loans}}_{ik}$ \, (0.60) &  &  & $0.836***$ \\
   &  &  & (0.160) \\[4pt]
  $F^{CP,\,\text{guar}}_{ik}$ \, (0.25) &  &  & $0.536*$ \\
   &  &  & (0.284) \\[4pt]
  \multicolumn{4}{@{}l}{\textit{Demand injection and exogenous health spending:}} \\[2pt]
  $F^{DI}_{i,k-1}$ \,[$\kappa_{DI}$] & $0.540*$ & $0.526*$ & $0.526*$ \\
   & (0.296) & (0.297) & (0.297) \\[4pt]
  $F^{H}_{i,k-1}$ \,[$\kappa_H$] & $0.893**$ & $0.908**$ & $0.908**$ \\
   & (0.377) & (0.378) & (0.378) \\[4pt]
  \midrule
  Country FE & $\checkmark$ & $\checkmark$ & $\checkmark$ \\
  Observations & 494 & 494 & 494 \\
  $R^2$ (adj.) & 0.185 & 0.188 & 0.188 \\
  Within $R^2$ & 0.232 & 0.236 & 0.236 \\
  \bottomrule
  \end{tabular}
  \vspace{4pt}
  \parbox{\textwidth}{\scriptsize\textit{Notes:} Dependent variable: adjusted quarterly change in the pandemic debt-gap state, $\Delta \tilde b_{ik}=\Delta
  b_{ik}-r_t b_{i,k-1}$, measured in percentage points of 2019 GDP. Sample: 38 OECD countries over the longer pandemic debt-estimation window $t \in [4,16]$ to capture the terminal liability. All specifications include country fixed effects. Standard errors clustered at the country level with small-sample correction are reported in parentheses.
  Column~(1) uses broad fiscal categories. Column~(2) uses headline commitments for loans and guarantees. Column~(3) is the preferred specification, using a
  60\% take-up rate for loans and a 25\% take-up rate for guarantees; based on \citep{esrb2021monitoring, europeanparliament2022guarantees}. Adjustments can be found in Appendix \ref{app:robustness_debt}. Demand injection enters with a one-quarter lag ($F^{DI}_{i,k-1}$), reflecting delayed budgetary pass-through.Health-related spending is included as an exogenous control with a one-quarter lag and is not part of the optimized fiscal instrument set.  * $p<0.1$, ** $p<0.05$, *** $p<0.01$.}
  \end{table}


\noindent The coefficient on the lagged output gap is stable and statistically significant across all three specifications. A one percentage point higher lagged output gap reduces debt accumulation by about $0.12$ percentage points of 2019 GDP. Conversely, a deeper recession increases the pandemic debt gap through lower revenues and higher automatic stabilizer payments. In the coarse specification, both above-the-line capacity preservation and below-the-line liquidity support increase debt accumulation by similar magnitudes. The coefficients are below one, which is plausible since fiscal measures do not translate mechanically one-for-one into the debt-gap variable: some measures are spread over several quarters, and some operate through existing budgetary frameworks \citep{imf2022fiscalmonitor}.\footnote{Additional split-sample estimates by the strength of pre-existing social safety nets suggest some institutional heterogeneity in debt pass-through, but do not alter the baseline conclusion that output losses and fiscal support both increase the pandemic debt gap.}

Columns~(2) and~(3) further split below-the-line support into loans and guarantees, yielding three statistically significant distinct debt-cost coefficients for capacity preservation. Loans have the largest estimated debt pass-through ($0.84$), followed by above-the-line spending ($0.64$) and guarantees ($0.536$). This ranking is economically plausible. Loans imply an actual cash outflow once taken up and often required newly financed public credit facilities, even if they also create a financial claim on the borrower. Above-the-line measures also generate direct budgetary outlays, but their pass-through can be lower because many schemes operated through existing fiscal structures and were spread over several quarters. Guarantees, by contrast, mainly create contingent liabilities and therefore have the lowest immediate impact on measured public debt, since they affect debt only when guarantees are called. The demand-injection side also shows a lagged positive and significant debt pass-through ($0.526$), a magnitude similar to guarantees and close to above-the-line capacity preservation, which provides the closest direct comparison in terms of debt pass-through. Health-related spending enters as an exogenous control and has a lagged positive and significant effect on debt accumulation. The coefficient is close to one, suggesting an almost one-for-one pass-through into the pandemic debt gap. This is plausible because health measures mainly reflect newly authorized spending for medical procurement, hospital capacity, and vaccination. These main qualitative patterns are stable across the relevant robustness checks (see Section \ref{app:robustness}).

Taken together, the output and debt estimates show the expected signs and economically plausible magnitudes. Containment reduces output, capacity preservation stabilizes output either with a delay or through cumulative liquidity provision, demand injection is state-dependent, and all fiscal instruments increase the pandemic debt gap. Combined with the health-side estimates, these results provide empirical support for the trilemma at the center of the model. Viral prevalence grows endogenously and becomes explosive without containment, making suppression necessary for health stabilization. Yet containment depresses output, while fiscal instruments that offset output losses widen the pandemic debt gap. The estimates therefore identify the central three-way tension between health, output, and debt.\footnote{The health-side dynamics are disciplined separately through the epidemiological parameters reported in Appendix~\ref{app:health_estimations} and \ref{app:calibration}.}

Viewed in isolation, the estimates suggest a clear ranking of instruments. Demand injection is attractive when containment is low, but loses effectiveness when the spending channel is constrained. Above-the-line capacity preservation is valuable for stabilization, but comes with relatively high fiscal costs. Loans and guarantees work through the same liquidity-preservation channel, yet differ strongly in their debt pass-through, making guarantees potentially cost-effective when they are credible enough to stabilize firms and lenders. This isolated ranking, however, is not sufficient to evaluate actual country strategies. The instruments interact with the state of the economy and with each other. A country with high containment, large demand-injection measures, and little capacity preservation may generate debt without much effective output stabilization. Conversely, liquidity guarantees in an unrestricted economy may have limited additional value if firms no longer face binding liquidity constraints and maybe if they come without wage subsidies they loose effectivness. The relevant question is therefore not which instrument is best in general, but which composition is best in a given state and at a given point in time.

This turns the empirical findings into a dynamic allocation problem. The social planner must choose not only the total amount of fiscal support, but also its composition and timing, taking into account containment, output persistence, liquidity needs, and the accumulated debt state. The next section calibrates the model with the estimated transition parameters and assesses whether the proposed transition system can reproduce the observed pandemic dynamics.

%\subsection{Heterogeneity}-> Maybe gar nicht nötig oder Appendix


%%%%

\subsection{Calibration}
\label{subsec:calibration}

Using the transition system from Eqs.~(\ref{eq:trans_y})--(\ref{eq:trans_d}) and the parameters in Table~\ref{tab:params}---I evaluate the system's fit on the 38 OECD economies. The objective of this step is to assess whether the parameterized transition system can reproduce the observed pandemic dynamics before it is used for representative and country-specific counterfactual policy evaluation. I initialize the system at the pre-pandemic steady state, with an output gap of zero, excess mortality equal to zero, and public debt equal to its country-specific pre-pandemic level. Formally, the debt baseline is given by $B^{\text{real}}_{i,2019\text{Q}4}$, and the simulated pandemic debt gap measures deviations from this baseline, normalized by 2019 GDP. I take the observed paths of containment, fiscal instruments, and the epidemiological state as given and roll the output, debt and health blocks forward over the pandemic horizon while the initial health shock in 2020Q2 is treated as the initial exogenous shock. Table~\ref{tab:calibration_fit} summarizes the model fit. The output block performs well: the median output-gap RMSE is 1.85 percentage points, or 40\% of the observed cross-country dispersion, so the model reproduces the main output dynamics. Debt dynamics are matched less tightly. The median RMSE of quarterly debt changes is $2.02$ percentage points, $85\%$ of the relevant dispersion. 

\begin{table}[H]
  \centering
  \caption{Model parameters: symbol, description, calibrated value, and equation block.}
  \label{tab:params}
  \scriptsize

  \begin{tabular}{lp{8cm}cl}
  \toprule
  \textbf{Symbol} & \textbf{Description} & \textbf{Value} & \textbf{Block} \\
  \midrule
  $\rho_y$                 & Persistence of output gap & $+0.231$ & Output \\
  $\alpha_S$               & Contemporaneous effect of stringency $S$ & $-0.095$ & Output \\
  $\alpha_{\text{above}}$  & Lag-2 effect of above-the-line fiscal spending & $+0.544$ & Output \\
  $\alpha_{\text{below}}$  & Contemporaneous effect of below-the-line stock (take-up adjusted) & $+0.261$ & Output \\
  $\alpha_{\text{DI}}$     & Lag-1 effect of direct income support & $+1.470$ & Output \\
  $\alpha_{S,\text{DI}}$   & Interaction of stringency with lag-1 direct income support & $-0.041$ & Output \\
  \midrule
  $r$                      & Quarterly sovereign interest rate & $+0.001$ & Debt \\
  $\gamma_y$               & Automatic-stabiliser sensitivity of debt & $-0.117$ & Debt \\
  $\kappa_{\text{above}}$  & Pass-through of above-the-line spending into debt & $+0.664$ & Debt \\
  $\kappa_{\text{loans}}$  & Pass-through of take-up-adjusted loan flows & $+0.836$ & Debt \\
  $\kappa_{\text{guar}}$   & Pass-through of take-up-adjusted guarantee flows & $+0.536$ & Debt \\
  $\kappa_{\text{DI}}$     & Pass-through of lag-1 direct income support & $+0.526$ & Debt \\
  \midrule
  $\rho_{\theta,k}$        & Infection persistence before/after vaccination break & $1.50/<1$ & Infection \\
  $\phi_S$                 & Stringency mitigation factor & $0.80$ & Infection \\
  $\delta_{\theta,k}$      & Wave-specific mortality mapping based on effective population IFRs & Appendix~\ref{app:health_estimations} & Mortality \\
  \bottomrule
  \end{tabular}
  \vspace{0.15cm}
  \begin{minipage}{0.95\textwidth}
  \scriptsize
  \textit{Notes:} The health parameters are calibrated using external epidemiological evidence, see Appendix \ref{app:health_estimations} and \ref{app:calibration} for the full discussion. 
  \end{minipage}
\end{table}


\noindent The model thus captures the direction and magnitude of the fiscal pass-throughs but is weaker on part of the initial short-run debt movement, owing to accounting differences (see Appendix~\ref{app:calibration} for the full discussion). The overall trend, the growth rate, and especially the levels are nonetheless reproduced closely: the debt level attains a median RMSE of $2.53$ percentage points and an RMSE/SD ratio of $0.06$, so the model tracks debt dynamics appropriately with a final mean debt residual of around 1 pp. 



\begin{table}[H]
\centering
\caption{Calibration fit of the transition system under observed policy paths.}
\label{tab:calibration_fit}
\scriptsize
\begin{tabular}{lccccp{3.8cm}}
\toprule
\textbf{Moment} & \textbf{Horizon} & \textbf{Median RMSE} & \textbf{Mean RMSE} & \textbf{RMSE/SD} & \textbf{Interpretation} \\
\midrule
Output gap & Q4.19--Q2.22 & 1.85 pp & 2.01 pp & 0.40 & Error is less than half of observed dispersion; dynamics are well reproduced. \\
Debt change & Q4.19--Q4.22 & 2.02 pp & 2.00 pp & 0.85 & Primary debt metric; weakest fitted component. \\
Debt level & Q4.19--Q4.22 & 2.53 pp & 3.02 pp & 0.06 & Debt levels are closely reproduced, partly because the initial debt stock is imposed. \\
Infection state $\theta$ & Q4.19--Q2.22 & 0.0000 & 0.0024 & 0.00 & See text. \\
Excess mortality $d$ & Q4.19--Q2.22 & 12.59 & 13.38 & 0.72 & Mortality dynamics are reproduced reasonably well, conditional on $\theta$. \\
\bottomrule
\end{tabular}
\vspace{0.15cm}
\begin{minipage}{0.95\textwidth}
\scriptsize
\textit{Notes:} RMSE/SD denotes the median RMSE relative to the cross-sectional standard deviation of the corresponding observed variable. Debt change is the primary debt validation metric because it tests the fiscal pass-through dynamics; debt levels are reported as a reference because the initial debt stock is imposed. 
\end{minipage}
\end{table}

\noindent The infection state $\theta$ enters the transition system in the median forward roll as a calibrated input, not as a fitted moment. By construction the calibrated innovations $\varepsilon_{\theta}$ absorb the entire gap between the simulated and observed infection path, so the reported $\theta$ RMSE of zero reflects this identity rather than model fit. Conditional on the observed infection path, the mortality equation maps $\theta$ into excess deaths with a median RMSE of $12.59$ deaths per million per week, or $72\%$ of the observed dispersion, so the wave-specific IFR mapping reproduces the level and timing of mortality reasonably well (Figure~\ref{fig:calib}).

Since in the first step of the main analysis the overall optimization is solved for the representative OECD economy, the mean trajectory is the object of interest, and here the planner faces a fully dynamic optimization problem. In this mean roll the infection state $\theta$ is no longer pinned to the observed path but enters as a state the planner can partly steer through containment, while the variant-driven innovations $\varepsilon_{\theta}$ remain exogenous. The exercise is therefore closer to a genuine optimization than a calibrated replay, yet the fit does not deteriorate: under observed mean inputs the representative economy reproduces the data at least as closely as the median roll---the mean output path with an RMSE of $1.24$ percentage points, the mean mortality path with an RMSE of $2.87$ deaths per million per week, better than the fixed, wave-adjusted median roll, and quarterly debt changes with an RMSE of $0.85$ percentage points (see Appendix \ref{app:calibration} for the full discussion). 


\begin{figure}[H]
    \centering
    \includegraphics[width=0.7\textwidth]{Files/text/main/Calibration V20.pdf}
    \caption{Calibration fit: OECD-median simulated and observed trajectories with IQR bands.}
    \label{fig:calib}
\end{figure}

\noindent Because the model is also used to solve country-specific planner problems, Table~\ref{tab:country_endpoint} in Appendix~\ref{app:calibration} reports the same output, mortality, and debt fit statistics at the individual-country level; the cross-country dispersion of these endpoint residuals is consistent with the aggregate moments above, supporting country-level as well as representative-economy counterfactual analysis. Overall, the calibrated transition system is sufficiently accurate for counterfactual policy evaluation.

\section{Optimization}
\label{sec:planner}

The planner chooses the control vector \(\bm{u}_k = (S_k,\, F^{\text{above}}_k,\, F^{\text{loans}}_k,\, F^{\text{guar}}_k,\, F^{\text{DI}}_k)'\) subject to the transition system in
Eqs.~(\ref{eq:trans_y})--(\ref{eq:trans_d}). Because direct income support disburses with a
one-quarter delay---entering both output and the budget at disbursement---and because the two
capacity-preservation channels operate through persistent stocks of accumulated fiscal
support, the state vector \((y_k, b_k, \theta_k, d_k)\) is augmented with the lagged
control and the two stocks to render the system first-order; Appendix~\ref{app:planner}
states the augmented system formally.

The two stocks are central to the model's economics. Below-the-line loans and guarantees
accumulate into a liquidity stock \(L_k\) and above-the-line transfers into a
capacity-preservation stock \(K_k\), each decaying geometrically:\footnote{The decay rates are calibrated in Section~\ref{subsec:calibration}; robustness checks with alternative and zero decay rates (Section~\ref{app:compl_results}) leave the main results unchanged.} 

\begin{equation}
  L_{k+1} = (1-\delta_L)\,L_k + F^{\text{loans}}_k + F^{\text{guar}}_k,
  \qquad
  K_{k+1} = (1-\delta_K)\,K_k + F^{\text{above}}_k,
  \label{eq:stocks}
\end{equation}

Because above- and below-the-line support operate through the same output channel but are distinguished in the debt dimension due to institutional differences---capacity
preservation---the model allows for an interaction between the two. Economically, the two
measures are complementary: above-the-line support (wage subsidies, for example) mainly
targets the preservation of the employment relationship, while below-the-line liquidity
targets the solvency of the firm itself, so each is more effective when the other is in
place, consistent with the argument made by \citet{QUELLE_ANGEBEN}. The complementarity is
captured by a bounded multiplier on the liquidity effect,

\begin{equation}
  m(K_{k+1}) = 1 + \chi\,\frac{K_{k+1}}{K_{k+1} + \bar K},
  \label{eq:complementarity}
\end{equation}
so that the output transition contains the term \(\alpha_L\, L_{k+1}\, m(K_{k+1})\) in
place of a simple additive liquidity effect, alongside a direct capacity effect
\(\alpha_K K_{k+1}\). The multiplier is bounded between \(1\) (no preserved capacity) and
\(1+\chi\) (capacity saturated relative to the empirical scale \(\bar K\)). The parameters
\(\chi\), \(\alpha_K\), and the half-life governing \(\delta_K\) are calibrated externally
(Section~\ref{subsec:calibration}), but the complementarity is not merely assumed:
[BEFUND], reported in [TABELLE/ABSCHNITT], shows the same pattern in the data.

The planner solves a finite-horizon discounted quadratic loss problem,

\begin{equation}
  \min_{\{\bm{u}_k\}_{k=0}^{N-1}}\;
  \frac{1}{2}\left\{
  \sum_{k=0}^{N-1} \beta^{k+1}
  \left[\tilde{\bm{x}}_{k+1}' \bm{Q}\, \tilde{\bm{x}}_{k+1}
  + \bm{u}_k' \bm{R}\, \bm{u}_k\right]
  + \beta^N \tilde{\bm{x}}_N' \bm{Q}_N\, \tilde{\bm{x}}_N
  \right\},
  \label{eq:objective}
\end{equation}
where \(\tilde{\bm{x}}_k \equiv \bm{x}_k - \bar{\bm{x}}\) denotes the deviation from the
target state and \(\beta = 0.99\) is the quarterly discount factor. The running loss is
evaluated on the states the policy brings about, \(\bm{x}_1, \ldots, \bm{x}_N\); the
predetermined initial state carries no loss. The model is solved for the main pandemic period, Q1~2020 through Q4~2021, where containment remained a relevant health instrument and it extends up to Q4~2022 to capture the fiscal legacy.

The problem is subject to two sets of policy
constraints. First, all controls are non-negative and bounded in their 99th percentile of the
corresponding in-support observed values across the OECD panel. Second, containment is subject to a ramping constraint, \(|S_k - S_{k-1}| \le \Delta S_{\max}\), with \(\Delta S_{\max}\) set to the 95th percentile of observed quarter-on-quarter changes in stringency: legal, administrative, and compliance frictions limited how fast restrictions could be tightened or lifted, and the constraint imposes this speed limit directly from the data. Because the transition system is nonlinear---through the complementarity multiplier and the interaction of
stringency with infection dynamics and lagged income support---the problem has no
closed-form solution; Appendix~\ref{app:planner} describes the iterative LQR algorithm,
the treatment of the constraints, and the safeguards against local optima.

The running and terminal losses are evaluated at the deviation from the target state
\(\bar{\bm{x}}\): for output, infection prevalence, and excess mortality the target is the
pre-pandemic benchmark of zero, while for public debt it is the pre-pandemic stock
\(b_0\)---the fiscal objective is to avoid accumulating debt beyond the level inherited at
the onset of the pandemic. The running state-cost matrix is diagonal,

\begin{equation}
  \bm{Q} = \mathrm{diag}\!\left(w_y,\; \tau_b\, w_b,\; 0,\; \lambda_d\, w_d,\; 0,\; 0,\; 0\right),
  \label{eq:Qmatrix}
\end{equation}


so the planner penalizes output-gap deviations, debt accumulation beyond \(b_0\), and
excess mortality. Infection prevalence receives zero direct weight because it is not a
welfare object in itself; it enters the loss only through its role in generating future
excess mortality via Eq.~(\ref{eq:trans_d}). The weights are normalized by economically
interpretable reference scales---an output gap of 5 percentage points, a debt increase of
10 percentage points of GDP, and excess mortality of 100 deaths per million per week---so
that the planner's preferences reduce to two free parameters: \(\tau_b\), the weight on
fiscal sustainability relative to output stabilization, and \(\lambda_d\), the weight on
mortality relative to normalized output loss while assuming in the baseline specification that the health dimension is the main target. Because the loss is quadratic,
\(\lambda_d\) is not itself the marginal mortality--output trade-off; I therefore anchor it to the value of a statistical life by setting \(\lambda_d\) such that the planner's marginal rate of substitution between deaths and output, evaluated at the observed sample means, equals a target VSL---following the approach of \citet{hall2020trading} and the pandemic planner analyses of \citet{alvarez2021simple} and \citet{acemoglu2021}.\footnote{The baseline \(\lambda_d = 75\) corresponds to a VSL of approximately \$6.5 million---conservative within the standard range---and Section~\ref{subsec:robustness} reports sweeps over \(\lambda_d \in \{58, 104, 127\}\), corresponding to VSLs of \$5, \$9, and \$11 million, together with alternative values of
\(\tau_b\) according to differences in age structure \citep{greenstone2020social,robinson2021benefits} }

The terminal term weights public debt alone. By the end of the horizon, widespread
vaccination has severed the link from infection to death and removed the output cost of
containment, so the only state that persists as a welfare-relevant legacy is the
accumulated fiscal burden. The terminal weight capitalizes the running debt loss into the
present value of a persistent post-horizon debt deviation
(Appendix~\ref{app:terminal}).\footnote{Alternative specifications---a Riccati-based
terminal value from the post-pandemic stationary problem and a terminal penalty on
cumulative excess mortality---are reported in Section~\ref{subsec:robustness}.} The
control-cost matrix \(\bm{R}\) plays no economic role: since fiscal costs are already
captured through the debt transition, it is set to a small uniform penalty on normalized
instrument use that keeps the problem well-posed while contributing negligibly (below 1
percent) to the objective at the reported solution.\footnote{The target variables and
transition parameters are summarized in Table~\ref{tab:planner_objects}.}

Because the model is not micro-founded, the weighted welfare solution serves as a benchmark rather than as the central welfare claim. The main analysis instead relies on the constrained policy frontier: fixing two of the three welfare-relevant dimensions at their model-implied observed level and minimizing the third, to test whether the observed policy is Pareto-dominated. These frontier problems are solved exactly, as nonlinearly constrained programs over the full control path, so no
welfare weights enter the frontier results; Section~\ref{sec:frontier} gives the formal
statement and Appendix~\ref{app:frontier_comp} the computational details.

%solved numerically using multistart nonlinear constrained optimization; I report the best feasible solution obtained across starting values.

This problem is solved at two levels: for the representative OECD economy (Section~\ref{sec:frontier}) and independently for each of the 38 countries. The structural parameters are common across all applications; only the exogenous inputs vary---the country-specific epidemiological innovations \(\varepsilon_\theta\), the initial output shock, the country intercepts \(\mu_y, \mu_b\), and the initial debt level \(b_0\).


\section{Results}
\label{sec:results}

This section presents the normative results in three steps. I first characterize the weighted planner benchmark---how the optimal policy differs from the observed average OECD response and what it delivers. I then turn to the central result, the constrained policy frontier, which tests whether the observed policy is Pareto-dominated without invoking welfare weights. Finally,
I extend the frontier analysis to the country level: solving the frontier for each of the 38 economies under their own pandemic conditions identifies which countries, given their pandemic burden, came closest to the best attainable combination of output stabilization and debt containment

% -----------------------------------------------------------------------------
\subsection{Planner Benchmark}
\label{subsec:results_benchmark}
% Inputs: fig_planner_policy_paths_baseline_aerstyle.pdf
%         fig_planner_outcome_paths_baseline_aerstyle.pdf
%         table_planner_summary_baseline.csv
% -----------------------------------------------------------------------------


Under the baseline preferences (health oriented) contrasted in Figure~\ref{fig:planner_policy}, the planner with perfect foresight\footnote{The value of this foresight is quantified in Section~\ref{subsec:robustness}}. sustains on average substantially more containment than observed until vaccination arrives (an average stringency of 48.7 against 34.1), building up at the ramping limit in the first two quarters and holding the acute-phase peak at the in-support cap---a planner sensitive to excess mortality protects lives first.\footnote{Relaxing the containment cap to the full index scale is examined in Section~\ref{subsec:robustness}.} To counter the output cost of containment, the planner deploys a slightly larger and more front-loaded above-the-line envelope (6.0 against 5.1 percent of GDP). The sharpest difference lies in liquidity support: the planner follows a similar timing profile but scales the effective volume up sevenfold (19.3 against 2.7 percent of GDP), mixing loans and guarantees and hitting the guarantee cap in two quarters.\footnote{Absent the caps, the planner would allocate the entire liquidity envelope to guarantees: both instruments carry the same output coefficient per effective unit, but guarantees enter the budget at \(\kappa_{\text{guar}} = 0.54\) against \(\kappa_{\text{loans}} = 0.84\). This is also consistent with the view that, with a functioning money market, the government should guarantee rather than originate credit.} This profile is in line with the framework of Section \ref{sec:trilemma}, where bridging liquidity over the restricted period is the key margin for keeping productive structures intact. Demand injection, by contrast, is never used despite its multiplier above one. The reason is the state-dependent multiplier: the output effect of direct income support, \(\alpha_{\text{DI}} + \alpha_{S,\text{DI}} S_k\), turns negative once stringency exceeds \(S^* = \alpha_{\text{DI}}/|\alpha_{S,\text{DI}}| \approx 36\)---a threshold the planner's containment path exceeds throughout the acute phase, and which the observed OECD average, at 34.1, only barely undercuts. Once containment is
unwound, the output gap has largely closed and demand injection has no remaining role. 

\begin{figure}[H]
    \centering
    \captionsetup{font=scriptsize}
    \includegraphics[width=0.7\textwidth]{Files/text/main/fig_planner_policy_paths_baseline_aerstyle.pdf}
    \caption{Planner benchmark and observed OECD average: policy paths.
    \textit{Notes:} Solid lines show the optimal policy under the baseline
    preferences \(\tau_b = 0.05\), \(\lambda_d = 75\) (VSL \(\approx\) \$6.5
    million; Section~\ref{sec:planner}); dashed lines show the observed OECD
    average policy. Fiscal instruments in percent of 2019 GDP; loans and
    guarantees are take-up-adjusted effective values. All instruments are subject
    to the in-support caps and the stringency ramping constraint of
    Section~\ref{sec:planner}. The pre-pandemic quarter Q4.2019 is locked at zero.}
    \label{fig:planner_policy}
\end{figure}


Figure~\ref{fig:planner_outcomes} shows the outcomes of these decisions against the model-implied benchmark. The planner achieves through the early ramp up a less severe trough and a faster recovery of the output gap, cutting the cumulative output loss by almost two thirds (\(-16.5\) against \(-46.1\) pp-quarters)---driven primarily by the front-loaded, larger liquidity envelope, complemented by a modestly larger but earlier above-the-line program. The price is fiscal: debt rises steeply during the acute phase and flattens quickly thereafter, ending 6.7 pp of GDP above the observed terminal level. This gap is the direct budgetary cost of the larger support package. The tighter containment, in turn, is what saves lives: cumulative excess mortality falls by 567 deaths per million (-23 percent).


\begin{figure}[H]
    \centering
    \captionsetup{font=scriptsize}
    \includegraphics[width=0.7\textwidth]{Files/text/main/fig_planner_outcome_paths_baseline_aerstyle.pdf}
\caption{Planner benchmark and observed OECD policy: model-implied outcome paths.
\textit{Notes:} Solid lines show outcomes under the optimal policy; dashed
lines show \emph{model-implied} outcomes under the observed OECD average
policy, i.e., the observed instruments fed through the calibrated transition
system under identical shocks. Holding the model and shock paths fixed
isolates the pure policy effect; the fit of the model-implied paths to the raw
data is documented in Section~\ref{subsec:calibration}. Both paths are
initialized at the common pre-pandemic state in Q4.2019; the debt gap in
Q1.2020 reflects the planner's front-loaded liquidity package. Baseline
preferences \(\tau_b = 0.05\), \(\lambda_d = 75\) (VSL \(\approx\) \$6.5
million).}
\label{fig:planner_outcomes}
\end{figure}



\begin{table}[H]
\centering
\scriptsize
\caption{Planner benchmark versus observed OECD policy: aggregate comparison.}
\label{tab:planner_summary}
\begin{tabular}{lccc}
\toprule
\textbf{Outcome} & \textbf{Observed (model-implied)} & \textbf{Planner} & \textbf{Difference} \\
\midrule
Cumulative output gap (pp-quarters) & $-46.1$ & $-16.5$ & $+29.6$ \\
Cumulative deaths per million       & $2{,}505$ & $1{,}938$ & $-567$ \\
Terminal debt (\% of GDP)           & $73.6$  & $80.3$  & $+6.7$ \\
Average stringency                  & $34.1$  & $48.7$  & $+14.5$ \\
Total above-the-line CP (\% GDP)    & $5.1$   & $6.0$   & $+0.9$ \\
Total below-the-line liquidity      & $2.7$   & $19.3$  & $+16.5$ \\
Total demand injection              & $2.6$   & $0.0$   & $-2.6$ \\
\bottomrule
\end{tabular}
\vspace{0.15cm}
\begin{minipage}{0.95\textwidth}
\scriptsize \textit{Notes:} Evaluation window Q1.2020--Q4.2022. Outcomes in the
first column are model-implied under the observed OECD average policy.
Below-the-line liquidity in take-up-adjusted effective terms; fiscal totals in
percent of 2019 GDP. Baseline preferences \(\tau_b = 0.05\), \(\lambda_d = 75\)
(VSL \(\approx\) \$6.5 million).
\end{minipage}
\end{table}


\noindent Table~\ref{tab:planner_summary} sums up the two margins behind the gains: the planner deploys a larger fiscal envelope---1.7 times the observed total when scaled by the containment intensity it compensates\footnote{Total fiscal support per point of average stringency: 0.52 percent of GDP for the planner against 0.31 observed.}---and shifts its composition from broad-based support toward effective and cheaper liquidity and capacity preservation, while demand injection is abandoned.


In sum, the planner benchmark displays the trilemma in action: it buys lives and output stabilization at the price of more debt, and it does so through an early ramp-up of fiscal support concentrated on capacity preservation combined
with liquidity provision. These findings are suggestive rather than conclusive, since they rest on a particular choice of welfare weights; whether the observed policy can be improved upon without taking a stand on preferences is the subject of the constrained policy frontier in the next subsection.



% -----------------------------------------------------------------------------
\subsection{Constrained Policy Frontier}
\label{sec:frontier}
% Inputs: table_constrained_frontier.csv
%         table_constrained_frontier_slacks.csv
%         table_frontier_identity.csv   (entscheidet die Formulierung!)
%         fig_frontier_comparison_baseline_aerstyle.pdf
%         table_chi_sweep.csv           (Appendix-Verweis)
% -----------------------------------------------------------------------------

The frontier analysis dispenses with welfare weights entirely. Let
\(\mathcal{L}_y(\bm{u})\), \(D(\bm{u})\), and \(b_N(\bm{u})\) denote the
discounted output loss, cumulative excess mortality, and terminal debt implied
by a policy path \(\bm{u}\) under the transition system, and let
\((\mathcal{L}_y^{\text{obs}}, D^{\text{obs}}, b_N^{\text{obs}})\) be their
values under the observed OECD policy. The three frontier programs each minimize
one dimension subject to the other two being no worse than observed:
\begin{equation}
  \min_{\bm{u} \in \mathcal{U}} \; \mathcal{L}_y(\bm{u})
  \quad \text{s.t.} \quad
  D(\bm{u}) \le D^{\text{obs}}, \;\; b_N(\bm{u}) \le b_N^{\text{obs}},
  \label{eq:frontier_program}
\end{equation}
and analogously for mortality and debt. The observed policy is Pareto-dominated if any strategy attains a strictly lower objective. Appendix~\ref{app:frontier_comp} describes the computation. Table~\ref{tab:frontier} reports the frontier results. 


\begin{table}[H]
\centering
\scriptsize
\caption{Constrained policy frontier: observed policy versus frontier solutions.}
\label{tab:frontier}
\begin{tabular}{lcccc}
\toprule
 & \textbf{Observed} & \textbf{A: min output} & \textbf{B: min deaths} & \textbf{C: min debt} \\
 & \textbf{(model-implied)} & & & \\
\midrule
Discounted output loss \(\mathcal{L}_y\) & $285.7$ & $\mathbf{60.8}$ & $285.7$ & $285.7$ \\
Cumulative output gap (pp-quarters)      & $-46.1$ & $-15.2$ & $-45.6$ & $-36.6$ \\
Cumulative deaths per million            & $2{,}505$ & $2{,}505$ & $\mathbf{1{,}947}$ & $2{,}505$ \\
Terminal debt (\% of GDP)                & $73.6$ & $73.6$ & $73.6$ & $\mathbf{67.3}$ \\
Improvement in target dimension (\%)     & --- & $78.7$ & $22.3$ & $8.6$ \\
\midrule
Average stringency                       & $34.1$ & $33.0$ & $45.1$ & $24.0$ \\
Total above-the-line (\% GDP)            & $5.1$ & $4.1$ & $1.5$ & $0.0$ \\
Total below-the-line, effective (\% GDP) & $2.7$ & $13.1$ & $10.7$ & $2.8$ \\
Total demand injection (\% GDP)          & $2.6$ & $0.0$ & $0.0$ & $0.0$ \\
\bottomrule
\end{tabular}
\vspace{0.15cm}
\begin{minipage}{0.95\textwidth}
\scriptsize \textit{Notes:} Each scenario minimizes one dimension (boldface)
subject to the other two being no worse than their model-implied values under
the observed OECD policy, Eq.~(\ref{eq:frontier_program}). The output
constraint is imposed on the discounted quadratic loss \(\mathcal{L}_y\); the
cumulative output gap is reported as a memo item. In all three scenarios both
constraints bind (satisfied with equality to numerical tolerance
\(10^{-8}\)). Evaluation window Q1.2020--Q4.2022; below-the-line liquidity in
take-up-adjusted effective terms; controls subject to the in-support caps and
the ramping constraint of Section~\ref{sec:planner}. Computation and
multistart: Appendix~\ref{app:frontier_comp}.
\end{minipage}
\end{table}

\noindent The observed OECD policy is strictly Pareto-dominated in every direction: holding the other two dimensions at their model-implied observed levels, a feasible policy path reduces the discounted output loss by 78.7 percent (scenario A), cumulative excess mortality by 22.3 percent, or 558 deaths per million (scenario B) and terminal debt by 6.3 percentage points of GDP (scenario C). In all three scenarios both constraints bind and the observed policy lies strictly in the interior of the possibility set along each dimension.

The composition of the frontier policies identifies the source of the gains, and it is the same recomposition the weighted planner selects. Scenario A raises effective liquidity support from 2.7 to 13.1 percent of GDP at slightly lower average stringency, so the output gain is almost entirely fiscal recomposition rather than reduced containment---the frontier selects precisely
the instrument with the smallest net debt footprint per unit of output identified analytically in Appendix~\ref{app:stability}. Scenario B buys the mortality reduction with tighter containment (average stringency 45.1 against 34.1) and compensates its output cost through the same channels, yielding similar results to the health oriented weighted planner. Scenario C is the purest illustration of the timing margin: it abandons above-the-line support entirely, keeps the comparatively cheap effective liquidity at its observed level, and lowers average stringency to 24.0---retiming containment into the acute waves---which improves the cumulative output gap by 9.5
pp-quarters at identical \(\mathcal{L}_y\) and identical deaths, and lets the automatic stabilizer deliver the debt reduction. 

Demand injection is zero in all three solutions, for the same state-dependence reason as in the planner benchmark. None of the three solutions requires instruments, scales, or adjustment speeds outside what OECD governments demonstrably implemented.\footnote{The dominance result does not hinge on the calibrated capacity--liquidity complementarity. Shutting it down
entirely (\(\chi = 0\)) leaves the improvements essentially unchanged, at
73.6, 22.3, and 9.0 percent; the complementarity accounts for roughly five
percentage points of the output improvement and none of the qualitative
pattern (Table~\ref{tab:chi_sweep}, Appendix~\ref{app:robustness_frontier}).
Its role is compositional: at \(\chi = 0\) the frontier policies abandon
above-the-line support altogether, so wage-subsidy-type instruments earn their
place in the optimum through their interaction with liquidity support.}

Taken together, these findings show that the weighted planner follows precisely
the strategy the frontier identifies as dominant: early front-loading with a
concentration on liquidity assistance and capacity preservation, buying output
stabilization and lives at the price of debt. It is therefore no surprise that
the weighted benchmark is robust across the preference scenarios and the VSL
sweep, as documented in Appendix~\ref{app:compl_results}.

Both results, however, are statements about the average OECD economy. In the
data, countries chose markedly heterogeneous fiscal strategies: some
implemented compositions and magnitudes close to what the average benchmark
identifies as optimal---Japan's liquidity-heavy package is the most striking
example---while others did not (Appendix~\ref{app:further_descriptives},
Table~\ref{tab:fm_per_country}). The final question is therefore whether these differences mattered: which countries, given their own pandemic burden, came closest to the best
attainable combination of output stabilization and debt containment.


\subsection{Cross-Country Results}
\label{subsec:results_countries}

The country-level analysis solves the three frontier programs of Eq.~(\ref{eq:frontier_program}) separately for each of the 38 economies. The structural parameters and the policy menu are common across countries by design---holding the technology fixed isolates differences in policy from differences in transmission---while the country-specific ingredients (epidemiological innovations, onset shocks, intercepts, initial debt) and, crucially, the constraint thresholds are each country's own: every program conditions on the country's own model-implied realized outcomes, so improvements are measured at the country's own pandemic burden. Model fit enters as a continuous diagnostic as shown in Appendix~\ref{app:country_analysis}, the frontier distances are exactly invariant to country-specific mortality-level misfit, and the ranking is rank-stable to excluding the worst-fitting quartile and further robustness checks (see Appendix~\ref{app:country_implementation}).

The headline result extends the average finding to nearly every economy: no OECD country operated on its own frontier (Figure~\ref{fig:countries_scatterp}). In 37 of 38 cases, a feasible policy
from the common menu improves each dimension at the country's own levels of the other two, with median improvements of 88 percent of the discounted output loss, 20 percent of cumulative excess mortality, and 6.1 percentage points of GDP in terminal debt. 

Figure~\ref{fig:countries_natural_units}, panel~(a), displays the composition margin (Frontier A) in natural units, holding each country's realized mortality and epidemiological shock path fixed: what each country paid in debt against the output stabilization it left unexploited at that debt. New Zealand, the United States, and the United Kingdom combine large realized debt increases (\(+21\), \(+19\), and \(+11\)~pp of GDP) with large foregone stabilization (42, 28, and 53 pp-quarters); Switzerland, Sweden, and Norway occupy the opposite corner---small debt increases, little left unexploited.

Panel~(b) prices the timing margin (Frontier C): at each country's realized output loss and mortality, how much debt was carried without corresponding returns. Denmark is the instructive case---inconspicuous in panel~(a) because its debt ratio even declined at the end of the observation period, it nonetheless carried 8.3~pp of GDP that a retimed policy would have saved at identical output and deaths; Chile, Japan, the United States, and New Zealand top this dimension with 10 to 12~pp, while T\"urkiye and Norway wasted the least (4 and 3~pp) despite very different absolute outcomes.\footnote{T\"urkiye and Ireland are special cases in the output dimension:
both recorded positive cumulative output gaps over the evaluation
window---Ireland's driven by the well-documented distortion of Irish GDP
through multinational relocations, T\"urkiye's by its high-growth,
high-inflation trajectory. Excluding both countries leaves the
ranking and the cross-sectional results unchanged
(Appendix~\ref{app:robustness}).}


%Noch given pandemic burden 
\begin{figure}[H]
    \centering
    \captionsetup{font=scriptsize}
    \begin{minipage}[t]{0.48\textwidth}
        \centering
        \includegraphics[width=\textwidth]{fig_countries_debt_vs_foregone_aerstyle.png}
        \subcaption{Panel A: Realized debt increase vs.\ foregone output stabilization}
        \label{fig:countries_debt_foregone}
    \end{minipage}\hfill
    \begin{minipage}[t]{0.48\textwidth}
        \centering
        \includegraphics[width=\textwidth]{fig_countries_outputloss_vs_wasteddebt_aerstyle.png}
        \subcaption{Panel B: Realized output loss vs.\ wasted debt}
        \label{fig:countries_outputloss_wasted}
    \end{minipage}
    \caption{Frontier distances in natural units.
    \textit{Notes:} Each point is one country; horizontal axes show realized
    outcomes from the raw data over Q1.2020--Q4.2022, vertical axes the
    unexploited improvement at the country's own frontier. Panel~(a):
    realized change in the debt-to-GDP ratio against the output stabilization
    attainable at the country's own mortality and debt, in pp-quarters.
    Panel~(b): realized cumulative output gap against the terminal debt
    reduction attainable at the country's own output loss and mortality
    (scenario~C), in pp of GDP. Reading: countries toward the upper right
    paid much and left much unexploited; positions near the origin reflect
    little room to convert at the realized position, not necessarily
    favorable outcomes. Italy omitted (Appendix~\ref{app:country_special}).}
    \label{fig:countries_natural_units}
\end{figure}


\noindent Combining the two margins yields the frontier distance: a single measure of how much a country left unexploited---in output stabilization and in debt---at its own realized position. It summarizes the quality of a strategy's composition and timing, conditional on the trilemma position the country chose, and it is the ranking of Table~\ref{tab:country_ranking} (For a visual representation see Appendix \ref{app:country_implementation} Figure \ref{fig:countries_scatterp}).


Ranking countries by their distance to their own frontier, Mexico, Sweden,
Portugal, Finland, France, and Norway come closest: given their fiscal space
and pandemic pressure, these countries chose the compositions and timing
that left the least output stabilization and debt reduction unexploited. New
Zealand, the United States, Australia, Japan, and Denmark are furthest
(Table~\ref{tab:country_ranking}). The measure is conditional: it identifies
the combination of policies that was least Pareto-improvable given the
country's own realized trilemma position, not the best absolute outcomes.
Mexico illustrates this exactly: its outcomes were poor, but its minimal
fiscal deployment left almost no room for improvement at its own debt
level---any material gain would have required accepting more debt, which the
conditional measure holds fixed. Another special case is Italy. It exceeds
the ramping limit pivotally: hit first and hardest, it tightened containment
faster than any other OECD government, at a speed the common menu does not
admit. Raising the ramping limit to the 99th percentile of observed
transitions renders Italy's programs feasible; under that relaxed menu, its
improvements (47 percent of output gap, 5.4 pp of debt) would place
it among the five countries closest to their frontier, between Portugal and
Finland---consistent with its efficient, liquidity-led composition.

\begin{table}[H]
\centering
\scriptsize
\caption{Distance to the country-specific frontier: closest and furthest countries.}
\label{tab:country_ranking}
\begin{tabular}{lc cc !{\vrule width 0.4pt} c}
\toprule
 & \textbf{Frontier distance} & \multicolumn{2}{c}{\emph{Ranking dimensions}} & \emph{Context} \\
\cmidrule(lr){3-4}
 & (avg.\ rank, 1--37) & \textbf{Output impr.} & \textbf{Debt impr.} & \textbf{Deaths impr.} \\
 & & (\% of own \(\mathcal{L}_y\)) & (pp of GDP) & (\% of own deaths) \\
\midrule
\multicolumn{5}{l}{\emph{Closest to own frontier}} \\
\quad Mexico    & $1.0$  & $35.3$ & $2.2$  & $16.6$ \\
\quad Sweden    & $4.0$  & $58.5$ & $3.7$  & $21.3$ \\
\quad Portugal  & $5.0$  & $48.7$ & $5.0$  & $16.2$ \\
\quad Finland   & $6.0$  & $62.3$ & $3.7$  & $13.3$ \\
\quad France    & $7.0$  & $52.1$ & $5.6$  & $13.3$ \\
\quad Norway    & $7.5$  & $77.9$ & $3.0$  & $12.5$ \\
\midrule
\multicolumn{5}{l}{\emph{Furthest from own frontier}} \\
\quad Chile       & $31.5$ & $92.9$  & $11.7$ & $27.4$ \\
\quad Denmark     & $32.0$ & $100.0$ & $8.3$  & $20.6$ \\
\quad Australia   & $33.5$ & $100.0$ & $8.9$  & $28.6$ \\
\quad Japan      & $33.5$ & $99.6$  & $10.8$ & $21.2$ \\
\quad New Zealand & $34.5$ & $100.0$ & $10.2$ & $38.7$ \\
\quad United States & $35.0$ & $100.0$ & $10.6$ & $34.6$ \\
\bottomrule
\end{tabular}
\vspace{0.15cm}
\begin{minipage}{0.95\textwidth}
\scriptsize \textit{Notes:} Frontier distance is the average of the country's
within-sample ranks in the output and debt improvements (37 ranked countries;
lower = closer to the own frontier). Each improvement column reports a
separate frontier program: the output and debt improvements hold the
country's mortality (and the respective other dimension) at its observed
level; the deaths improvement is from the program that instead minimizes
mortality at the observed output loss and debt---it is contextual and does
not enter the distance. Debt improvements are reported in absolute pp of GDP
rather than in percent of the observed debt increase. Italy is excluded (observed
policy not replicable within the menu, see text). Full 38-country table:
Appendix~\ref{app:country_tables}.
\end{minipage}
\end{table}


\noindent Relating the distances to the observed packages (Appendix \ref{app:country_analysis}  Table~\ref{tab:country_composition}) makes the ranking concrete. The frontier-near countries share a recognizable profile: liquidity-oriented, front-loaded packages with little demand injection with, on average, a liquidity share of 34 percent against 22 among the furthest, a demand-injection share of 24 percent against 35, and 84 against 75 percent of the package committed by Q3.2020. Sweden is the archetype: a package with a liquidity share of 34 percent, a demand-injection share of 13 percent, and 96 percent committed early. Norway pushed the composition furthest---half its package in liquidity, essentially no demand injection---and realized one of the highest returns in the panel. Switzerland shows that composition alone does not decide the ranking: with a nearly demand-injection-free, front-loaded package and one of the highest returns, it nonetheless sits in the midfield. Its distance is a scale margin, not a composition margin: although its realized debt increase was among the smallest---guarantees enter gross debt only upon being called, and the constitutional debt brake disciplined the envelope---the consolidation and stabilizer returns its policy left unused imply that four fifths of its (small) remaining loss was still removable at
that same terminal debt. Low realized debt is not penalized; the measure holds it fixed and asks only whether it was fully converted.

%%Nicht ganz zufrieden damit-> die ganze frontier thematik ist nur über die effizient also relativ zum verlust aber das ist auch richtig so, sonst ist es nicht vergleichbar
%returns erst weiter unten-> maybe zusammennehmen? was haben sie gemacht und returns

At the other end, there are two ways to be far from the frontier. The first is unfavorable composition and timing: the United States placed half its package in demand injection, committed
only 54 percent of it by Q3.2020, and gained little output for the debt it added; Chile combined the highest demand-injection share of the panel (0.67) with the most back-loaded profile (40 percent committed early) and realized a negative return; Australia and New Zealand concentrated their support in wage-subsidy commitments beyond the common menu while holding liquidity shares below a tenth of their packages. The second is unexploited fiscal space despite reasonable composition: Denmark and Japan ran liquidity-heavy, well-timed packages---Japan's is the most planner-like in the panel---but the fiscal space their realized debt paths left unused, combined with the output ceiling, implies that far more was attainable at their own debt levels than their scale delivered. The ranking is grounded in the realized magnitudes: Table~\ref{tab:country_raw_outcomes} juxtaposes each country's frontier distance with its raw outcomes and shows that the efficiency pattern is not an artifact of the model-implied thresholds---Sweden and the United States, for instance, suffered nearly identical raw output losses, but Sweden at one tenth of the debt increase and with fewer deaths.

%5Exchange rate sauberer verkaufen??
What separates the two ends is captured in a single summary statistic from the data: the realized fiscal exchange rate---the cumulative output gained per percentage point of fiscally induced debt, computed against a counterfactual that keeps each country's observed containment path and sets the fiscal instruments to zero yielding the return on stabilization per pp of debt. The median country earned 2.8 pp-quarters per point of debt; the high returns above belong to Norway (3.7),
Switzerland (3.6), and Sweden (3.1), the low ones to the United States (1.1)
and Chile (\(-0.2\)) Appendix~\ref{app:country_analysis}, Table~\ref{tab:country_composition} reports the full panel).



\begin{figure}[H]
    \centering
     \captionsetup{font=scriptsize}
    \includegraphics[width=0.6\textwidth]{Files/text/main/fig_countries_exchange_aerstyle.pdf}
    \caption{Realized fiscal exchange rate and demand-injection share.
    \textit{Notes:} Each point is one country. The exchange rate is the
    cumulative output gain of the observed fiscal package per percentage
    point of fiscally induced debt, computed against a counterfactual that
    keeps the country's observed containment path and sets all fiscal
    instruments to zero. Cross-country correlation with the observed
    demand-injection share: \(-0.89\). Solid line: linear fit; only selected
    countries are labeled. The four countries with a negative realized
    exchange rate (Chile, Colombia, Israel, Mexico) are those with the
    highest demand-injection shares.}
    \label{fig:countries_exchange}
\end{figure}


\noindent The gradient across countries is almost entirely a demand-injection gradient: the exchange rate correlates at \(-0.89\) with the observed demand-injection share, and it is negative in exactly the four countries with the highest such shares (Chile, Colombia, Israel, and Mexico, at 56 to 67 percent of their packages; Figure~\ref{fig:countries_exchange}).

This is the cross-country footprint of the state-dependent multiplier of Section~\ref{sec:estimation}: demand injection deployed under binding containment adds debt without adding output---the same mechanism that drives the planner's abandonment of the instrument, visible here in the realized data rather than in the optimization.

Across all three levels of the analysis, the same answer emerges. The
weighted planner, the weight-free frontier, and the cross-country comparison
identify the identical dominant strategy for confronting the trilemma:
fiscal support that is front-loaded and recomposed toward effective
liquidity provision and capacity preservation, containment concentrated in the acute waves, and demand injection withheld while restrictions bind. Countries whose packages came
closest to this composition as Switzerland, Norway and Sweden realized the highest output return per unit of debt; those that relied on scale, late commitment, or demand injection paid
for output they did not receive. No country's policy---whatever its size---reached its own frontier: the binding constraint on pandemic policy was not the fiscal envelope but its composition and timing.


\section{Conclusion}
\label{sec:conclusion}


%%DRAFT
This dissertation asked whether the fiscal and containment strategies
observed during COVID-19 were efficient, given that pandemic policy is
constrained by a trilemma between excess mortality, output stabilization, and
public debt. The answer, within the estimated transition system, is no. The
observed OECD response was strictly Pareto-dominated: a feasible policy
path---using no instrument, scale, or adjustment speed beyond what OECD
governments demonstrably implemented---would have reduced the discounted
output loss by 79 percent, excess mortality by 22 percent, or terminal debt
by 6.3 percentage points of GDP, holding the other dimensions at their
observed levels, and no individual economy operated on its own frontier. The
trilemma itself, however, is binding: debt is dynamically unstable and cannot
be stabilized by any instrument, no fiscal measure is self-financing, and
containment protects lives only at an output cost. The gains therefore come
not from escaping the trade-off but from navigating it efficiently---through
composition and timing rather than volume. Front-loaded liquidity provision and capacity preservation, containment concentrated in the acute waves, and the abandonment of demand
injection under binding restrictions recur at every level of the analysis: in the weighted planner, in the weight-free frontier, and in the cross-country comparison, where the countries closest to their own frontier ran precisely such packages and where the output return per unit of debt
falls sharply---and eventually turns negative---in the demand-injection share of the observed package.

For future crises that share the pandemic's structure---binding supply and mobility constraints rather than deficient demand---these results translate into a compositional heuristic rather than a spending target. A government facing such a shock should build the liquidity architecture first and early: guarantee and lending facilities preserve productive structures through the restricted period at the lowest budgetary cost per unit of output, and their effect operates through the outstanding stock, which rewards anticipation. Wage-subsidy-type instruments earn their place primarily in combination with that liquidity channel. Demand injection, by contrast, should remain off while restrictions bind, because its output effect is state-dependent and turns negative under tight containment---the instrument is not wasteful in general, but mistimed when deployed as crisis relief during lockdowns. Containment itself is best ramped early and concentrated in the waves, where its epidemiological return is highest in line with the findings of \citep{acemoglu2021}. None of this requires fiscal capacity beyond what OECD governments actually mobilized; it requires deploying it differently.

These conclusions carry three limitations. First, the transition parameters are reduced-form panel estimates, not
causally identified deep parameters, and are therefore exposed to the Lucas
critique when evaluated under counterfactual policies. The exposure is
limited: the policy menu is restricted to the observed support, so the
counterfactuals interpolate within---rather than extrapolate beyond---the
variation that identifies the parameters, and the exercise is a
retrospective, episode-specific evaluation, for which the critique carries
less force than for prospective rule design. Above all, the two mechanisms
that drive the normative results---the state-dependent demand-injection
multiplier and the dominance of liquidity-type support under
containment---emerge as equilibrium properties of the micro-founded models
discussed in Section~\ref{sec:trilemma}
\citep{guerrieri2022,bayer2023c,fariaecastro2021}, in which behavior
is derived from optimization and the critique does not apply by
construction. The transition system reproduces this structural policy
ranking in reduced form; what it adds is the joint quantification with the
debt side of the trilemma. Second, the planner solves under perfect foresight of the
realized shock paths. The no-anticipation variant, which fixes the
pre-shock quarter at the observed policy, prices this foresight at roughly
a quarter of the weighted planner's welfare gain and shows that the
dominance result does not rest on it: all three frontier improvements
remain strictly positive, with unchanged policy composition
(Appendix~\ref{app:compl_results}).
%% NOTE: NA-Levels erst nach konsistentem Re-Run (eps_y(Q1.20)-Check)
%% in den Appendix uebernehmen.
Third, the model is estimated on aggregate country-quarter data: it is
silent on the distributional incidence of the instruments and on
within-country heterogeneity, and the common transition technology---the
design choice that makes cross-country policy differences comparable in the
first place---abstracts from genuine differences in fiscal multipliers and
health systems, which the fit diagnostics flag rather than resolve.
 
The paper's contribution is a framework, a dataset, and a benchmark. The
trilemma formalizes why pandemic policy cannot be understood as standard
stabilization policy; the measure-level database of 2,301 fiscal
interventions makes the composition of the pandemic fiscal response
observable; and the estimated planner problem converts both into a
model-implied standard against which realized strategies can be judged.
The natural next steps follow from the limitations: micro-founding the
capacity-preservation channel, extending the calibration to
country-specific mortality and take-up profiles, and embedding real-time
uncertainty to ask how much of the frontier a planner without foresight
can reach. The retrospective answer, however, already stands: in a
pandemic, the binding constraint on fiscal policy is not how much
governments can spend, but how well what they spend matches the state of
the crisis.
%%%OLD
\newpage




\bibliographystyle{plainnat}
\bibliography{references}


\floatbarrier

\newpage
\section{Appendix}

\appendix

%==============================================================
\section{Data Appendix}
\label{app:data}

All used materials and data are available on the corresponding GitHub page \href{https://github.com/pesentiA/Pandemic-Trilemma}
{\textit{Pandemic Trilemma}}. The exact location of the following main elements are summarized in the following table:
\begin{table}[H]
\centering
\caption{Replication Files and Documentation}
\label{tab:replication_files}
\renewcommand{\arraystretch}{1.25}

\begin{tabular}{@{}ll@{}}
\toprule
\textbf{Component} & \textbf{Repository Location} \\
\midrule
README
&
\href{https://github.com/pesentiA/Pandemic-Trilemma}
{\texttt{README.md}}
\\

Fiscal Measures Codebook
&
\href{https://github.com/pesentiA/Pandemic-Trilemma/tree/main/Files/text/codebook}
{\texttt{Files/text/codebook}}
\\

Data
&
\href{https://github.com/pesentiA/Pandemic-Trilemma/tree/main/Files/data}
{\texttt{Files/data}}
\\

R Code
&
\href{https://github.com/pesentiA/Pandemic-Trilemma/tree/main/Files/code/R}
{\texttt{Files/code/R}}
\\

MATLAB Code
&
\href{https://github.com/pesentiA/Pandemic-Trilemma/tree/main/Files/code/matlab}
{\texttt{Files/code/matlab}}
\\

Figures and Tables
&
\href{https://github.com/pesentiA/Pandemic-Trilemma/tree/main/Files/output}
{\texttt{Files/output}}
\\



\bottomrule
\end{tabular}
\end{table}




%% ============================================================
\section{External Validation: The Trilemma in Comparative Perspective}
\label{sec:comparison}
%% ============================================================

To establish that the pandemic trilemma is not simply a generic feature of crisis management, I apply the four constitutive properties from Section~\ref{sec:trilemma} to four alternative crisis types: armed conflict, natural disasters, financial crises, and climate change. The purpose is diagnostic. A crisis satisfies the pandemic-style trilemma only if the threat is self-reinforcing, the suppression instrument is contractionary, fiscal compensation generates debt, and the output loss feeds back into public debt through automatic stabilizers. Table~\ref{tab:comparison} summarizes the comparison.

\begin{landscape}

\begin{table}[p]
\centering
\caption{Diagnostic Comparison of Crisis Types}
\label{tab:comparison}

\renewcommand{\arraystretch}{1.45}
\setlength{\tabcolsep}{6pt}

\begin{threeparttable}

\begin{tabular}{@{}P{5.4cm}
                    C{2.7cm}
                    C{2.7cm}
                    C{2.7cm}
                    C{2.7cm}
                    C{4.2cm}@{}}
\toprule
&
\textbf{Pandemic}
&
\shortstack{\textbf{Armed}\\\textbf{conflict}}
&
\shortstack{\textbf{Natural}\\\textbf{disaster}}
&
\shortstack{\textbf{Financial}\\\textbf{crisis}}
&
\shortstack{\textbf{Climate}\\\textbf{change}}
\\
\midrule

Exogenous self-reinforcing threat
& \checkmark
& Partial
& $\times$
& $\times$
& \checkmark
\\

Contractionary suppression instrument
& \checkmark
& $\times$
& n/a
& $\times$
& Partial
\\

Debt-generating fiscal compensation
& \checkmark
& \checkmark
& \checkmark
& \checkmark
& \checkmark
\\

Endogenous fiscal feedback
& \checkmark
& Partial
& Partial
& \checkmark
& \checkmark
\\

\midrule

\textbf{Pandemic-style trilemma}
& \checkmark
& $\times$
& $\times$
& $\times$
& \shortstack{Analogous, but not\\directly comparable}
\\

\bottomrule
\end{tabular}

\begin{tablenotes}[flushleft]
\footnotesize
\item \textit{Notes:}
\checkmark{} indicates that the property is satisfied;
$\times$ indicates that the property is not satisfied;
\textit{Partial} indicates that the property is satisfied only in a restricted or non-structural sense;
and \textit{n/a} denotes that the property is not applicable.
Climate change is structurally closest to the pandemic case but differs in its time scale,
observability, and terminal condition.
\end{tablenotes}

\end{threeparttable}
\end{table}

\end{landscape}

%% --- Armed Conflict ---
\subsection*{Armed Conflict}

Armed conflict shares some surface-level features with the pandemic. It can generate fiscal pressures through military expenditure and may trigger automatic stabilizers if economic activity is disrupted. However, it does not satisfy the full trilemma structure.

Property~1 is at most partially satisfied. While armed conflicts can escalate, this escalation is not governed by an autonomous biological growth process comparable to viral prevalence with $\rho_\theta>1$. Military escalation is mediated by strategic decisions, military capacity, and political choices, rather than by a mechanical reproduction process that compounds independently of policy.

The main difference is Property~2. The instruments used to address the threat---military mobilization, defense production, and related public spending---are not inherently contractionary in the pandemic sense. They often increase demand, absorb idle resources, and raise employment. The instrument that addresses the threat therefore works with, rather than against, short-run economic stabilization. This breaks the circular mechanism that defines the pandemic trilemma. In addition, war often destroys physical capital, whereas the pandemic primarily suppresses economic flows while leaving much of the capital stock intact. The post-war problem is therefore a reconstruction problem; the pandemic problem is a stabilization problem under binding containment.

%% --- Natural Disasters ---
\subsection*{Natural Disasters}

Natural disasters are the clearest example of an exogenous shock, but they fail the self-reinforcing-dynamics condition. An earthquake, flood, or storm strikes, causes damage, and then dissipates. The shock does not grow endogenously with a persistence parameter above unity. There is therefore no equivalent of viral prevalence that must be suppressed to prevent future growth.

As a result, Property~2 is not applicable. The planner does not need to impose a contractionary suppression instrument to stop the shock from reproducing. The relevant policy problem is reconstruction: restoring infrastructure, housing, and productive capacity after the event. Reconstruction spending may generate debt and output losses may activate automatic stabilizers, so Properties~3 and~4 can be present. But without a self-reinforcing threat and without a contractionary suppression instrument, the circular trilemma structure does not arise.

%% --- Financial Crisis ---
\subsection*{Financial Crises}

Financial crises fail Property~1 because the shock is endogenous to the economic system. It emerges from leverage, balance-sheet fragility, asset-price misalignment, or institutional failure within the financial sector. Even when crises propagate internationally, the transmission operates through financial linkages and market expectations rather than through an external biological process.

Property~2 is also violated. The main crisis-fighting instruments---liquidity provision, bank recapitalization, lender-of-last-resort interventions, and asset purchases---are stabilizing or expansionary for the real economy. They are designed to prevent a collapse in credit, demand, and employment. This reverses the pandemic logic. In a pandemic, the instrument that suppresses the threat directly reduces output. In a financial crisis, the instrument that addresses the threat is intended to prevent output from falling further. Subsequent austerity or fiscal consolidation may be contractionary, but these are responses to the fiscal legacy of the crisis, not the crisis-fighting instruments themselves. The pandemic-style trilemma therefore does not emerge.

%% --- Climate Change ---
\subsection*{Climate Change}

Climate change is the closest structural analogue to the pandemic trilemma, but it is not directly comparable.

Property~1 is broadly satisfied. Greenhouse gas accumulation is an external process with self-reinforcing features: emissions accumulate in the atmosphere, feedback loops can amplify warming, and tipping points may make the process harder to reverse. The health and welfare costs are also real, operating through heat stress, extreme weather, food-system disruption, and other climate-related damages.

Property~2 is partially satisfied. Climate mitigation can be contractionary for carbon-intensive sectors and may impose transitional costs on output, employment, and relative prices. However, unlike pandemic containment, climate policy does not mechanically shut down broad economic activity. Mitigation also induces substitution, innovation, and sectoral reallocation. The suppression instrument is therefore not contractionary in the same immediate and economy-wide sense as pandemic containment.

Properties~3 and~4 are present. The transition to a low-carbon economy can require substantial fiscal compensation, public investment, and support for affected households and firms. These measures can generate public liabilities. If mitigation or transition costs depress output, automatic stabilizers can further worsen the fiscal position. In this sense, climate change contains an analogous debt-output-policy tension.

The differences are nevertheless decisive. First, the relevant time scale is measured in decades rather than quarters. This turns the discount factor from a technical modeling choice into a normative question of intergenerational welfare. Second, the health consequences are diffuse, statistically attributed, and temporally remote, whereas pandemic mortality is immediate and politically salient. Third, there is no terminal condition analogous to vaccination. In the pandemic, vaccination reduces the fatality and transmission constraints and thereby dissolves the trilemma over a finite horizon. Climate change has no comparable exogenous terminal date; mitigation must be sustained over a much longer horizon.

Climate change therefore represents a potential analogue rather than a direct application of the pandemic framework. A climate trilemma would require different assumptions about the planning horizon, discounting, technological change, and terminal conditions.

%% --- Conclusion ---
\vspace{5mm}

\noindent The comparison shows that the pandemic case is structurally distinct. Armed conflict, natural disasters, and financial crises each activate parts of the fiscal trade-off, but they do not combine a self-reinforcing external threat with a contractionary suppression instrument. Climate change is the closest analogue, because mitigation can generate output costs and requires fiscal compensation, but its horizon, observability, and absence of a vaccine-like terminal condition make it a different planning problem.

The pandemic therefore provides a uniquely tractable setting in which all four properties bind simultaneously within a finite planning horizon. This is why the pandemic trilemma can be modeled as a dynamic allocation problem over health, output, and public debt, while the comparison cases either lack one of the constitutive properties or require a substantially different modeling framework.

AB HIER: DURCH DEN TEXT GEHEN; WAS BRAUCHE ICH WO?

\section{Further Descriptives}
\label{app:further_descriptives}

\subsection{Fiscal Measures per Category per Country}

\begin{table}[H]
\centering
\caption{\label{tab:tab:fiscal_sums_country}Fiscal Policy Sums by Country}
\label{tab:fm_per_country}
\centering
\resizebox{\ifdim\width>\linewidth\linewidth\else\width\fi}{!}{
\fontsize{8}{10}\selectfont
\begin{threeparttable}
\begin{tabular}[t]{lrrrrrl}
\toprule
Country & Capacity Preservation & Demand Injection & Health & Above & Loans & Guarantees\\
\midrule
AUS & 11.94 & 4.09 & 0.59 & 10.17 & 0.76 & 1.01\\
AUT & 9.88 & 1.19 & 1.68 & 7.60 & 0.29 & 1.99\\
BEL & 16.75 & 1.70 & 1.79 & 5.60 & 0.22 & 10.93\\
CAN & 9.40 & 6.47 & 1.71 & 5.24 & 0.63 & 3.53\\
CHE & 11.17 & 0.18 & 0.71 & 4.91 & 0.14 & 6.11\\
\addlinespace
CHL & 13.55 & 11.89 & 0.38 & 2.68 & 1.62 & 9.20\\
COL & 3.92 & 2.64 & 0.50 & 1.18 & 0.54 & 2.20\\
CRI & 7.65 & 1.43 & 0.06 & 0.39 & 4.77 & 2.48\\
CZE & 20.34 & 2.54 & 1.38 & 4.93 & 0.09 & 15.32\\
DEU & 27.61 & 2.63 & 1.70 & 7.12 & 3.10 & 17.39\\
\addlinespace
DNK & 22.40 & 4.77 & 0.83 & 5.47 & 13.26 & 3.67\\
ESP & 19.74 & 1.89 & 0.08 & 3.96 & 2.80 & 12.98\\
EST & 16.15 & 0.48 & 1.68 & 7.98 & 2.59 & 5.59\\
FIN & 9.35 & 2.32 & 1.86 & 2.40 & 1.43 & 5.53\\
FRA & 22.38 & 1.21 & 0.39 & 7.05 & 0.84 & 14.48\\
\addlinespace
GBR & 24.25 & 2.09 & 4.73 & 10.09 & 0.92 & 13.24\\
GRC & 15.23 & 2.48 & 0.42 & 7.54 & 5.77 & 1.42\\
HUN & 9.21 & 1.69 & 4.48 & 3.60 & 1.36 & 4.25\\
IRL & 7.18 & 4.14 & 1.54 & 4.73 & 1.13 & 1.31\\
ISL & 7.27 & 3.54 & 0.03 & 4.73 & 1.08 & 1.46\\
\addlinespace
ISR & 5.75 & 5.72 & 1.65 & 2.90 & 1.42 & 1.42\\
ITA & 29.10 & 0.93 & 1.32 & 7.65 & 1.79 & 19.66\\
JPN & 29.79 & 5.34 & 1.90 & 4.68 & 13.18 & 11.94\\
KOR & 10.87 & 2.58 & 0.76 & 2.48 & 4.93 & 3.46\\
LTU & 7.83 & 3.52 & 0.21 & 4.68 & 1.03 & 2.12\\
\addlinespace
LUX & 8.52 & 0.11 & 0.40 & 3.85 & 0.24 & 4.43\\
LVA & 12.75 & 0.25 & 0.03 & 5.49 & 2.00 & 5.27\\
MEX & 0.88 & 1.01 & 0.09 & 0.41 & 0.24 & 0.23\\
NLD & 9.24 & 1.15 & 0.09 & 5.45 & 0.39 & 3.40\\
NOR & 6.69 & 0.12 & 0.30 & 1.62 & 1.52 & 3.55\\
\addlinespace
NZL & 13.60 & 4.99 & 1.99 & 10.21 & 1.24 & 2.15\\
POL & 10.14 & 1.90 & 0.33 & 6.60 & 0.29 & 3.25\\
PRT & 12.49 & 1.41 & 0.00 & 2.93 & 1.50 & 8.05\\
SVK & 7.64 & 0.13 & 0.08 & 5.11 & 2.45 & 0.08\\
SVN & 16.97 & 0.72 & 2.31 & 11.62 & 0.63 & 4.73\\
\addlinespace
SWE & 7.76 & 0.70 & 0.75 & 2.87 & 1.72 & 3.16\\
TUR & 10.87 & 0.50 & 0.45 & 2.56 & 0.02 & 8.29\\
USA & 11.74 & 8.95 & 4.66 & 6.03 & 2.85 & 2.86\\
\bottomrule
\end{tabular}
\begin{tablenotes}
\item \textit{Note: } 
\item Notes: Values are summed over the full sample period. CP = capacity preservation, DI = demand injection, H = health-related spending. Loans and guarantees refer to capacity-preservation below-the-line instruments.
\end{tablenotes}
\end{threeparttable}}
\end{table}





% =============================================================================
%  CONSOLIDATED APPENDIX: Analytical Framework and iLQR Derivation
%  Replaces: "Analyticial Analysis" (old, single-channel)
%            + "Derivation of the Optimal Control Law" (generic LQR)
%            + "iLQR Derivation for the Augmented System" (correct but partial)
%  Aligned with: main.tex Section 3 and pandemic_sequentialALV6.m
%  Last updated: April 2026
% =============================================================================
\section{Identification Strategy}
\label{app:identification}

\subsection{Identification Concerns and Dynamic-Panel Interpretation}
\label{app:concerns}

Three features of the pandemic setting complicate the interpretation of the
estimated transition parameters. First, containment and fiscal support may both
respond to pandemic severity: countries experiencing worse outbreaks may impose
stricter containment measures and deploy larger fiscal packages, generating a
correlation between policy variables and unobserved shocks to output. Second,
containment $S_{ik}$ and fiscal support are implemented simultaneously in the
same country-quarter units. This creates a multiple-treatment setting in which
estimating the conditional contribution of one instrument requires conditioning
on the other. Third, treatment intensity is continuous, time-varying, and
global, which rules out standard $2\times2$ difference-in-differences designs.

The maintained assumption is therefore not strict exogeneity, but conditional
sequential exogeneity in a dynamic panel setting. The output equation includes
the lagged dependent variable $y_{i,k-1}$, which is required by the
transition-system interpretation of the model. With a lagged dependent
variable, strict exogeneity is generally not available because innovations to
current output enter future values of the dependent variable. The relevant
assumption is instead that, conditional on country fixed effects, lagged
economic states, observed pandemic conditions, and the timing structure implied
by the model, the remaining output innovation is not systematically correlated
with the predetermined policy variables entering the equation.

This interpretation is consistent with the timing of the fiscal variables.
Above-the-line capacity preservation enters with a two-quarter lag,
$F^{\text{above}}_{i,k-2}$, while demand injection enters with a one-quarter
lag, $F^{\text{DI}}_{i,k-1}$. These measures were typically announced and
implemented before the corresponding quarterly GDP outcomes were observed, and
often before those outcomes were publicly available. Below-the-line support
enters through the cumulative liquidity stock $K^{\text{below}}_{ik}$,
reflecting the availability of liquidity rather than a one-period expenditure
flow. The lag structure therefore aligns the empirical specification with the
economic transmission channels derived in the theoretical framework and limits,
although it cannot fully eliminate, reverse causality from realized output to
fiscal deployment.

The data provide useful variation for this parameterization strategy. Containment intensity and fiscal deployment are only weakly correlated over the sample period ((r=0.04)), and fiscal responses vary substantially more across countries than containment intensity (see Appendix~\ref{sec:correlation_matrix} for the full discussion). This mitigates, but does not eliminate, the concern that stricter containment mechanically induces larger fiscal compensation. Countries facing broadly similar pandemic conditions adopted markedly different fiscal strategies, particularly in the composition of support across demand injection, above-the-line capacity preservation, and below-the-line liquidity instruments. This compositional heterogeneity is the main source of variation used to discipline the channel-specific transition parameters.

A direct way to assess whether cross-country selection on crisis severity
contaminates the estimates is to decompose the fiscal coefficients into their
within- and between-country components following \citet{mundlak1978}. The
within and between estimators differ in sign for the key fiscal channels: the
between estimator is dominated by the fact that countries with deeper
contractions deployed larger fiscal packages, so that the cross-country
covariation between fiscal support and output is contaminated by selection on
crisis severity. The within estimator, which exploits only the
timing and composition of fiscal deployment within a country over the pandemic
trajectory, is free of this contamination. This sign reversal is precisely the
pattern that motivates the fixed-effects specification: it demonstrates that the
identifying variation must come from within-country differences in fiscal
composition rather than from cross-country comparisons of total support. Full
decomposition results are reported in Section~\ref{sec:robustness}.

Several remaining concerns are addressed directly in robustness checks. First, voluntary reductions in economic activity may not be fully captured by mandated stringency (\(S_{ik}\)). I therefore add specification where I include excess mortality (\(d_{ik}\)) to capture the fear channel implied by the theoretical framework, namely behavioral responses to observed mortality. The coefficient on containment is consequently interpreted as the conditional output association of mandated restrictions given the observed health state. Empirically, the results are not sensitive to this choice: (\(d_{ik}\)) is not statistically significant, and excluding it leaves the coefficients of interest largely unchanged. I therefore retain it in the baseline for theoretical consistency and report specifications without the fear term as robustness checks. Mobility-based measures of activity provide an additional robustness check for the containment channel.

Second, vaccination may confound the estimated relationships if countries that vaccinated faster also adjusted fiscal composition differently. I therefore control for vaccination dynamics where relevant and report robustness checks restricting the sample to the pre-vaccination period. Third, the classification of fiscal measures into demand injection and capacity preservation involves judgment. The classification protocol mitigates this concern by assigning measures according to their primary transmission mechanism, and robustness checks under alternative classification rules assess whether the results are driven by borderline cases. Fourth, cross-country spillovers through trade and supply chains may generate correlated shocks across countries. Common shocks are addressed through robustness specifications with alternative time controls, while asymmetric spillovers remain a maintained limitation. Robustness checks controlling for trade openness assess the quantitative relevance of this concern.

Inference is based on cluster-robust HC1 standard errors clustered at the country level, allowing for arbitrary heteroskedasticity and serial correlation within countries over time. Since the number of country clusters is moderate (\(N=38)\), I additionally report wild cluster bootstrap (p)-values clustered at the country level as a small-sample robustness check. These inference procedures address the precision of the estimates but do not correct the dynamic-panel bias generated by the inclusion of (\(y_{i,k-1}\)) together with country fixed effects.


The inclusion of the lagged output gap follows directly from the
transition-system interpretation of the model. Pandemic output losses were
persistent, and the planner problem requires an estimate of this persistence
because current output gaps affect future states. Omitting $y_{i,k-1}$ would
force the fiscal and containment variables to absorb part of this persistence
and could therefore distort the parameterization of the output transition.
Since the output equation~\eqref{eq:output_est} includes a lagged dependent
variable together with country fixed effects in a short panel, the within
estimator of $\rho_y$ is subject to Nickell bias of order $1/T$. I therefore
report bias-corrected dynamic-panel estimates in Section~\ref{sec:robustness}.
The persistence-channel results are robust to these corrections.

The estimates should therefore be interpreted as conditional transition parameters under the maintained timing and sequential-exogeneity assumptions rather than as randomized treatment effects. Their role is to discipline the model quantitatively and to assess whether the signs, timing, and relative magnitudes implied by the transition system are consistent with the data.

\subsection{Correlation Matrix and Variation Decomposition}
\label{sec:correlation_matrix}


\begin{table}[H]
\centering
\caption{Within vs.\ Between Variation Decomposition}
\label{tab:within_between}
\scriptsize
\renewcommand{\arraystretch}{1.15}
\begin{tabular}{@{} l r r r r r @{}}
\toprule
\textbf{Variable} & \textbf{Overall SD} & \textbf{Between SD} & \textbf{Within SD}
  & \textbf{Within (\%)} & \textbf{Between (\%)} \\
\midrule
$y_{ik}$ (pp of potential)       & 4.766 & 3.018 & 3.688 & 77.4 & 63.3 \\
$S_{ik}$ (0--100)                & 18.696 & 6.143 & 17.658 & 94.4 & 32.9 \\
$F^{CP}_{ik}$ (pp GDP)          & 3.032 & 0.691 & 2.953 & 97.4 & 22.8 \\
$F^{DI}_{ik}$ (pp GDP)          & 0.705 & 0.251 & 0.659 & 93.5 & 35.6 \\
$F^{H}_{ik}$ (pp GDP)           & 0.381 & 0.123 & 0.361 & 94.7 & 32.3 \\
$d_{ik}$ (P-score, \%)          & 14.873 & 7.603 & 12.783 & 85.9 & 51.1 \\
$\theta_{ik}$ (\% of pop.)      & 1.889 & 0.533 & 1.813 & 95.9 & 28.2 \\
$\text{Vax}_{ik}$ (\% of pop.)  & 33.038 & 4.212 & 32.768 & 99.2 & 12.7 \\
\bottomrule
\end{tabular}
\vspace{4pt}

\parbox{\textwidth}{\scriptsize\textit{Notes:} 38 OECD countries,
Q1\,2020--Q2\,2022 ($N = 380$). Overall SD: total standard deviation.
Between SD: standard deviation of country means. Within SD: standard
deviation of deviations from country means. Within (\%) = Within SD /
Overall SD $\times$ 100. Country fixed effects absorb between-variation;
within-variation identifies the coefficients. All policy variables
exhibit within-shares above 93\%, confirming that the identifying
variation survives absorption of country fixed effects.}
\end{table}


\begin{table}[H]
\centering
\caption{Pairwise Correlation Matrix}
\label{tab:corr_matrix}
\scriptsize
\begin{tabular}{@{} l r r r r r r r @{}}
\toprule
 & $y_{ik}$ & $S_{ik}$ & $F^{CP}_{ik}$ & $F^{DI}_{ik}$
 & $d_{ik}$ & $\theta_{ik}$ & $\text{Vax}_{ik}$ \\
\midrule
$y_{ik}$           & 1    &       &       &       &       &       &  \\
$S_{ik}$           & $-$0.43 & 1  &       &       &       &       &  \\
$F^{CP}_{ik}$      & $-$0.25 & 0.13  & 1  &       &       &       &  \\
$F^{DI}_{ik}$      & $-$0.11 & 0.12  & 0.25  & 1  &       &       &  \\
$d_{ik}$           & $-$0.09 & 0.20  & $-$0.15 & $-$0.04 & 1  &       &  \\
$\theta_{ik}$      & 0.20  & $-$0.39 & $-$0.19 & $-$0.17 & 0.32  & 1  &  \\
$\text{Vax}_{ik}$  & 0.43  & $-$0.56 & $-$0.36 & $-$0.27 & 0.04  & 0.53  & 1 \\
\bottomrule
\end{tabular}
\vspace{4pt}

\parbox{\textwidth}{\scriptsize\textit{Notes:} Pearson correlations,
estimation sample (38 countries, Q1\,2020--Q2\,2022, $N = 380$).
Lower triangle reported. The correlations
$r(S, F^{CP}) = 0.13$ and $r(S, F^{DI}) = 0.12$ confirm that fiscal
composition is only weakly correlated with containment intensity---the
key identifying condition for the $S \times F^{CP}$ interaction term.}
\end{table}


\section{Health Equation}
\label{app:health_estimations}

Discussion about literature and estimates

\section{Robustness}
\label{app:robustness}


alle robustness checks-> zeigen wo, robust (alles relevante) gegenüber...

Nicht kausal deshalb so...-> sagen was sonst gemacht werden könnte und warum nicht, was kann ich mehr
Müsste als System geschätzt werden...
coefficie


%Vorschlag Struktur


%%DRAFT
% =====================================================================
%  APPENDIX: ROBUSTNESS ANALYSIS  (output-gap equation, V14 main spec)
%  Requires: booktabs. Place after \appendix.
%  All figures verified against the full re-run on 2026-06-15.
% =====================================================================
\subsection{Robustness Analysis Output}
\label{app:robustness_output}
This appendix documents the full battery of checks behind the preferred
output-gap specification (V14, eq.~\ref{eq:output_est}): the country
fixed-effects estimator
\[
y_{ik}=\rho_y\,y_{i,k-1}+\alpha_S\,S_{ik}
       +\alpha_{\text{above}}\,F^{CP,\text{above}}_{i,k-2}
       +\alpha_{\text{below}}\,K^{CP,\text{below}}_{ik}
       +\alpha_{DI}\,F^{DI}_{i,k-1}
       +\alpha_{S\cdot DI}\,S_{ik}F^{DI}_{i,k-1}
       +\gamma_i+\varepsilon_{ik},
\]
estimated on 38 OECD countries over $t\in[4,14]$ (Q4.2019--Q2.2022,
$N=418$), with standard errors clustered by country (CRV1, small-sample
corrected). The exercises span six concerns: (i)~identification and
collinearity, (ii)~estimator and dynamic-panel bias, (iii)~heterogeneous-
treatment / negative-weight contamination, (iv)~time effects and sample
composition, (v)~instrument construction and functional form, and
(vi)~small-cluster inference. Table~\ref{tab:robustness_summary} lists each
check, what it probes, and its outcome.

The two headline results are robust throughout: the above-the-line
capacity-preservation multiplier ($\alpha_{\text{above}}\approx0.54$, flow,
two-quarter lag) and the demand-injection ``push-on-string'' slope
($\alpha_{S\cdot DI}\approx-0.041$, break-even at $S^\star\approx36$) retain
their sign, magnitude and significance across every specification, sample
window and inference method. The below-the-line stock channel
($\alpha_{\text{below}}\approx0.26$) is the only fragile coefficient: it is
sensitive to the inclusion of time trends and to the pre-pandemic
zero-stock anchor, and is therefore reported as a theoretically motivated
lower bound rather than a precisely identified estimate.

\begin{table}[H]
\centering
\scriptsize
\renewcommand{\arraystretch}{1.15}
\caption{Robustness battery for the V14 output-gap specification}
\label{tab:robustness_summary}
\begin{tabular}{@{}p{0.255\textwidth} p{0.37\textwidth} p{0.315\textwidth}@{}}
\toprule
\textbf{Check} & \textbf{What it probes} & \textbf{Outcome} \\
\midrule
\multicolumn{3}{@{}l}{\textit{A.\ Identification \& collinearity}}\\
$S$--composition orthogonality & Fiscal mix endogenous to severity? & $r(\bar S,\bar F^{CP})\!\approx\!0.13$; composition driven by institutions/safety nets, not contemporaneous severity. \\
\addlinespace
\multicolumn{3}{@{}l}{\textit{B.\ Estimator \& dynamic-panel bias}}\\
\texttt{feols} vs.\ \texttt{plm} & Estimator invariance & Identical point estimates. \\
Nickell-bias assessment & Short-$T$ dynamic-panel bias & Bias $<0.7$pp, negligible vs.\ SEs; corroborated by Arellano--Bond GMM. \\
Mundlak / Chamberlain CRE & Fixed vs.\ random effects & Within-effects $=$ V14; joint test $\chi^2=25.7$, $p<0.001\Rightarrow$ FE required. \\
\addlinespace
\multicolumn{3}{@{}l}{\textit{C.\ Heterogeneous-treatment robustness}}\\
de--Chaisemartin--D'Haultf\oe uille & Negative TWFE weights & Negative-weight shares: Above $8.7\%$, DI $2.7\%$ (benign); Below $26.8\%$ (flags its fragility). \\
\addlinespace
\multicolumn{3}{@{}l}{\textit{D.\ Time effects \& sample composition}}\\
Year-FE / linear trend / Quarter-FE & Common time variation & Above-Flow \& DI$\times S$ trend-robust; Below-Stock trend-confounded (sign flips under Year-FE absorbing the same variation). \\
Start/end restrictions ($t\!\ge\!5,\ge\!6,\le\!13$) & Pre-pandemic anchor & Above/DI$\times S$ stable; Below inflates to $5.87$ at $t\!\ge\!6$ --- the Q4.2019 zero-stock anchor is essential. \\
Estimation--window (Narrow/Wide/Only-2020) & Window choice & Above $0.37$--$0.73$, DI$\times S$ $-0.04$ to $-0.05$ stable; Below significant only in baseline. \\
Outlier exclusion (TUR, IRL) & Inflation- / MNC-driven GDP & Core unchanged: Above $0.45^{*}$, DI$\times S$ $-0.035^{*}$. \\
\addlinespace
\multicolumn{3}{@{}l}{\textit{E.\ Instrument construction \& functional form}}\\
Take-up grid (loans 40/60/80\%, guar.\ 25/35/50\%) & Below-stock take-up assumptions & Above ($0.54$) \& DI$\times S$ ($-0.04$) invariant; Below $0.13$--$0.26$ (baseline largest). \\
Lag selection (Above 0--3, DI 0--2) & Disbursement timing & Only Above lag-2 and DI lag-1 yield sign-correct, theory-consistent estimates. \\
Alternative containment (peak $S$, transit mobility) & Measurement of $S$ & Above and push-on-string survive both. \\
\addlinespace
\multicolumn{3}{@{}l}{\textit{F.\ Mechanism \& heterogeneity}}\\
Tightening vs.\ loosening asymmetry & Hysteresis / ratchet & $\Delta S^{+}=-0.125^{***}$ vs.\ $\Delta S^{-}=-0.005$ ($\approx26\times$): structural support for AR(1) persistence. \\
Sample splits ($S$, income, pre-debt, safety net) & Where channels operate & Above strongest in high-$S$/high-debt; DI push-on-string in all splits; Below only in high-$S$ \& weak-safety-net groups what makes sense. \\
\addlinespace
\multicolumn{3}{@{}l}{\textit{G.\ Small-cluster inference} ($N_{\text{cl}}=38$)}\\
Wild-cluster restricted bootstrap-$t$ (CGM 2008, Rademacher) & CRV1 over-rejection & Above-Flow $p=0.008$; Below \& DI$\times S$ marginal ($p\approx0.06$); qualitative pattern unchanged. \\
Cluster-jackknife (leave-one-country-out) & Influential clusters & Above-Flow stable; no single country drives the result. \\
\bottomrule
\end{tabular}
\vspace{4pt}
\parbox{\textwidth}{\scriptsize\textit{Notes:} Dependent variable: output gap $y_{ik}$ (pp of potential GDP). All
fiscal variables in pp of 2019 GDP; $S$ on a $0$--$100$ scale. All models include country fixed effects with
country-clustered standard errors. Significance: $^{*}p<0.1$, $^{**}p<0.05$, $^{***}p<0.01$. Above $=$
$F^{CP,\text{above}}_{i,k-2}$ (above-the-line flow); Below $=$ $K^{CP,\text{below}}_{ik}$ (below-the-line
accumulated stock); DI$\times S=$ the demand-injection $\times$ stringency interaction.}
\end{table}

Taken together, the robustness checks show that the paper's two main transition-system claims do not depend on any single modeling choice, sample window, instrument-construction assumption, or inference procedure: capacity preservation has a positive lagged effect on output, and the effect of demand injection is state-dependent, becoming ineffective under tight containment. The below-the-line stock channel is the weakest element of the specification, since the cumulative liquidity stock is mechanically trending and therefore harder to separate from common pandemic time dynamics once richer time controls are introduced.

% =====================================================================
%  APPENDIX: ROBUSTNESS ANALYSIS  (debt equation, V15 main spec)
%  Requires: booktabs. Place after \appendix.
%  All figures verified against the full re-run on 2026-06-15 (N = 494).
% =====================================================================
\subsection{Robustness Analysis: Debt Equation}
\label{app:robustness_debt}

This appendix documents the robustness battery for the preferred debt-transition
specification (V15, eq.~\ref{eq:debt_est}). The dependent variable is the
\emph{interest-adjusted} change in the real debt-to-GDP ratio,
$\Delta \tilde b_{ik}=\Delta b_{ik}-r_{k}\,b_{i,k-1}$, which nets out the
mechanical debt-service component. The estimator is country fixed effects,
\[
\Delta \tilde b_{ik}= \gamma_y\,y_{i,k-1}
   +\kappa_{\text{above}}\,F^{CP,\text{above}}_{ik}
   +\kappa_{\text{loans}}\,L_{ik}
   +\kappa_{\text{guar}}\,G_{ik}
   +\kappa_{DI}\,F^{DI}_{ik}
   +\kappa_{H}\,F^{H}_{i,k-1}
   +\delta_i+\varepsilon_{ik},
\]
estimated on 38 OECD countries over $t\in[4,16]$ (Q4.2019--Q4.2022, $N=494$),
with standard errors clustered by country. Below-the-line commitments enter at
the central take-up calibration (loans 60\%, guarantees 25\%). The preferred
estimates are
$\hat\gamma_y=-0.13^{**}$, $\hat\kappa_{\text{above}}=0.62^{**}$,
$\hat\kappa_{\text{loans}}=0.81^{***}$, $\hat\kappa_{\text{guar}}=0.51$,
$\hat\kappa_{DI}=0.33$, and $\hat\kappa_{H}=1.08^{*}$
(within $R^2=0.23$). Table~\ref{tab:robustness_debt} lists every check.

The \emph{robust core} of the debt equation comprises the loan channel
($\kappa_{\text{loans}}$), which is significant in every specification and
implies an effective pass-through of $\approx 0.49$ (loans convert roughly
one-for-one into realised debt at $\approx 40$--$60\%$ drawdown); the lagged
health channel ($\kappa_{H}$); and the automatic-stabiliser term
($\gamma_y$), which is robust in real terms. The above-the-line channel
($\kappa_{\text{above}}$) is robust to most checks but weakens once a time
trend or year/quarter fixed effects absorb the common debt drift. Guarantees
($\kappa_{\text{guar}}$) and demand injection ($\kappa_{DI}$) are the weak
channels---consistent with low contingent-liability conversion of guarantees
and with DI being a small, transfer-based flow.

\begin{table}[H]
\centering
\scriptsize
\renewcommand{\arraystretch}{1.15}
\caption{Robustness battery for the V15 debt-transition specification}
\label{tab:robustness_debt}
\begin{tabular}{@{}p{0.255\textwidth} p{0.37\textwidth} p{0.315\textwidth}@{}}
\toprule
\textbf{Check} & \textbf{What it probes} & \textbf{Outcome} \\
\midrule
\multicolumn{3}{@{}l}{\textit{A.\ Outcome \& output-state definition}}\\
Unadjusted debt change (drop debt-service netting) & DV $=\Delta b$ instead of $\Delta\tilde b$ & Near-identical: $\kappa_{\text{loans}}=0.79^{***}$, $\kappa_{\text{above}}=0.57^{*}$, $\kappa_{H}=1.03^{*}$; netting is not driving results. \\
Output-gap timing & $y_{i,k-1}$ (main) vs.\ $y_{ik}$ vs.\ $y_{i,k-2}$ & Fiscal channels stable; $y_{ik}$ absorbs more ($\gamma=-0.20^{***}$), $y_{i,k-2}$ null ($-0.01$). Predetermined $y_{i,k-1}$ is the conservative choice. \\
Nominal outcome & DV $=$ nominal debt change & $\kappa_{\text{loans}}=0.91^{***}$ stable \\
\addlinespace
\multicolumn{3}{@{}l}{\textit{C.\ Instrument construction \& timing}}\\
Take-up grid (loans 40/60/80\%, guar.\ 25/35/50\%) & Below-the-line take-up assumptions & Effective pass-through invariant ($0.40\times1.22=0.60\times0.81\approx0.49$ for loans); coefficient scales inversely. \\
Lag below-the-line 1Q & Loans/guarantees lagged & Lagged terms insignificant ($0.55$, $0.12$) $\Rightarrow$ below-the-line affects debt \emph{contemporaneously} (booked at announcement). \\
Health-spending timing & $F^H$ vs.\ $F^H_{k-1}$ vs.\ both & Contemporaneous $F^H$ null ($-0.08$); $F^H_{k-1}=1.08^{*}$ $\Rightarrow$ lag-1 (delayed budgetary recording) is correct. \\
\addlinespace
\multicolumn{3}{@{}l}{\textit{D.\ Time effects \& sample composition}}\\
Estimation window ($t\!\ge\!5$; $\le\!14$; $\le\!15$) & Inclusion of late pandemic quarters & All coefficients stable ($\kappa_{\text{loans}}$ $0.79$--$0.82^{***}$, $\kappa_{H}$ $0.99$--$1.08^{*}$). \\
Linear trend / Year-FE / Quarter-FE & Common debt drift \& deployment cycle & $\kappa_{\text{loans}}$ \& $\kappa_{H}$ survive throughout; $\kappa_{\text{above}}$ weakens under Year-FE ($0.39$, n.s.) and Quarter-FE ($0.26$, n.s.). Quarter-FE over-absorbs the deployment cycle. \\
Outlier exclusion (TUR, IRL) & Inflation- / MNC-driven debt \& GDP & Core unchanged. \\
\addlinespace
\multicolumn{3}{@{}l}{\textit{E.\ Small-cluster inference} ($N_{\text{cl}}=38$)}\\
Wild-cluster restricted bootstrap-$t$ (CGM 2008, Rademacher) & CRV1 over-rejection & $\gamma_y$ ($p=0.005$), $\kappa_{\text{above}}$ ($0.014$), $\kappa_{\text{loans}}$ ($0.036$), $\kappa_{H}$ ($0.014$) remain significant; $\kappa_{\text{guar}}$ ($0.13$) and $\kappa_{DI}$ ($0.38$) not. \\
Leave-one-country-out jackknife & Influential clusters & $\kappa_{\text{loans}}\in[0.74,1.03]$, $\kappa_{\text{above}}\in[0.52,0.77]$, $\gamma_y\in[-0.14,-0.11]$; no single country overturns the result (most influential: JPN). \\
\bottomrule
\end{tabular}
\vspace{4pt}
\parbox{\textwidth}{\scriptsize\textit{Notes:} Dependent variable: interest-adjusted real debt change $\Delta\tilde b_{ik}$
(pp of 2019 GDP). All fiscal variables in pp of 2019 GDP. All models include country fixed effects with
country-clustered standard errors; small-sample correction applied. Significance: $^{*}p<0.1$, $^{**}p<0.05$,
$^{***}p<0.01$. Loans $=L_{ik}$ (60\% take-up), guarantees $=G_{ik}$ (25\% take-up), $F^H_{k-1}$ lagged one
quarter. Effective pass-through $=$ coefficient $\times$ assumed take-up rate.}
\end{table}

Overall, the debt equation's central message---that realised debt accumulation
is driven mechanically by below-the-line \emph{loans} (near one-for-one) and by
\emph{health} outlays recorded with a one-quarter delay, partly offset by
automatic stabilisers, while guarantees and demand injection contribute less---it is invariant to the outcome definition, the timing of the output control, instrument-construction assumptions, sample window, and small-cluster inference. The above-the-line channel is robust except under the most aggressive time-fixed-effects controls, which absorb the common debt drift that also identifies it.


% =============================================================================
% =====================================================================
%  Appendix subsection: Wave-specific effective infection fatality rates
%  Drop-in for the calibration appendix. Requires booktabs, siunitx (optional),
%  and a bibliography providing the \citep keys listed in the NOTE below.
%
%  CITATION KEYS USED (map these to your .bib):
%    levin2020, meyerowitzkatz2020, brazeau2020, covidforecast2022,
%    marziano2023, ward2024, erikstrup2022, recovery2021,
%    davies2021nature, davies2021severity, challen2021, nyberg2022,
%    veneti2022, sievers2022, wang2023
%  Where a single empirical claim is supported by more than one study,
%  multiple keys are grouped in one \citep{...,...}.
% =====================================================================

\subsection{Wave-specific effective infection fatality rates}
\label{app:ifr}

The structural inversion that constructs the prevalence proxy,
$\hat{\theta}_w = d_{w+\ell}/\delta_\theta$, requires a wave-specific
infection fatality rate $\delta_\theta$, the probability of death conditional
on infection. Because a lower $\delta_\theta$ implies more infections for the
same observed deaths, the level of $\delta_\theta$ pins the scale of
$\hat{\theta}$ within each wave, while its movement across waves governs the
shape of the imputed prevalence path. We therefore calibrate $\delta_\theta$
to the empirical IFR literature rather than treating it as a free parameter.

\paragraph{Effective, not intrinsic, fatality.}
The object we require is the \emph{effective} population IFR realized in each
wave, not the intrinsic per-infection lethality of the dominant variant in an
immune-naive host. The deaths that the inversion divides by $\delta_\theta$
are the deaths that actually occurred, in a population whose vaccination
coverage, accumulated natural immunity, and standard of clinical care all
evolved over the sample. The relevant divisor is the fatality rate that this
real, time-varying population experienced. This distinction resolves what
would otherwise appear to be a puzzle in the calibrated sequence: the
effective IFR falls monotonically from the ancestral strain through Omicron,
\emph{even though} the intrinsic virulence of the Alpha and Delta variants
exceeded that of the wild type. The fall is driven by immunity and treatment
improving faster than intrinsic severity rose; once ``effective'' is read as
deaths-realized-in-the-actual-population, the two facts are consistent.

\paragraph{Evidence on intrinsic severity.}
Intrinsic severity rose across the pre-Omicron variants. Relative to the wild
type, the Alpha variant carried an adjusted all-cause mortality hazard ratio
of roughly $1.6$--$1.7$ in large community cohorts
\citep{challen2021,nyberg2022}, and was independently estimated to be on the
order of $30$--$50\%$ more severe \citep{davies2021severity}. Delta was more
severe still than Alpha on both hospitalization and death
\citep{marziano2023}. Omicron then reversed the trend sharply, with adjusted
hospitalization and death hazards roughly $65$--$73\%$ below Delta
\citep{veneti2022,sievers2022,wang2023}. The same vaccination-stratified
studies attribute the Omicron decline jointly to reduced intrinsic
pathogenicity and to vaccination, with vaccine efficacy against death
conditional on infection on the order of $90\%$ \citep{wang2023}. It is
precisely this protection effect, rising steeply across the Alpha and Delta
quarters, that pushes the effective IFR down even as intrinsic virulence
climbs.

\paragraph{Calibrated values.}
Table~\ref{tab:ifr_by_wave} reports the calibrated $\delta_\theta$ by wave
together with the literature anchors. For the ancestral strain in spring 2020,
the OECD-restricted age-IFR meta-analysis of \citet{levin2020} implies a
population IFR near $0.9\%$ for a representative high-income age structure,
with supporting estimates of $0.68\%$ \citep{meyerowitzkatz2020} and $1.15\%$
\citep{brazeau2020}; we adopt a value at the lower end of this band,
consistent with an OECD-wide rather than single-country average. The summer
2020 lull combines the availability of dexamethasone from June 2020
\citep{recovery2021} with a younger infected age profile, both lowering the
effective IFR. The autumn--winter wild-type wave is set modestly higher, as
the surge re-loaded older and care-home populations and re-stressed hospital
capacity; this is an age-composition-and-capacity effect rather than a change
in the virus, and is consistent with the ancestral-phase effective IFR in
\citet{marziano2023} and the pre-vaccine estimates of
\citet{covidforecast2022}. The Alpha and Delta waves carry higher intrinsic
severity \citep{davies2021nature,challen2021,nyberg2022} offset by rising
vaccination coverage, yielding net effective IFRs of roughly $0.4\%$ and
$0.3\%$; the by-variant effective series of \citet{marziano2023} and the
real-time England estimates of \citet{ward2024} place these phases at
comparable levels and declining. For Omicron, the effective IFR collapses:
\citet{erikstrup2022} estimate $0.006\%$ in an age-restricted Danish
serosurveillance sample, \citet{marziano2023} report approximately $0.05\%$
in Italy, and \citet{ward2024} reach approximately $0.01\%$ by spring 2022.

The central IFR for the late-pandemic wave is set at $0.04\%$, below the
REACT-1 England estimate of $0.097\%$ for early Omicron (BA.1), reflecting the
progressively milder BA.2 and BA.5 subvariants that dominated the latter half
of 2022 and the substantially higher hybrid-immunity levels in later quarters.

\paragraph{Role in the estimation.}
Because $\delta_\theta$ enters the inversion as a divisor, the imputed
$\hat{\theta}$ panel is monotone in $\delta_\theta$ within each wave. The
second-stage fiscal estimates, however, depend on $\hat{\theta}$ only through
its cross-sectional and cross-wave \emph{ranking}, consistent with the role of
$\hat{\theta}$ as a reduced-form pressure variable rather than a biological
estimate. The effective-IFR levels in Table~\ref{tab:ifr_by_wave} fix the
scale of $\hat{\theta}$; the wave ordering that the empirical strategy
requires is robust to moving any single wave within its literature band. We
report sensitivity over the IFR bands alongside the $\rho_\theta$ sensitivity
in the calibration report.

\begin{table}[htbp]
\centering
\caption{Wave-specific effective infection fatality rates}
\label{tab:ifr_by_wave}
\begin{tabular}{llcl}
\toprule
Wave & Period & $\delta_\theta$ (\%) & Principal anchors \\
\midrule
W1 (ancestral)      & Q4.2019--Q1.2020 & 0.50 & \citet{levin2020,meyerowitzkatz2020,brazeau2020} \\
W1 summer lull      & Q2.2020          & 0.20 & \citet{recovery2021}; younger age profile \\
W2 wild type        & Q3.2020          & 0.70 & \citet{covidforecast2022,marziano2023} \\
W2 Alpha            & Q4.2020--Q1.2021 & 0.40 & \citet{davies2021nature,challen2021,nyberg2022} \\
W3 Delta            & Q2.2021          & 0.30 & \citet{marziano2023,ward2024} \\
W4 Omicron          & Q3.2021--Q4.2021 & 0.20 & \citet{marziano2023,ward2024,erikstrup2022} \\
Late Omicron (BA.2/BA.5) & Q1.2022--Q4.2022 & 0.02 & \citet{ward2024,erikstrup2022} \\
\bottomrule
\end{tabular}
\begin{minipage}{0.92\linewidth}
\vspace{4pt}
\footnotesize
\textit{Notes:} $\delta_\theta$ is the effective population infection fatality
rate used in the inversion $\hat{\theta}_w = d_{w+\ell}/\delta_\theta$, and
incorporates intrinsic variant virulence, treatment improvements, vaccination
coverage, and prior immunity at the time of each wave. Effective rates fall
monotonically across waves despite higher intrinsic severity of Alpha and
Delta relative to the wild type, because protection from vaccination and
prior infection rose faster than intrinsic severity. Values are calibrated to
the cited literature; the second-stage estimates depend on $\hat{\theta}$ only
through its cross-wave and cross-country ranking.
\end{minipage}
\end{table}

% =====================================================================
%  Appendix subsection: Transmission parameter across resolutions
%  Drop-in for the calibration appendix, to follow the IFR subsection.
%  Requires the same bibliography. CITATION KEYS used:
%    hart2022, manica2022, madewell2023, brauner2021, flaxman2020
%  (Map these to your .bib; they correspond to refs [21]-[25] in the R code.)
% =====================================================================

\subsection{The transmission parameter across resolutions}
\label{app:rho}

The weekly imputation in Appendix~\ref{app:ifr} and the quarterly structural
model use the same infection-transition equation,
$\theta_{t+1} = \rho_\theta\,(1 - \phi_S S_t)\,\theta_t$, but assign
$\rho_\theta$ at different resolutions. The weekly construction employs a
fully wave-specific $\rho_\theta^w$ that rises across variants in line with
the transmissibility ordering---ancestral, Alpha, Delta, Omicron---reflecting
shorter generation times and higher basic reproduction numbers for later
variants \citep{hart2022,manica2022,madewell2023}. The quarterly model instead
uses a parsimonious two-regime $\rho_\theta$, constant before the vaccination
rollout and halved thereafter. This subsection explains why the two are
consistent, so that the coarser model specification is a deliberate
simplification rather than a discrepancy.

\paragraph{The imputed prevalence does not depend on $\rho_\theta$.}
The prevalence proxy is constructed by inverting observed mortality,
$\hat{\theta}_w = d_{w+\ell}/\delta_\theta$, and therefore does not use
$\rho_\theta$ at any point. The transmission parameter enters only the
\emph{transition} equation that links successive prevalence values, not the
\emph{level} of $\hat{\theta}$ itself. Consequently, the weekly $\rho_\theta^w$
is a parameter for assessing the plausibility of the transition equation
against the imputed panel, not an input to the data that the model is
calibrated to. The choice of $\rho_\theta$ resolution cannot alter the
$\hat{\theta}$ series.

\paragraph{The fitted prevalence path is invariant to $\rho_\theta$.}
Within the structural model, the transition is closed by wave-specific
innovations $\varepsilon_\theta$ defined as the residual between observed and
expected prevalence. Because these innovations are constructed from, and added
back to, the same deterministic term $\rho_\theta\,(1-\phi_S S)\,\theta$, the
parameter cancels in the fitted path: any value of $\rho_\theta$ leaves the
reproduced $\hat{\theta}$ trajectory unchanged, shifting only the share of
each period's movement attributed to deterministic persistence versus
innovation. A forward roll with the innovations switched off decays to
negligible prevalence while the observed series rises by more than two orders
of magnitude, confirming that the innovations, not $\rho_\theta$, carry the
epidemic dynamics. The two-regime quarterly specification therefore yields the
same fitted path as a fully wave-specific one; the additional weekly
granularity would add free parameters without changing any reproduced moment.

\paragraph{Why the model retains a vaccination break.}
The single feature of the weekly $\rho_\theta^w$ profile that is economically
material at quarterly resolution is the post-vaccination decline in effective
transmission, which is what disciplines the planner's containment problem: it
governs how much stringency is required to stabilize prevalence once
vaccination weakens the chain of transmission. The two-regime specification
preserves exactly this feature, calibrated to the suppression effectiveness
$\phi_S$ estimated from non-pharmaceutical-intervention studies
\citep{brauner2021,flaxman2020}, while discarding the within-regime variation
that the invariance result renders immaterial. The model is thus consistent
with the weekly construction on the dimension that matters for the optimal
policy and simplified only where simplification has no effect on the validated
trajectory.




\section{Calibration Details and Additional Validation}
\label{app:calibration}

This appendix reports additional details on the calibration exercise described in Section~\ref{subsec:calibration}. The purpose of the calibration is not to estimate the transition system again, but to assess whether the estimated and calibrated model can reproduce the main pandemic dynamics under observed policy paths. The validation is conducted on the full panel of 38 OECD economies over Q4 2019--Q4 2022. Output is evaluated over Q4 2019--Q2 2022, corresponding to the main pandemic macroeconomic window used in the output transition, while debt is evaluated through Q4 2022 to capture the fiscal legacy of the pandemic response.

\subsection{Health Calibration}

Country-level infection states are matched using calibrated innovations, reflecting the fact that observed infection dynamics were strongly affected by variant arrivals, testing regimes, vaccination, and other epidemiological shocks that are not fully captured by the parsimonious transition system. The mortality mapping is wave-specific and converts the infection state into excess deaths using effective population IFRs. These effective IFRs incorporate intrinsic variant severity, treatment improvements, vaccination coverage, and prior immunity. The pre-vaccine IFR calibration is anchored in age-specific and population-level IFR estimates from early meta-analyses and global evidence \citep{levin2020assessing,meyerowitzkatz2020systematic,brazeau2020infection,covidforecasting2022variation}. Treatment-related declines are motivated by evidence on dexamethasone \citep{recovery2021dexamethasone}, while variant-specific severity and the decline in fatality risk over Delta and Omicron phases are based on evidence from Alpha severity studies and later variant-period IFR estimates \citep{davies2021increased,marziano2023estimating,ward2024realtime,erikstrup2022seroprevalence}. Vaccination and prior-immunity effects enter through the effective IFR interpretation and through the post-vaccination reduction in infection persistence \citep{degier2021vaccine}.

\begin{table}[H]
\centering
\captionsetup{font=scriptsize}
\caption{Wave-specific effective population IFRs used to map the infection state into excess deaths. Effective IFRs incorporate variant severity, treatment improvements, vaccination coverage, and prior immunity.}
\label{tab:ifr_waves}
\scriptsize
\begin{tabular}{llcc}
\toprule
\textbf{Wave} & \textbf{Interpretation} & \textbf{Effective IFR} & \textbf{Quarters} \\
\midrule
W1 & First wave, pre-treatment, no vaccination & $0.50\%$ & Q4.19--Q1.20 \\
W1-summer & Lower overload and treatment improvement & $0.20\%$ & Q2.20 \\
W2-wildtype & Autumn/winter 2020--21 wildtype phase & $0.70\%$ & Q3.20 \\
Alpha & Alpha phase, partly offset by early vaccination & $0.40\%$ & Q4.20--Q1.21 \\
Delta & Delta phase with increasing vaccination coverage & $0.20\%$ & Q2.21 \\
Early Omicron & Omicron transition phase & $0.20\%$ & Q3.21--Q1.22 \\
Late Omicron & High immunity, lower fatality phase & $0.02\%$ & Q2.22--Q4.22 \\
\bottomrule
\end{tabular}
\end{table}


\subsection{Country-specific Inputs}

\begin{table}[H]
\centering
\caption{Country-specific inputs to the calibration.}
\label{tab:country_inputs}
\scriptsize
\begin{tabular}{lcccc}
\toprule
\textbf{Country} & $B_{i,2019\text{Q}4}$ & $\mu_{y,i}$ & $\mu_{b,i}$ & $\varepsilon^{y}_{i,2020\text{Q}2}$ \\
 & (\% GDP) & (pp) & (pp) & (pp) \\
\midrule
AUS & 37.6 & +1.106 & -0.853 & -3.62 \\
AUT & 69.7 & -1.040 & -0.600 & -8.55 \\
BEL & 77.9 & +0.301 & -0.873 & -6.46 \\
CAN & 44.2 & -0.098 & -0.265 & -7.79 \\
CHE & 16.3 & +1.299 & -0.431 & -4.73 \\
CHL & 29.6 & +1.325 & -1.029 & -10.10 \\
COL & 47.7 & +1.989 & +0.549 & -11.20 \\
CRI & 53.2 & +0.147 & +0.121 & -3.75 \\
CZE & 30.2 & -3.562 & -0.657 & -5.10 \\
DEU & 34.0 & -1.738 & -1.095 & -5.98 \\
DNK & 34.1 & -0.083 & -1.661 & -1.95 \\
ESP & 86.0 & -4.896 & -0.962 & -9.57 \\
EST & 11.8 & -1.833 & -0.481 & -2.43 \\
FIN & 54.5 & -1.483 & -0.631 & -1.97 \\
FRA & 85.4 & -1.869 & -0.967 & -8.49 \\
GBR & 107.0 & -3.327 & -1.372 & -12.50 \\
GRC & 205.0 & +0.291 & -1.304 & -9.94 \\
HUN & 52.1 & -2.015 & -0.967 & -8.99 \\
IRL & 54.3 & +8.319 & -0.270 & -1.10 \\
ISL & 69.9 & -4.849 & -0.817 & -4.70 \\
ISR & 56.9 & +2.367 & -0.342 & -0.66 \\
ITA & 122.0 & -0.501 & -1.193 & -10.90 \\
JPN & 199.0 & -1.956 & -1.657 & -6.82 \\
KOR & 35.1 & +0.614 & -0.514 & -1.83 \\
LTU & 29.3 & +0.700 & -0.482 & -0.03 \\
LUX & 26.6 & +2.357 & +0.119 & -5.31 \\
LVA & 39.7 & -0.683 & +0.048 & -7.43 \\
MEX & 7.2 & -3.206 & -0.908 & -7.97 \\
NLD & 43.9 & +1.060 & -0.427 & -5.17 \\
NOR & 16.6 & +1.222 & -0.109 & -2.45 \\
NZL & 73.0 & -1.062 & +0.039 & -5.49 \\
POL & 45.6 & -0.597 & -0.693 & -1.80 \\
PRT & 111.0 & -2.775 & -1.724 & -9.75 \\
SVK & 48.4 & -0.258 & -0.131 & -4.87 \\
SVN & 72.1 & -1.828 & -0.991 & -6.78 \\
SWE & 27.7 & +1.035 & -0.511 & -4.70 \\
TUR & 54.0 & +4.066 & -0.315 & -10.40 \\
USA & 97.3 & +1.099 & -0.225 & -4.79 \\
\bottomrule
\end{tabular}

\vspace{0.15cm}
\begin{minipage}{0.9\textwidth}
\scriptsize
\textit{Notes:} $B_{i,2019\text{Q}4}$ is the observed pre-pandemic public debt level (\% of 2019 GDP), used as the initial debt stock. $\mu_{y,i}$ and $\mu_{b,i}$ are the estimated country fixed effects in the output and debt equations. $\varepsilon^{y}_{i,2020\text{Q}2}$ is the country-specific initial pandemic output shock in 2020Q2. The fixed effects and initial shock are estimated on the panel; the debt baseline is taken from observed data.
\end{minipage}
\end{table}


\subsection{Representative OECD Economy}

The planner is solved for a representative OECD economy constructed from cross-country mean observed policy inputs and initial conditions. It is therefore important to validate the same object used in the optimization exercise. Under the observed mean policy path, the representative-economy rollout reproduces the aggregate pandemic dynamics reasonably well. The output RMSE is 1.24 percentage points, the quarterly debt-change RMSE is 0.85 percentage points, and the excess-mortality RMSE is 2.87 deaths per million per week what is even better as the median fit. This supports the use of the representative OECD economy as the baseline object for the counterfactual planner exercise.

\begin{table}[H]
\centering
\caption{Representative OECD economy validation under observed mean policy inputs.}
\label{tab:rep_economy_validation}
\scriptsize
\begin{tabular}{lcc}
\toprule
\textbf{Moment} & \textbf{Horizon} & \textbf{RMSE} \\
\midrule
Output gap & Q4.19--Q2.22 & 1.24 pp \\
Debt change & Q4.19--Q4.22 & 0.85 pp \\
Debt level & Q4.19--Q4.22 & 2.96 pp \\
Infection state $\theta$ & Q4.19--Q2.22 & 0.0024 \\
Excess mortality $d$ & Q4.19--Q2.22 & 2.87 \\
\bottomrule
\end{tabular}
\end{table}

\noindent Figure~\ref{fig:rep_oecd_validation} reports the validation for the representative OECD economy used in the planner exercise. The output-gap trajectory is closely aligned with the observed mean path, and the fit for excess mortality is even stronger than in the country-level panel validation. The debt trajectory is also well matched until around 2021Q2. Afterwards, however, the model starts to overestimate the debt level, reaching a gap of roughly 5 percentage points by the end of 2022. This difference is not large relative to the level of public debt, but it should be stated explicitly. Since the simulated output gap remains negative, and since the transition system does not include an additional consolidation, repayment, or stock-flow adjustment mechanism, there is no model-implied force that brings debt back down after the main fiscal expansion. This limitation does not undermine the efficiency comparison in the planner exercise, because the observed and counterfactual policy paths are evaluated within the same representative-economy transition system. The comparison therefore remains internally consistent, even if the model somewhat overstates the terminal debt level of the representative economy.

Figure~\ref{fig:calib_appendix_debtchange} shows the corresponding validation for quarterly debt changes. The model tracks the quarterly debt-change dynamics reasonably well, with the main deviation occurring in the first pandemic quarter. This is consistent with the construction of the fiscal variables: loans and guarantees enter the model as take-up-adjusted flows, while the observed debt data partly reflect the announcement and accounting recognition of large fiscal envelopes. The forward roll therefore underestimates the initial jump in debt, but captures the subsequent debt-change dynamics more closely.

\begin{figure}[H]
    \centering
    \includegraphics[width=\textwidth]{Files/text/main/mean_vali v20.pdf}
    \caption{Representative OECD economy validation under observed mean policy inputs. The figure compares the simulated representative-economy path to the observed OECD mean path.}
    \label{fig:rep_oecd_validation}
\end{figure}


\begin{figure}[H]
    \centering
    \includegraphics[width=0.7\textwidth]{Files/text/main/debt change V20.pdf}
    \caption{Quarterly debt-change validation under observed policy paths. Debt changes are the primary debt validation object and the weakest fitted component of the transition system.}
    \label{fig:calib_appendix_debtchange}
\end{figure}



\subsection{Calibration Fit per Country}

\begin{longtable}{lccccccccc}
\caption{Country-level endpoint fit: observed, simulated, and residuals.}
\label{tab:country_endpoint}\\
\toprule
 & \multicolumn{3}{c}{Output gap (pp)} & \multicolumn{3}{c}{Mortality ($d$)} & \multicolumn{3}{c}{Debt (\% GDP)} \\
\cmidrule(lr){2-4}\cmidrule(lr){5-7}\cmidrule(lr){8-10}
\textbf{Country} & Obs & Sim & Resid & Obs & Sim & Resid & Obs & Sim & Resid \\
\midrule
\endfirsthead
\multicolumn{10}{l}{\scriptsize\itshape Table~\ref{tab:country_endpoint} continued}\\
\toprule
 & \multicolumn{3}{c}{Output gap (pp)} & \multicolumn{3}{c}{Mortality ($d$)} & \multicolumn{3}{c}{Debt (\% GDP)} \\
\cmidrule(lr){2-4}\cmidrule(lr){5-7}\cmidrule(lr){8-10}
\textbf{Country} & Obs & Sim & Resid & Obs & Sim & Resid & Obs & Sim & Resid \\
\midrule
\endhead
\midrule
\multicolumn{10}{r}{\scriptsize continued on next page}\\
\endfoot
\bottomrule
\endlastfoot
\scriptsize
AUS & -0.02 & -0.94 & +0.92 & 11.45 & 9.41 & +2.04 & 41.28 & 40.39 & +0.89 \\
AUT & -1.81 & -4.77 & +2.96 & 10.45 & 7.10 & +3.35 & 80.58 & 79.71 & +0.87 \\
BEL & -0.06 & +0.89 & -0.95 & 7.99 & 5.07 & +2.92 & 81.56 & 80.19 & +1.37 \\
CAN & -1.58 & -2.08 & +0.50 & 8.92 & 5.76 & +3.16 & 58.08 & 57.62 & +0.46 \\
CHE & +1.07 & +1.32 & -0.26 & 2.82 & 4.06 & -1.24 & 18.44 & 18.58 & -0.13 \\
CHL & +0.48 & +0.10 & +0.39 & 8.10 & 19.46 & -11.36 & 32.13 & 32.81 & -0.68 \\
COL & +3.75 & +0.02 & +3.73 & 0.62 & 10.87 & -10.25 & 63.81 & 64.16 & -0.35 \\
CRI & -2.60 & +0.30 & -2.90 & 3.28 & 5.46 & -2.18 & 67.23 & 65.44 & +1.80 \\
CZE & -8.51 & -4.63 & -3.88 & 5.46 & 3.56 & +1.90 & 42.23 & 41.42 & +0.81 \\
DEU & -4.24 & -1.46 & -2.78 & 11.74 & 2.61 & +9.13 & 39.77 & 38.95 & +0.82 \\
DNK & -1.80 & +2.53 & -4.33 & 11.27 & 3.32 & +7.95 & 30.99 & 29.06 & +1.93 \\
ESP & -6.16 & -6.59 & +0.42 & 9.08 & 2.82 & +6.26 & 97.40 & 95.92 & +1.48 \\
EST & -6.63 & -3.43 & -3.19 & 8.63 & 13.41 & -4.77 & 21.99 & 20.95 & +1.04 \\
FIN & -4.57 & -2.44 & -2.13 & 25.70 & 12.36 & +13.34 & 60.66 & 58.77 & +1.89 \\
FRA & -3.52 & -2.04 & -1.49 & 8.50 & 4.76 & +3.75 & 92.23 & 90.72 & +1.52 \\
GBR & -3.20 & -3.87 & +0.67 & 10.53 & 0.42 & +10.11 & 118.13 & 116.70 & +1.43 \\
GRC & -0.33 & -2.46 & +2.13 & 20.89 & 21.13 & -0.24 & 212.81 & 208.10 & +4.71 \\
HUN & -4.93 & -2.37 & -2.57 & 10.86 & 4.78 & +6.09 & 59.00 & 58.24 & +0.76 \\
IRL & +10.90 & +10.73 & +0.17 & 11.65 & 3.34 & +8.30 & 52.38 & 50.78 & +1.60 \\
ISL & -5.89 & -6.30 & +0.42 & 13.32 & 21.23 & -7.91 & 80.58 & 78.70 & +1.88 \\
ISR & +1.65 & +2.06 & -0.41 & 3.79 & 11.21 & -7.41 & 65.61 & 63.44 & +2.17 \\
ITA & +1.27 & -1.15 & +2.43 & 12.03 & 9.42 & +2.61 & 128.11 & 126.00 & +2.11 \\
JPN & -2.30 & -2.25 & -0.05 & 1.85 & 8.11 & -6.26 & 206.86 & 201.18 & +5.68 \\
KOR & -2.24 & +0.34 & -2.58 & 13.86 & 21.01 & -7.16 & 40.57 & 38.92 & +1.65 \\
LTU & -4.25 & +0.60 & -4.85 & 8.98 & 22.89 & -13.92 & 34.30 & 32.84 & +1.45 \\
LUX & -2.03 & +1.06 & -3.09 & 7.71 & 0.83 & +6.87 & 34.12 & 33.89 & +0.24 \\
LVA & -2.78 & -1.18 & -1.60 & 11.58 & 17.16 & -5.58 & 53.30 & 52.56 & +0.73 \\
MEX & -4.98 & -6.16 & +1.18 & 1.74 & 15.06 & -13.32 & 8.19 & 9.05 & -0.85 \\
NLD & +0.73 & +0.18 & +0.55 & 2.02 & 4.80 & -2.78 & 49.45 & 48.75 & +0.70 \\
NOR & -0.00 & +1.68 & -1.68 & 13.48 & 7.61 & +5.87 & 20.57 & 20.48 & +0.09 \\
NZL & -3.78 & -3.05 & -0.73 & 11.55 & 5.29 & +6.26 & 94.19 & 92.32 & +1.86 \\
POL & -4.05 & -1.70 & -2.35 & 3.29 & 7.03 & -3.74 & 51.71 & 50.13 & +1.58 \\
PRT & -4.95 & -3.70 & -1.25 & 18.30 & 3.59 & +14.70 & 108.34 & 106.38 & +1.96 \\
SVK & -5.14 & -0.36 & -4.78 & 12.29 & 11.54 & +0.75 & 59.60 & 57.97 & +1.63 \\
SVN & -5.31 & -2.88 & -2.42 & 6.34 & 5.79 & +0.56 & 82.17 & 80.26 & +1.91 \\
SWE & -1.28 & +1.45 & -2.73 & 6.73 & 2.92 & +3.81 & 29.55 & 28.89 & +0.66 \\
TUR & +5.55 & +5.20 & +0.35 & 0.81 & 10.90 & -10.09 & 53.97 & 53.62 & +0.34 \\
USA & -1.33 & -0.34 & -0.99 & 8.43 & 11.58 & -3.15 & 116.57 & 113.89 & +2.67 \\
\midrule
\textit{Mean} & -1.97 & -0.99 & -0.98 & 9.11 & 8.89 & +0.22 & 68.12 & 66.78 & +1.33 \\
\textit{Median} & -2.27 & -1.16 & -0.97 & 8.95 & 7.06 & +1.32 & 58.54 & 57.80 & +1.44 \\
\textit{SD} & +3.64 & +3.24 & +2.11 & 5.41 & 6.22 & +7.33 & 44.88 & 43.85 & +1.24 \\
\end{longtable}



\subsection{Additional Fit Diagnostics}

The calibration also matches several non-targeted moments. The simulated output series has a mean cross-country SD ratio of 1.154, indicating that the model does not flatten the observed dispersion. Output persistence is close to the data, with an observed mean AC(1) of 0.244 and a simulated mean AC(1) of 0.155. The endpoint residuals are small: the median output-gap residual is $-0.97$ pp in Q2 2022, while the mean and median terminal debt residuals are $+1.33$ pp and $+1.44$ pp, respectively. The weakest component remains the quarter-to-quarter debt change, consistent with the main text.

\begin{table}[H]
\centering
\caption{Additional calibration diagnostics.}
\label{tab:additional_calib_diagnostics}
\scriptsize
\begin{tabular}{lc}
\toprule
\textbf{Diagnostic} & \textbf{Value} \\
\midrule
Output SD ratio & 1.154 \\
Output AC(1), observed & 0.244 \\
Output AC(1), simulated & 0.155 \\
Output AC(1) gap & 0.089 \\
Output ICC, observed & 0.356 \\
Output ICC, simulated & 0.420 \\
Debt ICC, observed & 1.018 \\
Debt ICC, simulated & 1.020 \\
Terminal debt residual, mean & +1.33 pp \\
Terminal debt residual, median & +1.44 pp \\
Output endpoint residual, mean & -0.98 pp \\
Output endpoint residual, median & -0.97 pp \\
\bottomrule
\end{tabular}
\end{table}

\subsection{Fiscal Channel Decomposition}

The fiscal channel decomposition confirms that the signs of the estimated policy channels are economically consistent. Over the main 11-quarter validation window, the median cumulative contribution of above-the-line support is $+3.45$ percentage points, while the median contribution of below-the-line liquidity support is $+6.35$ percentage points. Demand injection contributes negatively at the median, $-1.21$ percentage points, reflecting the estimated interaction with containment. Total fiscal support is positive in 33 of 38 countries and equals $+8.05$ percentage points at the median.

\begin{table}[H]
\centering
\caption{Fiscal channel decomposition over the main validation window.}
\label{tab:fiscal_channel_decomposition}
\scriptsize
\begin{tabular}{lccc}
\toprule
\textbf{Channel} & \textbf{Median contribution} & \textbf{Positive countries} & \textbf{Negative countries} \\
\midrule
Above-line fiscal support & +3.45 pp & 38 & 0 \\
Below-line liquidity support & +6.35 pp & 38 & 0 \\
Demand injection & -1.21 pp & 6 & 32 \\
Total fiscal support & +8.05 pp & 33 & 5 \\
Health drag & +0.00 pp & 0 & 0 \\
\bottomrule
\end{tabular}
\end{table}








\section{Solution Method and Augmented State Space}
\label{app:planner}

% =============================================================================

\subsection{Why the Policy Trade-off Is Unavoidable}
\label{app:necessity}

Mortality is predetermined within the quarter: from Eq.~(\ref{eq:trans_d}),
$d_{k+1} = \delta_\theta(k)\,\theta_k$ depends only on prevalence already
realized. The planner controls \emph{future} mortality by suppressing current
transmission. Combining the infection and mortality transitions,
%
\begin{equation}
    d_{k+2} = \delta_\theta(k+1)\,\rho_\theta(k)\bigl(1 - \phi_S\, S_k/100\bigr)\,\theta_k,
    \label{eq:excess_dynamics}
\end{equation}
%
so the marginal effect of containment on mortality two quarters out is
%
\begin{equation}
    \frac{\partial d_{k+2}}{\partial S_k}
    = -\,\delta_\theta(k+1)\,\rho_\theta(k)\,\frac{\phi_S}{100}\,\theta_k .
    \label{eq:dnecessity}
\end{equation}
%
The benefit of containment vanishes when the variant is non-lethal
($\delta_\theta = 0$): suppression then saves no lives and the planner sets
$S_k = 0$ to protect output. It is strictly positive only when the variant is
simultaneously lethal ($\delta_\theta > 0$) and transmissible
($\rho_\theta > 0$) during an active outbreak ($\theta_k > 0$). The planner
accepts the output cost of containment solely because of this conjunction. This
is the mechanism behind the post-vaccination unwinding of containment in
Section~\ref{sec:results}: once vaccination drives $\delta_\theta$ toward zero,
Eq.~(\ref{eq:dnecessity}) collapses and the optimal $S_k$ falls to zero even
while prevalence remains high.

\subsection{Augmented State Space}
\label{app:augmented}

The transition system in Section~\ref{sec:nonlinear} is not first-order Markov:
direct income support disburses with a one-quarter delay, entering both output
and the budget through $F^{\text{DI}}_{k-1}$, and the two capacity-preservation
channels operate through persistent stocks of accumulated fiscal support. iLQR
requires a first-order system, so I augment the four economic states
$(y_k, b_k, \theta_k, d_k)$ with three bookkeeping states that carry the lag and
the two stocks:
%
\begin{equation}
    \bm{x}_k = \bigl(y_k,\; b_k,\; \theta_k,\; d_k,\;
    z_k,\; L_k,\; K_k\bigr)' \in \mathbb{R}^7,
    \label{eq:state7}
\end{equation}
%
where $z_k = F^{\text{DI}}_{k-1}$, $L_k$ is the below-the-line liquidity stock,
and $K_k$ the above-the-line capacity-preservation stock. The control vector is
$\bm{u}_k = (S_k, F^{\text{above}}_k, F^{\text{loans}}_k, F^{\text{guar}}_k,
F^{\text{DI}}_k)' \in \mathbb{R}^5$. The bookkeeping recursions are
$z_{k+1} = F^{\text{DI}}_k$ and the stock accumulation equations of
Eq.~(\ref{eq:stocks}),
%
\begin{align}
    L_{k+1} &= (1-\delta_L)\,L_k + F^{\text{loans}}_k + F^{\text{guar}}_k,
    \label{eq:stockL_app}\\
    K_{k+1} &= (1-\delta_K)\,K_k + F^{\text{above}}_k.
    \label{eq:stockK_app}
\end{align}
%
Loan and guarantee flows are take-up-adjusted at the data-construction stage:
the realization rates $c_{\text{lo}} = 0.6$ and $c_{\text{gu}} = 0.25$ convert
authorized envelopes into the fraction that materializes as effective support,
so the flows entering Eq.~(\ref{eq:stockL_app}) and the debt transition are
already effective values and no take-up factor appears inside the dynamics.

Output loads on the \emph{end-of-period} stocks, which absorb the current flows:
the transition contains $\alpha_K\, K_{k+1}$ and the complementarity-augmented
liquidity term $\alpha_L\, L_{k+1}\, m(K_{k+1})$, with the bounded multiplier
$m(K) = 1 + \chi\, K/(K + \bar K)$ of Eq.~(\ref{eq:complementarity}). Above-the-line
support is therefore both a flow and a stock: the flow is absorbed
contemporaneously into $K_{k+1}$, which then governs output directly and raises
the effectiveness of outstanding liquidity. This capacity-stock channel replaces
the lag-2 flow specification used in estimation; the coefficient
$\alpha_K = 0.30\,\alpha_{\text{above}}$ reinterprets the estimated flow-lag
effect as a calibrated persistent margin, deliberately discounted to avoid
double-counting persistence already implicit in the lag structure
(Section~\ref{subsec:calibration}).

The fiscal timing embedded in the augmentation reflects the institutional
sequence. Above- and below-the-line support commit the budget on authorization:
$\kappa_{\text{above}} F^{\text{above}}_k$, $\kappa_{\text{loans}}
F^{\text{loans}}_k$, and $\kappa_{\text{guar}} F^{\text{guar}}_k$ enter
$b_{k+1}$ contemporaneously. Direct income support is disbursement-based on both
margins: it reaches aggregate spending with the one-quarter delay identified
empirically in Section~\ref{sec:estimation}
($\alpha_{\text{DI}}\, z_k$ in $y_{k+1}$) and hits the budget at the same
disbursement date ($\kappa_{\text{DI}}\, z_k$ in $b_{k+1}$).

\subsection{Jacobian Matrices}
\label{app:jacobians}

Linearizing around a nominal trajectory $(\bar{\bm{x}}_k, \bar{\bm{u}}_k)$ gives
$\delta\bm{x}_{k+1} \approx \bm{A}_k\,\delta\bm{x}_k + \bm{B}_k\,\delta\bm{u}_k$.
The system is nonlinear through the bilinear terms
$\alpha_{S,\text{DI}}\, S_k F^{\text{DI}}_{k-1}$ in output and
$\phi_S\, S_k \theta_k$ in infection, and through the complementarity multiplier
$m(K_{k+1})$, so $\bm{A}_k$ and $\bm{B}_k$ depend on the nominal point. Write
$\bar m_k = m(\bar K_{k+1})$ and
$\bar m'_k = \chi \bar K / (\bar K_{k+1} + \bar K)^2$ for the multiplier and its
derivative evaluated at the nominal end-of-period capacity stock. With states
ordered as in Eq.~(\ref{eq:state7}) and controls as
$(S, F^{\text{above}}, F^{\text{loans}}, F^{\text{guar}}, F^{\text{DI}})$, the
nonzero entries are
%
\begin{align}
    \bm{A}_k:\quad
    & A_{yy} = \rho_y, \quad
      A_{yd} = -\beta_d, \quad
      A_{y z} = \alpha_{\text{DI}} + \alpha_{S,\text{DI}}\,\bar{S}_k, \notag \\
    & A_{y L} = \alpha_L\,\bar m_k\,(1-\delta_L), \quad
      A_{y K} = \bigl(\alpha_K + \alpha_L\,\bar L_{k+1}\,\bar m'_k\bigr)(1-\delta_K), \notag \\
    & A_{by} = -\gamma_y, \quad A_{bb} = 1+r, \quad A_{bz} = \kappa_{\text{DI}}, \notag \\
    & A_{\theta\theta} = \rho_\theta(k)\bigl(1 - \phi_S\,\bar{S}_k/100\bigr), \quad
      A_{d\theta} = \delta_\theta(k), \notag \\
    & A_{LL} = 1-\delta_L, \quad A_{KK} = 1-\delta_K,
    \label{eq:Ak}
\end{align}
%
\begin{align}
    \bm{B}_k:\quad
    & B_{yS} = \alpha_S + \alpha_{S,\text{DI}}\,\bar{z}_k, \quad
      B_{y,\text{ab}} = \alpha_K + \alpha_L\,\bar L_{k+1}\,\bar m'_k, \notag \\
    & B_{y,\text{lo}} = B_{y,\text{gu}} = \alpha_L\,\bar m_k, \notag \\
    & B_{b,\text{ab}} = \kappa_{\text{above}}, \quad
      B_{b,\text{lo}} = \kappa_{\text{loans}}, \quad
      B_{b,\text{gu}} = \kappa_{\text{guar}}, \notag \\
    & B_{\theta S} = -\rho_\theta(k)\,\phi_S\,\bar{\theta}_k/100, \quad
      B_{z,\text{DI}} = 1, \notag \\
    & B_{L,\text{lo}} = B_{L,\text{gu}} = 1, \quad
      B_{K,\text{ab}} = 1 .
    \label{eq:Bk}
\end{align}
%
The nonlinear couplings appear in three places. First, $A_{yz}$ carries
$\alpha_{S,\text{DI}}\,\bar{S}_k$ and $B_{yS}$ carries
$\alpha_{S,\text{DI}}\,\bar{z}_k$, the two partial derivatives of the
$S_k F^{\text{DI}}_{k-1}$ term. Second, $B_{\theta S}$ carries the derivative of
$S_k\theta_k$; the stringency entries scale by $1/100$ because containment enters
the infection equation in fractional units while entering output in points.
Third, the capacity entries $A_{yK}$ and $B_{y,\text{ab}}$ combine the direct
effect $\alpha_K$ with the indirect effect through the multiplier,
$\alpha_L \bar L_{k+1} \bar m'_k$: because $m$ is bounded, $\bar m'_k$ shrinks as
capacity saturates, so the marginal return to above-the-line support declines in
the outstanding stock.

\subsection{The iLQR Algorithm}
\label{app:ilqr}

I solve the problem with the iterative Linear-Quadratic Regulator (iLQR) of
\citet{li2004iterative}, which successively linearizes the dynamics around the
current trajectory and solves the resulting time-varying linear-quadratic (LQ)
subproblem in closed form.
\citet{sharma2025sqp} show iLQR is equivalent to Sequential Quadratic Programming
on the discrete-time control problem, with cost descent at each iteration; the
notation follows \citet[Chapters~3--4]{bertsekas2017dynamic}.

\paragraph{Backward pass.}
Let $V_x, V_{xx}$ denote the gradient and Hessian of the cost-to-go. Starting from
the terminal condition (Section~\ref{app:terminal}), the recursion at each $k$ is
%
\begin{align}
    \bm{Q}_x  &= \bm{A}_k'\,V_x', &
    \bm{Q}_u  &= \bm{l}_u + \bm{B}_k'\,V_x', \notag \\
    \bm{Q}_{xx} &= \bm{A}_k'\,V_{xx}'\,\bm{A}_k, &
    \bm{Q}_{uu} &= \bm{l}_{uu} + \bm{B}_k'\,V_{xx}'\,\bm{B}_k + \mu\,\bm{D}, \\
    \bm{Q}_{ux} &= \bm{B}_k'\,V_{xx}'\,\bm{A}_k, \notag
\end{align}
%
where $V_x' = \bm{l}_x + V_x$ and $V_{xx}' = \bm{l}_{xx} + V_{xx}$ fold in the
running cost, and $\bm{l}_x, \bm{l}_u, \bm{l}_{xx}, \bm{l}_{uu}$ are the
discounted gradients and Hessians of the stage loss. Because the stage cost is
separable in $(\bm{x}, \bm{u})$, there is no cost cross term ($\bm{l}_{ux} = 0$),
which keeps $\bm{Q}_{uu}$ well-conditioned. The feedforward and feedback gains are
$\bm{k}_k = -\bm{Q}_{uu}^{-1}\bm{Q}_u$ and
$\bm{K}_k = -\bm{Q}_{uu}^{-1}\bm{Q}_{ux}$.

\paragraph{Regularization.}
The explosive eigenvalues $1+r > 1$ and
$\rho_\theta(k)(1-\phi_S\bar{S}_k/100) > 1$ in early waves inflate $V_{xx}$ and
can render $\bm{Q}_{uu}$ indefinite. Following \citet{tassa2014control}, I add
Levenberg--Marquardt damping $\mu\,\bm{D}$ with
$\bm{D} = \mathrm{diag}(\bar u_1^{-2}, \ldots, \bar u_5^{-2})$, scaled by the
instrument caps so the damping is uniform in normalized control units despite
the instruments' very different natural scales. The parameter $\mu$ is increased
by a factor of ten when the Cholesky factorization of $\bm{Q}_{uu}$ fails, and
decreased geometrically after each accepted iteration.

\paragraph{Forward pass.}
Given the gains, the updated controls are applied to the \emph{nonlinear} dynamics
with line-search step $\alpha \in (0,1]$:
%
\begin{equation}
    \bm{u}_k^{\text{new}} = \bm{u}_k + \alpha\,\bm{k}_k
    + \bm{K}_k\bigl(\bm{x}_k^{\text{new}} - \bm{x}_k\bigr),
    \label{eq:fwd}
\end{equation}
%
starting from $\bm{x}_0^{\text{new}} = \bm{x}_0$. The step $\alpha$ is halved from
unity until total cost decreases; if no decrease is found, $\mu$ is increased and
the backward pass repeats.

\paragraph{Constraints.}
Stringency is bounded by its natural scale, $S_k \in [0, 100]$; the four fiscal
instruments are bounded below at zero and above at the 99th percentile of the
observed positive values for each instrument; and the Q4.2019 controls are locked
at zero. The bounds are enforced by projection (clipping) in the forward pass. In
addition, quarter-on-quarter changes in stringency are capped at
$|S_k - S_{k-1}| \le \Delta S_{\max}$, where $\Delta S_{\max}$ is the 95th
percentile of observed $|\Delta S|$ across the panel, enforced by clamping $S_k$
relative to the previously updated $S_{k-1}$ within the same forward pass. This
speed limit is imposed as a hard constraint rather than a quadratic penalty: a
$\Delta S$ penalty would require augmenting the state with $S_{k-1}$ and
introduce a cost cross term ($\bm{l}_{ux} \neq 0$) that, combined with the
bilinear dynamics, ill-conditions $\bm{Q}_{uu}$. The constraint has no such cross
term and is economically transparent---containment cannot ramp faster than
governments managed in-sample.

\paragraph{Certainty equivalence.}
At each iteration the subproblem is linear-quadratic, so certainty equivalence
holds \citep[Chapter~3.1]{bertsekas2017dynamic}: the feedback gains are identical
to the deterministic case, and the epidemiological shocks $\varepsilon_k$ raise
expected cost without altering the optimal policy rule. The realized shock path is
fed through the nonlinear forward pass.

\paragraph{Convergence and global optimality.}
The algorithm terminates when the relative cost reduction falls below
$10^{-9}$. The bilinear terms and the complementarity multiplier make the problem
non-convex, so iLQR converges to a local optimum. I therefore restart the solver
from four economically distinct initializations---the observed average policy,
the zero-control path, and two corner allocations tilted toward above- and
below-the-line support---and report the solution with the lowest cost.
Trajectories whose backward pass terminates with saturated regularization
($\mu$ at its ceiling) are treated as failed exits rather than solutions.
[CHECK: verifizieren, ob im Baseline-Run alle vier Starts auf dasselbe $J^*$
konvergieren; falls ja, entsprechenden Satz wieder aufnehmen.]

\subsection{Terminal Condition}
\label{app:terminal}

The terminal cost weights public debt alone. By the end of the horizon,
vaccination has severed the link from infection to death and removed the output
cost of containment, so the only state that persists as a welfare-relevant legacy
is accumulated debt. The terminal weight capitalizes the discounted stream of
future debt-loss deviations under a stationary post-pandemic continuation,
$Q_N^{(b,b)} = M\,\tau_b\, w_b$ with $M = 1/(1-\beta\,\lambda^2)$, where $\lambda = 0.99$
is the assumed post-horizon persistence of the terminal debt deviation. All other
terminal weights are zero.

\subsection{Computation of the Constrained Policy Frontier}
\label{app:frontier_comp}

The frontier problems of Section~\ref{sec:frontier} are solved as exact
nonlinearly constrained programs over the full control path
$\bm{u} \in \mathbb{R}^{5 \times N}$, using sequential quadratic programming. In
each scenario, one of the three welfare-relevant outcomes---discounted output
loss, cumulative excess mortality, or terminal debt---is minimized subject to
the other two being no worse than their model-implied values under the observed
policy, evaluated through the full nonlinear transition system. The box bounds
and the ramping constraint of Section~\ref{app:ilqr} are imposed as bounds and
nonlinear inequality constraints, respectively; a vanishing tie-breaking penalty
on normalized instrument use ($10^{-8}$) selects among cost-equivalent policies
and carries no welfare interpretation. Because the constrained problems inherit
the non-convexity of the dynamics, each scenario is solved from five starting
policies---observed policy, the weighted-planner solution, zero policy, and two
mixed allocations tilted toward above-the-line support---and the best feasible
local solution across starts is reported, with feasibility verified ex post
against the observed thresholds.

% =============================================================================
\subsection{Eigenvalue Structure and Fiscal Self-Financing}
\label{app:stability}
% =============================================================================

\subsubsection*{Eigenvalues}

The augmented system matrix $\bm{A}_k$ in Eq.~(\ref{eq:Ak}) is block-triangular, so
its eigenvalues are its diagonal entries. Two are zero (the bookkeeping states
$d$ and $z$, which carry no own-persistence), and five govern the dynamics:
%
\begin{align}
    \lambda_y      &= \rho_y < 1 && \text{(output: self-stabilizing)}, \\
    \lambda_b      &= 1 + r > 1 && \text{(debt: unconditionally unstable)}, \\
    \lambda_\theta &= \rho_\theta(k)\bigl(1 - \phi_S\, S_k/100\bigr)
                      && \text{(infection: conditionally unstable)}, \\
    \lambda_L      &= 1 - \delta_L < 1 && \text{(liquidity stock: stable)}, \\
    \lambda_K      &= 1 - \delta_K < 1 && \text{(capacity stock: stable)}.
\end{align}
%
Two features drive the planner's problem. First, the debt eigenvalue exceeds unity
for any positive interest rate and is invariant to policy: no control enters
$A_{bb}$, so debt cannot be stabilized by any available instrument. During the
crisis both forces on debt are explosive---the automatic stabilizer
$-\gamma_y\, y_k > 0$ when output is below potential, and discretionary support
$\kappa\cdot F_k > 0$. Second, the infection eigenvalue exceeds unity, and the
epidemic grows, whenever stringency falls short of the threshold
%
\begin{equation}
    \lambda_\theta \ge 1
    \quad\Longleftrightarrow\quad
    S_k < \underline{S}(k) \equiv
    \frac{100}{\phi_S}\Bigl(1 - \tfrac{1}{\rho_\theta(k)}\Bigr).
    \label{eq:S_threshold}
\end{equation}
%
The threshold is wave-specific through $\rho_\theta(k)$ but independent of current
prevalence: the \emph{level} of $\theta_k$ sets the death toll, while the
\emph{stability} of the epidemiological state depends only on $\rho_\theta$ and
$\phi_S$. At the pre-vaccination calibration ($\rho_\theta = 1.5$,
$\phi_S = 0.8$), $\underline{S} \approx 42$ points: keeping the epidemic from
diverging requires substantial containment, of the order of the average
stringency actually observed during the acute waves. Once vaccination lowers
$\rho_\theta$ below unity, $\underline{S} \le 0$, the infection state is stable
without intervention, and Eq.~(\ref{eq:S_threshold}) no longer binds.

The planner therefore faces a cascade: epidemiological instability requires
containment; containment depresses output through $\alpha_S$; the output gap widens
the deficit through $-\gamma_y\, y_k$; and the resulting debt cannot be stabilized
by any control. Each stabilizing action on one dimension feeds the instability of
another. This is the trilemma as a dynamical-systems property.

\subsubsection*{Fiscal self-financing}

A fiscal instrument is self-financing if the automatic-stabilizer savings from the
output it generates outweigh its direct budgetary cost. Consider demand injection,
which is disbursement-based on both margins. Authorized in quarter $0$, it reaches
the budget and output at disbursement one quarter later:
$\kappa_{\text{DI}}\, z_1 = \kappa_{\text{DI}}\, F^{\text{DI}}_0$ enters $b_2$ and
$\alpha_{\text{DI}}\, z_1$ enters $y_2$. The higher $y_2$ feeds back into the
budget through the automatic stabilizer, reducing $b_3$ by
$\gamma_y\,\alpha_{\text{DI}}$. The net debt effect is
%
\begin{equation}
    \frac{\partial b_3}{\partial F^{\text{DI}}_0}
    = (1+r)\,\kappa_{\text{DI}} - \gamma_y\,\alpha_{\text{DI}},
    \label{eq:DI_selffinance}
\end{equation}
%
so demand injection is self-financing when
$\gamma_y\,\alpha_{\text{DI}} \ge (1+r)\,\kappa_{\text{DI}}$. At the estimated
parameters, $\gamma_y\,\alpha_{\text{DI}} = 0.17$ against
$(1+r)\,\kappa_{\text{DI}} = 0.53$.

Above-the-line support operates through the persistent capacity stock, so the
relevant comparison credits the entire output stream the stock generates.
Authorized in quarter $0$, the flow commits the budget contemporaneously
($\kappa_{\text{above}} F^{\text{above}}_0$ in $b_1$) and raises output through
$\alpha_K K_{k+1}$ in every subsequent quarter as the stock decays. Ignoring
discounting and interest compounding---both of which work against
self-financing---the cumulative stabilizer saving per unit of flow is bounded by
$\gamma_y\,\alpha_K \sum_{j\ge 0}(1-\delta_K)^j = \gamma_y\,\alpha_K/\delta_K
\approx 0.18$, against a budgetary cost of $\kappa_{\text{above}} = 0.66$. The
same bound for below-the-line liquidity, credited with the full complementarity
uplift $m \le 1+\chi$, gives
$\gamma_y\,\alpha_L(1+\chi)/\delta_L \approx 0.46$ against
$\kappa_{\text{loans}} = 0.84$ and $\kappa_{\text{guar}} = 0.54$: still short of
self-financing, but by the smallest margin of any instrument---a first indication
of why the frontier recomposes support toward the below-the-line channel.

Every instrument therefore leaves a positive net debt footprint even under
assumptions deliberately tilted in its favor. This is the analytical counterpart
to the paper's central result. Because no instrument repays itself through output
and debt cannot be stabilized ($\lambda_b > 1$), fiscal policy cannot resolve the
trilemma---it can only redistribute the burden across instruments and time. The
planner's gains come not from making support self-financing but from
\emph{recomposing} it toward instruments with a smaller net debt footprint per
unit of output and mortality relief. The normative content of the model is debt
efficiency at given output and health outcomes, not output stabilization through
fiscal expansion.

\section{Complementary Results}
\label{app:compl_results}
% -----------------------------------------------------------------------------
\subsection{The Weighted Welfare Benchmark}
\label{subsec:results_weighted}
% Inputs: table_weighted_sweep.csv
%         table_welfare_gain_decomposition_baseline.csv
%         fig_welfare_gain_decomposition_baseline_aerstyle.pdf
% -----------------------------------------------------------------------------
Table~\ref{tab:weight_sweep} reports the planner solution across the four preference scenarios and the VSL-anchored \(\lambda_d\) grid.

Three patterns stand out. First, the composition of the optimal policy is
invariant across all seven preference configurations: effective liquidity
support dominates the fiscal package (15.5 to 21.5 percent of GDP against 4.6
to 6.7 above the line), demand injection is never used, and containment is
sustained well above the observed average---the recomposition is a property of
the transition system, not of any particular weighting. Second, the
comparative statics move as theory predicts: a higher value of life shifts the
solution monotonically along the trilemma, buying fewer deaths (1,965 down to
1,889 per million across the VSL grid) with tighter containment, a larger
fiscal envelope, and more terminal debt, while a doubled debt weight
(\(\tau_b = 0.10\)) compresses fiscal support and accepts a substantially
worse output path in exchange for the lowest debt level. Notably, the
precautionary planner ends with more deaths than the baseline despite
twice the mortality weight: the tighter fiscal stance raises the output cost
of containment, illustrating that the instruments are not separable within the
trilemma. Third, the welfare gain over the observed policy is not a knife-edge
result---it ranges from 30 to 46 percent of the observed loss across the
scenarios, with the preferred baseline at 35 percent.


\begin{table}[H]
\centering
\scriptsize
\caption{Weighted planner across preference scenarios and the VSL sweep.}
\label{tab:weight_sweep}
\begin{tabular}{lcccccccc}
\toprule
\textbf{Scenario} & $\tau_b$ & $\lambda_d$ & Cum. $y$ & Deaths/M & $b_N$ &
Total above & Total below & Gain (\%) \\
\midrule
Output-oriented  & 0.025 & 25  & $-8.1$  & $2{,}013$ & $82.0$ & $6.7$ & $21.5$ & $46.2$ \\
Balanced         & 0.050 & 50  & $-15.1$ & $1{,}987$ & $79.5$ & $5.8$ & $18.7$ & $33.9$ \\
Baseline         & 0.050 & 75  & $-16.5$ & $1{,}938$ & $80.3$ & $6.0$ & $19.3$ & $34.6$ \\
Precautionary    & 0.100 & 150 & $-30.2$ & $1{,}943$ & $77.3$ & $4.6$ & $15.5$ & $30.2$ \\
\midrule
VSL \$5M         & 0.050 & 58  & $-15.5$ & $1{,}965$ & $79.8$ & $5.9$ & $18.9$ & $34.1$ \\
VSL \$9M         & 0.050 & 104 & $-17.8$ & $1{,}907$ & $81.3$ & $6.2$ & $20.5$ & $35.3$ \\
VSL \$11M        & 0.050 & 127 & $-18.8$ & $1{,}889$ & $81.8$ & $6.4$ & $21.1$ & $35.8$ \\
\bottomrule
\end{tabular}
\vspace{0.15cm}
\begin{minipage}{0.95\textwidth}
\scriptsize \textit{Notes:} Each row solves the planner problem of
Section~\ref{sec:planner} under the stated preferences. Cum.\ \(y\) in
pp-quarters; \(b_N\) in percent of GDP; fiscal totals in percent of 2019 GDP,
below-the-line in take-up-adjusted effective terms. Demand injection is zero in
every scenario and is omitted. Gain in percent of the observed total welfare
loss evaluated under the same weights. VSL mapping of \(\lambda_d\): see
Section~\ref{sec:planner}.
\end{minipage}
\end{table}

The gain is carried by the two dimensions the planner deliberately prioritizes
and paid for by the one it deliberately sacrifices: lower mortality
contributes 73 percent of the total gain and the improved output path 56
percent, while higher running debt and the larger terminal debt burden offset
29 percent. This is the weighted counterpart of the frontier logic---the same
liquidity-led recomposition that the frontier identifies without weights, here
pushed beyond the observed outcome levels: unlike the frontier solutions, the
weighted optimum is not itself a Pareto improvement, since plausible
preferences find it worthwhile to buy additional lives and output at a debt
level above the observed one. The frontier establishes that improvement
required no sacrifice; the weighted benchmark shows where welfare-maximizing
preferences would go beyond it.


\begin{figure}[H]
    \centering
    \captionsetup{font=scriptsize}
    \includegraphics[width=0.7\textwidth]{fig_welfare_gain_decomposition_baseline_aerstyle.pdf}
    \caption{Decomposition of the baseline welfare gain by loss component.
    \textit{Notes:} [Notiz: in \% des beobachteten Gesamtverlusts; positive
    Werte = geringerer Verlust unter dem Planner.]}
    \label{fig:welfare_gain}
\end{figure}

% =============================================================================
% APPENDIX: CROSS-COUNTRY ANALYSIS -- technische Details zu Subsection 6.3
% =============================================================================

\section{Cross-Country Analysis: Implementation and Diagnostics}
\label{app:country_analysis}

\subsection{Country-Level Implementation}
\label{app:country_implementation}

For each country \(i\), the three frontier programs of
Eq.~(\ref{eq:frontier_program}) are solved with the common structural
parameters of Table~\ref{tab:params} and the common policy menu---the
in-support caps and the ramping limit \(\Delta S_{\max}\) of
Section~\ref{sec:planner}, both computed from the pooled OECD panel---while
five ingredients are country-specific: the epidemiological innovation path
\(\varepsilon_{\theta,i}\), backed out from the country's own estimated
prevalence series under the common transmission parameters; the onset shock
\(\varepsilon_{y,i}\) in Q2.2020; the country intercepts \(\mu_{y,i}\) and
\(\mu_{b,i}\) from the panel fixed effects; and the initial debt stock
\(b_{0,i}\), which also serves as the country's debt target. The constraint
thresholds \((\mathcal{L}_y^{\text{obs},i}, D^{\text{obs},i},
b_N^{\text{obs},i})\) are the model-implied outcomes of the country's
\emph{unconstrained} observed policy: the menu restricts only the
counterfactual, never the evaluation of the country itself. Debt is reported
throughout as the change relative to \(b_{0,i}\), since initial levels
ranging from 7 to 205 percent of GDP are not comparable across countries.
Each program is solved by SQP from up to five starting policies, including
the country's weighted-planner solution, with feasibility verified ex post
against all constraints (outcome thresholds, box bounds, and ramping);
solutions failing any ex-post check are discarded
(Appendix~\ref{app:frontier_comp}).

\subsection{Invariance of Frontier Distances to Mortality-Level Misfit}
\label{app:ifr_invariance}

The common infection-fatality sequence \(\delta_q\) over- or under-predicts
cumulative deaths for individual countries; the multiplicative scale factor
\(s_i\) that matches model deaths to the raw data has a median of 0.89 and a
range of \([0.38,\, 1.59]\). The following result shows that this misfit
cannot contaminate the frontier distances.

\begin{proposition}
Under the maintained \(\beta_d = 0\), the feasible sets and optimal policies
of all three frontier programs are invariant to the country-specific
rescaling \(\delta_{q,i} = s_i\, \delta_q\), \(s_i > 0\). Reported death
levels scale by \(s_i\); the percentage mortality improvement and the output
and debt improvements are unchanged.
\end{proposition}

\begin{proof}
Deaths are a linear readout of prevalence, \(d_{k+1} = \delta_q(k)\,
\theta_k\), and the prevalence dynamics do not involve \(\delta_q\). With
\(\beta_d = 0\), neither the output nor the debt transition involves
\(d_k\). Hence for any policy path \(\bm{u}\), rescaling \(\delta_q\) by
\(s_i\) scales the entire death path, and thus \(D(\bm{u})\), by \(s_i\),
while \(\mathcal{L}_y(\bm{u})\) and \(b_N(\bm{u})\) are unchanged. In
scenarios A and C the mortality constraint becomes \(s_i D(\bm{u}) \le s_i
D^{\text{obs}}\), which defines the same feasible set; objectives are
unaffected. In scenario B the objective is scaled by the positive constant
\(s_i\), leaving the minimizer unchanged, and the percentage improvement
\((s_i D^{\text{obs}} - s_i D^{*})/(s_i D^{\text{obs}})\) is invariant.
\end{proof}

The invariance is exact for multiplicative, time-constant rescaling; it does
not cover misfit in the \emph{time profile} of \(\delta_q\), and it relies on
the absence of a mortality feedback into output (\(\beta_d = 0\)). It was
verified numerically for the two worst-fitting countries (Switzerland,
\(s = 0.45\), and the Netherlands, \(s = 0.38\)): the frontier improvements
under the rescaled sequence coincide with the baseline to within
\(5 \times 10^{-12}\), i.e., to solver precision. A consequence worth noting
is that the mortality-fit diagnostics reported below matter only for the
\emph{reported death levels} and for the weighted-planner metric, in whose
loss \(\lambda_d w_d d^2\) the rescaling enters quadratically---not for the
frontier distances on which the ranking is based.

\subsection{Fit Diagnostics and Ranking Stability}
\label{app:country_fit}

Model fit is summarized by a continuous score,
\(\text{fit}_i = \max\!\big(|\log \text{DR}_i|/\text{MAD}_{|\log DR|},\;
|\Delta y_i|/\text{MAD}_{|\Delta y|}\big)\), where \(\text{DR}_i\) is the
ratio of model-implied to raw cumulative deaths and \(\Delta y_i\) the gap
between model-implied and raw cumulative output, each normalized by the
panel median absolute deviation. The worst-fitting quartile comprises
Belgium, Switzerland, Colombia, Spain, Estonia, Iceland, Lithuania, the
Netherlands, and T\"urkiye; the dominant pattern
is over-prediction of deaths under the common infection-fatality sequence,
which by the proposition above does not affect the frontier distances.
Output misfit, which is not covered by the invariance, is material for three
countries (COL, EST, TUR). The ranking is insensitive to fit-based
exclusion: re-ranking after dropping the five (ten) worst-fitting countries
yields a Spearman rank correlation of 0.998 (0.997) with the baseline on the
surviving sample. Country-specific rescaling of the infection-fatality
sequence is thus unnecessary for the ranking; extending the calibration to
country-specific mortality profiles remains a natural refinement for the
reported death levels.

\subsection{Demand Injection by Containment State}
\label{app:di_decomposition}

The exchange-rate gradient reported in
Section~\ref{subsec:results_countries} is sharpened by decomposing each
country's demand injection according to the containment state in which it
becomes effective. Because demand injection acts with a one-quarter lag, the
relevant state for support committed in quarter \(k\) is stringency in
\(k+1\), where the multiplier \(\alpha_{\text{DI}} +
\alpha_{S,\text{DI}}\,S_{k+1}\) applies; the decomposition therefore splits
each package's demand injection into the share deployed above and below the
threshold \(S^{*} = \alpha_{\text{DI}}/|\alpha_{S,\text{DI}}| \approx 36\)
at the effective quarter. Both components vary substantially across
countries: the below-threshold share ranges from zero to 0.28 of the total
package, with ten countries above 0.05. The model's prediction is that only
the above-threshold component should be costly, and the cross-section
confirms it exactly: the realized fiscal exchange rate correlates at
\(-0.89\) with the above-threshold component---identical to the correlation
with the total demand-injection share---and at \(+0.03\) with the
below-threshold component. Finland illustrates the split: it deployed one of
the larger demand-injection shares of the frontier-near group, but over four
fifths of it below the threshold, and realized an above-median exchange
rate. The decomposition also speaks against a country-characteristics
confound: an institutional or informality account of the exchange-rate
gradient would burden both components alike, whereas the data assign the
entire penalty to the component deployed in the state where the estimated
multiplier is negative.

\subsection{Special Cases: Off-Menu Policies and the Output Ceiling}
\label{app:country_special}

Seven countries committed at least one instrument beyond the common caps in
at least one quarter: concentrated wage-subsidy commitments in Q1.2020
(Australia, New Zealand), a demand-injection spike in Q1.2021 (Chile), a
loan-program spike in Q2.2020 (Denmark), and sustained near-maximal
containment (Colombia in Q2.2020; Ireland and the United Kingdom in
Q1.2021); Table~\ref{tab:country_exceedance} reports the instrument, quarter,
and magnitude. All seven are nonetheless dominated from within the common
menu, so their off-menu intensity was dispensable for their realized
outcomes. Italy is the converse case: sixteen countries exceeded the ramping
limit \(\Delta S_{\max}\) in Q2.2020, but only for Italy is this speed
pivotal---without it, the mortality threshold is unattainable from the menu
and the frontier programs are infeasible. Finally, seven countries reach the
output ceiling (\(\text{ImpA} \ge 99\) percent): for Australia, Denmark,
Israel, Japan, Korea, New Zealand, and the United States, the frontier
eliminates the policy-amenable output loss entirely, so their output
distances are bounded by the transition technology rather than by fiscal
space, and their positions in the ranking should be read accordingly.

\begin{table}[H]
\centering
\scriptsize
\caption{Observed policies beyond the common menu: instrument, quarter, and magnitude.}
\label{tab:country_exceedance}
\begin{tabular}{llcc}
\toprule
 & \textbf{Instrument} & \textbf{Quarter} & \textbf{Exceedance} \\
\midrule
Australia      & Above-the-line (wage subsidy) & Q1.2020 & $3.53$ \\
New Zealand    & Above-the-line (wage subsidy) & Q1.2020 & $2.04$ \\
Chile          & Demand injection              & Q1.2021 & $4.65$ \\
Denmark        & Loans                         & Q2.2020 & $0.67$ \\
Colombia       & Stringency                    & Q2.2020 & $7.61$ \\
Ireland        & Stringency                    & Q1.2021 & $6.27$ \\
United Kingdom & Stringency                    & Q1.2021 & $1.61$ \\
\bottomrule
\end{tabular}
\vspace{0.15cm}
\begin{minipage}{0.9\textwidth}
\scriptsize \textit{Notes:} Largest exceedance of the common in-support caps
per country: fiscal instruments in percent of 2019 GDP above the cap
(effective, take-up-adjusted terms), stringency in index points above the
cap of 78.0. Fiscal exceedances reflect commitments in the quarter of
announcement. All seven countries are nonetheless dominated from within the
common menu (Section~\ref{subsec:results_countries}).
\end{minipage}
\end{table}

\subsection{Full Country Tables}
\label{app:country_tables}

\begin{table}[H]
\centering
\scriptsize
\caption{Distance to the country-specific frontier: all countries.}
\label{tab:country_ranking_full}
\begin{tabular}{lc cc !{\vrule width 0.4pt} c}
\toprule
 & \textbf{Frontier distance} & \multicolumn{2}{c}{\emph{Ranking dimensions}} & \emph{Context} \\
\cmidrule(lr){3-4}
 & (avg.\ rank, 1--37) & \textbf{Output impr.} & \textbf{Debt impr.} & \textbf{Deaths impr.} \\
 & & (\% of own \(\mathcal{L}_y\)) & (pp of GDP) & (\% of own deaths) \\
\midrule
Mexico                    & $1.0$  & $35.3$  & $2.2$  & $16.6$ \\
Sweden                    & $4.0$  & $58.5$  & $3.7$  & $21.3$ \\
Portugal                  & $5.0$  & $48.7$  & $5.0$  & $16.2$ \\
Finland                   & $6.0$  & $62.3$  & $3.7$  & $13.3$ \\
France                    & $7.0$  & $52.1$  & $5.6$  & $13.3$ \\
Norway                    & $7.5$  & $77.9$  & $3.0$  & $12.5$ \\
Iceland\(^{+}\)           & $9.5$  & $60.8$  & $5.6$  & $28.5$ \\
Belgium\(^{+}\)           & $11.0$ & $61.3$  & $5.6$  & $22.0$ \\
Luxembourg                & $11.5$ & $83.9$  & $4.0$  & $22.0$ \\
Switzerland\(^{+}\)       & $12.0$ & $82.1$  & $4.5$  & $23.5$ \\
Colombia\(^{\#+}\)        & $12.5$ & $80.7$  & $5.1$  & $19.9$ \\
Netherlands\(^{+}\)       & $13.0$ & $69.2$  & $5.7$  & $17.1$ \\
Hungary                   & $14.0$ & $80.5$  & $5.6$  & $17.0$ \\
Spain\(^{+}\)             & $14.5$ & $59.0$  & $7.7$  & $22.2$ \\
Austria                   & $15.0$ & $67.9$  & $6.3$  & $16.4$ \\
Costa Rica                & $16.0$ & $89.2$  & $5.3$  & $21.9$ \\
Germany                   & $16.5$ & $76.7$  & $6.6$  & $15.6$ \\
United Kingdom\(^{\#}\)   & $17.0$ & $59.7$  & $8.2$  & $21.9$ \\
T\"urkiye\(^{+}\)         & $17.5$ & $96.0$  & $4.0$  & $21.7$ \\
Czechia                   & $18.5$ & $80.3$  & $7.6$  & $17.6$ \\
Latvia                    & $19.0$ & $89.1$  & $5.9$  & $18.9$ \\
Slovak Republic           & $19.0$ & $88.9$  & $6.0$  & $18.1$ \\
Estonia\(^{+}\)           & $21.0$ & $90.6$  & $6.1$  & $20.7$ \\
Lithuania\(^{+}\)         & $22.0$ & $94.8$  & $5.9$  & $20.2$ \\
Ireland\(^{\#}\)          & $23.0$ & $90.9$  & $7.0$  & $16.3$ \\
Slovenia                  & $25.0$ & $87.8$  & $8.6$  & $15.0$ \\
Greece                    & $28.5$ & $93.2$  & $8.6$  & $14.2$ \\
Korea\(^{*}\)             & $28.5$ & $99.9$  & $7.8$  & $24.0$ \\
Poland                    & $28.5$ & $96.3$  & $7.9$  & $19.6$ \\
Canada                    & $29.0$ & $91.8$  & $9.2$  & $27.0$ \\
Israel\(^{*}\)            & $30.5$ & $100.0$ & $7.8$  & $27.5$ \\
Chile\(^{\#}\)            & $31.5$ & $92.9$  & $11.7$ & $27.4$ \\
Denmark\(^{*\#}\)         & $32.0$ & $100.0$ & $8.3$  & $20.6$ \\
Australia\(^{*\#}\)       & $33.5$ & $100.0$ & $8.9$  & $28.6$ \\
Japan\(^{*}\)             & $33.5$ & $99.6$  & $10.8$ & $21.2$ \\
New Zealand\(^{*\#}\)     & $34.5$ & $100.0$ & $10.2$ & $38.7$ \\
United States\(^{*}\)     & $35.0$ & $100.0$ & $10.6$ & $34.6$ \\
\midrule
Italy\(^{\dagger}\)       & ---    & ---     & ---    & ---    \\
\bottomrule
\end{tabular}
\vspace{0.15cm}
\begin{minipage}{0.95\textwidth}
\scriptsize \textit{Notes:} Countries in ascending order of frontier distance,
the average of the within-sample ranks in the output and debt improvements
(37 ranked countries; lower = closer to the own frontier). Each improvement
column reports a separate frontier program: the output and debt improvements
hold the country's mortality (and the respective other dimension) at its
observed level; the deaths improvement is from the program that instead
minimizes mortality at the observed output loss and debt---it is contextual
and does not enter the distance. Debt improvements in absolute pp of GDP; the
distance uses within-column ranks and is invariant to any monotone choice of
units. \(^{*}\)~Output ceiling reached (\(\text{ImpA} \ge 99\) percent).
\(^{\#}\)~Observed policy committed at least one instrument beyond the common
menu (Table~\ref{tab:country_exceedance}). \(^{+}\)~Worst-fitting quartile of
the continuous fit score (Appendix~\ref{app:country_fit}); by the invariance
result of Appendix~\ref{app:ifr_invariance}, mortality-level misfit does not
affect the reported distances. \(^{\dagger}\)~Italy: observed policy not
replicable within the common menu (first-wave ramp-up exceeds
\(\Delta S_{\max}\)); the output and debt programs are infeasible and the
mortality program cannot improve on the observed outcome, so no distance is
defined.
\end{minipage}
\end{table}


\begin{table}[H]
\centering
\scriptsize
\caption{Frontier ranking and realized outcomes in the raw data.}
\label{tab:country_raw_outcomes}
\begin{tabular}{lc ccc}
\toprule
 & \textbf{Frontier distance} & \textbf{Cum.\ output gap} & \textbf{Excess deaths} & \textbf{Debt change} \\
 & (avg.\ rank, 1--37) & (pp-quarters, raw) & (per million, raw) & (pp of GDP, raw) \\
\midrule
Mexico                  & $1.0$  & $-101.2$ & $2{,}630$ & $+0.9$  \\
Sweden                  & $4.0$  & $-29.6$  & $2{,}169$ & $+1.9$  \\
Portugal                & $5.0$  & $-95.0$  & $2{,}491$ & $-2.7$  \\
Finland                 & $6.0$  & $-57.3$  & $1{,}418$ & $+6.2$  \\
France                  & $7.0$  & $-68.5$  & $2{,}407$ & $+6.8$  \\
Norway                  & $7.5$  & $-17.6$  & $869$     & $+4.0$  \\
Iceland\(^{+}\)         & $9.5$  & $-105.1$ & $766$     & $+10.7$ \\
Belgium\(^{+}\)         & $11.0$ & $-33.0$  & $2{,}949$ & $+3.7$  \\
Luxembourg              & $11.5$ & $-13.4$  & $1{,}925$ & $+7.5$  \\
Switzerland\(^{+}\)     & $12.0$ & $-17.6$  & $1{,}730$ & $+2.1$  \\
Colombia\(^{\#+}\)      & $12.5$ & $-27.2$  & $2{,}899$ & $+16.1$ \\
Netherlands\(^{+}\)     & $13.0$ & $-36.2$  & $1{,}391$ & $+5.5$  \\
Hungary                 & $14.0$ & $-82.2$  & $5{,}002$ & $+6.9$  \\
Spain\(^{+}\)           & $14.5$ & $-118.5$ & $2{,}746$ & $+11.4$ \\
Austria                 & $15.0$ & $-69.6$  & $2{,}435$ & $+10.9$ \\
Costa Rica              & $16.0$ & $-52.9$  & $1{,}784$ & $+14.0$ \\
Germany                 & $16.5$ & $-63.5$  & $1{,}962$ & $+5.8$  \\
United Kingdom\(^{\#}\) & $17.0$ & $-88.8$  & $3{,}320$ & $+11.1$ \\
T\"urkiye\(^{+}\)       & $17.5$ & $+22.3$  & $1{,}226$ & $-0.0$  \\
Czechia                 & $18.5$ & $-104.7$ & $3{,}988$ & $+12.0$ \\
Latvia                  & $19.0$ & $-50.7$  & $3{,}229$ & $+13.6$ \\
Slovak Republic         & $19.0$ & $-55.0$  & $3{,}810$ & $+11.2$ \\
Estonia\(^{+}\)         & $21.0$ & $-65.9$  & $2{,}150$ & $+10.2$ \\
Lithuania\(^{+}\)       & $22.0$ & $-39.5$  & $3{,}376$ & $+5.0$  \\
Ireland\(^{\#}\)        & $23.0$ & $+78.9$  & $1{,}668$ & $-1.9$  \\
Slovenia                & $25.0$ & $-78.3$  & $3{,}347$ & $+10.1$ \\
Greece                  & $28.5$ & $-54.9$  & $3{,}240$ & $+7.8$  \\
Korea\(^{*}\)           & $28.5$ & $-36.2$  & $620$     & $+5.5$  \\
Poland                  & $28.5$ & $-64.8$  & $3{,}186$ & $+6.1$  \\
Canada                  & $29.0$ & $-51.9$  & $1{,}304$ & $+13.9$ \\
Israel\(^{*}\)          & $30.5$ & $-25.1$  & $1{,}300$ & $+8.7$  \\
Chile\(^{\#}\)          & $31.5$ & $-37.5$  & $3{,}272$ & $+2.5$  \\
Denmark\(^{*\#}\)       & $32.0$ & $-28.5$  & $1{,}344$ & $-3.1$  \\
Australia\(^{*\#}\)     & $33.5$ & $-26.1$  & $661$     & $+3.7$  \\
Japan\(^{*}\)           & $33.5$ & $-46.6$  & $456$     & $+7.9$  \\
New Zealand\(^{*\#}\)   & $34.5$ & $-45.1$  & $460$     & $+21.2$ \\
United States\(^{*}\)   & $35.0$ & $-26.3$  & $3{,}307$ & $+19.3$ \\
\midrule
Italy\(^{\dagger}\)     & ---    & $-50.6$  & $3{,}379$ & $+6.1$  \\
\bottomrule
\end{tabular}
\vspace{0.15cm}
\begin{minipage}{0.95\textwidth}
\scriptsize \textit{Notes:} Countries in ascending order of frontier distance
(Table~\ref{tab:country_ranking_full}); outcomes are the \emph{raw} realized
data over Q1.2020--Q4.2022---cumulative output gap from the estimation panel,
cumulative excess mortality from the weekly mortality data, and the change in
the debt-to-GDP ratio from reported levels---not the model-implied values on
which the frontier thresholds are based (closure gaps are documented in
Table~\ref{tab:country_composition}). The table is not an outcome ranking:
frontier distance measures conversion efficiency at the country's own
realized position, so frontier-near countries need not have favorable
outcomes (Mexico) and frontier-far countries need not have unfavorable ones
(Denmark). What the raw data do display is the efficiency pattern itself:
Sweden and the United States realized nearly identical output losses, but
Sweden with fewer excess deaths and one tenth of the debt increase; Norway
and Switzerland reached small losses at low fiscal cost, while the furthest
group paid the largest debt increases of the panel (United States \(+19.3\),
New Zealand \(+21.2\)~pp) for moderate output outcomes. T\"urkiye and Ireland
recorded positive cumulative output gaps (growth over the window); flags as
in Table~\ref{tab:country_ranking_full}.
\end{minipage}
\end{table}



\begin{table}[H]
\centering
\scriptsize
\caption{Observed policy composition, realized fiscal exchange rate, and fit diagnostics.}
\label{tab:country_composition}
\begin{tabular}{l ccc !{\vrule width 0.4pt} cc !{\vrule width 0.4pt} c}
\toprule
 & \multicolumn{3}{c}{\emph{Observed composition}} & \multicolumn{2}{c}{\emph{Counterfactuals}} & \emph{Diagnostics} \\
\cmidrule(lr){2-4} \cmidrule(lr){5-6}
 & \textbf{Liquidity} & \textbf{DI} & \textbf{Front-} & \textbf{Fiscal} & \textbf{Deaths} & \textbf{Fit} \\
 & \textbf{share} & \textbf{share} & \textbf{loading} & \textbf{exch.\ rate} & \textbf{avoided /M} & \textbf{score} \\
\midrule
Mexico                  & $0.12$ & $0.62$ & $1.00$ & $-0.14$ & $9{,}027$  & $1.7$ \\
Sweden                  & $0.34$ & $0.13$ & $0.96$ & $3.12$  & $7{,}065$  & $0.8$ \\
Portugal                & $0.40$ & $0.19$ & $0.82$ & $2.81$  & $6{,}342$  & $1.2$ \\
Finland                 & $0.32$ & $0.33$ & $0.84$ & $2.95$  & $1{,}018$  & $1.1$ \\
France                  & $0.33$ & $0.10$ & $0.61$ & $3.74$  & $10{,}289$ & $0.8$ \\
Norway                  & $0.51$ & $0.03$ & $0.83$ & $3.74$  & $1{,}183$  & $1.7$ \\
Iceland\(^{+}\)         & $0.11$ & $0.38$ & $0.46$ & $1.71$  & $1{,}834$  & $4.4$ \\
Belgium\(^{+}\)         & $0.28$ & $0.17$ & $0.88$ & $2.94$  & $17{,}181$ & $3.0$ \\
Luxembourg              & $0.24$ & $0.02$ & $0.69$ & $2.93$  & $3{,}594$  & $0.6$ \\
Switzerland\(^{+}\)     & $0.24$ & $0.03$ & $0.80$ & $3.57$  & $7{,}361$  & $4.4$ \\
Colombia\(^{\#+}\)      & $0.19$ & $0.56$ & $0.73$ & $-0.02$ & $4{,}906$  & $3.0$ \\
Netherlands\(^{+}\)     & $0.14$ & $0.15$ & $0.72$ & $2.17$  & $12{,}006$ & $5.4$ \\
Hungary                 & $0.26$ & $0.24$ & $0.53$ & $2.29$  & $5{,}453$  & $1.6$ \\
Spain\(^{+}\)           & $0.46$ & $0.18$ & $0.74$ & $2.82$  & $24{,}359$ & $4.5$ \\
Austria                 & $0.07$ & $0.13$ & $0.66$ & $2.42$  & $4{,}826$  & $1.3$ \\
Costa Rica              & $0.66$ & $0.27$ & $0.62$ & $1.88$  & $1{,}519$  & $2.1$ \\
Germany                 & $0.39$ & $0.16$ & $0.82$ & $3.62$  & $3{,}443$  & $1.0$ \\
United Kingdom\(^{\#}\) & $0.24$ & $0.13$ & $0.84$ & $2.92$  & $12{,}511$ & $1.2$ \\
T\"urkiye\(^{+}\)       & $0.41$ & $0.10$ & $0.86$ & $4.09$  & $2{,}664$  & $2.7$ \\
Czechia                 & $0.34$ & $0.22$ & $0.66$ & $2.96$  & $6{,}377$  & $1.2$ \\
Latvia                  & $0.31$ & $0.03$ & $0.75$ & $3.65$  & $3{,}118$  & $1.1$ \\
Slovak Republic         & $0.22$ & $0.02$ & $0.63$ & $2.72$  & $5{,}319$  & $1.1$ \\
Estonia\(^{+}\)         & $0.26$ & $0.04$ & $0.86$ & $3.52$  & $1{,}629$  & $2.6$ \\
Lithuania\(^{+}\)       & $0.12$ & $0.38$ & $0.78$ & $1.40$  & $8{,}014$  & $3.2$ \\
Ireland\(^{\#}\)        & $0.10$ & $0.42$ & $0.62$ & $0.30$  & $5{,}821$  & $0.7$ \\
Slovenia                & $0.11$ & $0.05$ & $0.64$ & $2.86$  & $5{,}722$  & $1.3$ \\
Greece                  & $0.28$ & $0.18$ & $0.58$ & $2.08$  & $2{,}913$  & $1.8$ \\
Korea\(^{*}\)           & $0.43$ & $0.29$ & $0.81$ & $2.40$  & $348$      & $0.9$ \\
Poland                  & $0.10$ & $0.20$ & $0.77$ & $1.77$  & $5{,}368$  & $1.7$ \\
Canada                  & $0.10$ & $0.50$ & $0.74$ & $0.90$  & $2{,}020$  & $0.8$ \\
Israel\(^{*}\)          & $0.12$ & $0.58$ & $0.92$ & $-0.50$ & $1{,}858$  & $1.2$ \\
Chile\(^{\#}\)          & $0.18$ & $0.67$ & $0.40$ & $-0.18$ & $4{,}250$  & $2.6$ \\
Denmark\(^{*\#}\)       & $0.46$ & $0.25$ & $0.97$ & $2.81$  & $923$      & $2.1$ \\
Australia\(^{*\#}\)     & $0.05$ & $0.27$ & $0.71$ & $1.76$  & $796$      & $1.0$ \\
Japan\(^{*}\)           & $0.52$ & $0.26$ & $0.70$ & $3.30$  & $230$      & $2.6$ \\
New Zealand\(^{*\#}\)   & $0.08$ & $0.30$ & $0.83$ & $2.32$  & $242$      & $0.5$ \\
United States\(^{*}\)   & $0.14$ & $0.51$ & $0.54$ & $1.12$  & $6{,}036$  & $0.3$ \\
\midrule
Italy\(^{\dagger}\)     & $0.41$ & $0.06$ & $0.74$ & $3.88$  & $20{,}994$ & $2.3$ \\
\bottomrule
\end{tabular}
\vspace{0.15cm}
\begin{minipage}{0.95\textwidth}
\scriptsize \textit{Notes:} Countries in the order of Table~\ref{tab:country_ranking_full}
(ascending frontier distance). Liquidity and DI shares of the total observed
fiscal package (take-up-adjusted effective terms); front-loading is the share
committed in Q1--Q3.2020. The fiscal exchange rate is the cumulative output
gain per pp of fiscally induced debt against a counterfactual that keeps the
observed containment path and sets fiscal instruments to zero; deaths avoided
(per million) are computed against a full-zero counterfactual (no containment,
no fiscal support) and thus reflect the entire policy package, containment
included. Fit score as defined in Appendix~\ref{app:country_fit} (higher =
worse fit); by the invariance result of Appendix~\ref{app:ifr_invariance},
mortality-level misfit does not affect the frontier distances.
\(^{*}\)~output ceiling; \(^{\#}\)~instrument committed beyond the common
menu; \(^{+}\)~worst-fitting quartile; \(^{\dagger}\)~Italy: not ranked (see
Table~\ref{tab:country_ranking_full}); its composition, exchange rate, and
counterfactuals remain well-defined.
\end{minipage}
\end{table}

\begin{figure}[H]
    \centering
     \captionsetup{font=scriptsize}
    \includegraphics[width=0.95\textwidth]{Files/text/main/fig_countries_frontier_bars_aerstyle.pdf}
    \caption{Output improvement at the country-specific frontier, by country.
    \textit{Notes:} Improvement in the discounted output loss
    \(\mathcal{L}_y\) at the country's own realized mortality and terminal
    debt, in percent, sorted ascending. The cluster at 100 percent marks the
    output-ceiling countries, for which the frontier eliminates the
    policy-amenable output loss entirely within the estimated linear
    transition technology. Italy omitted (Appendix~\ref{app:country_special}).}
    \label{fig:countries_bars}
\end{figure}

\begin{figure}[H]
    \centering
     \captionsetup{font=scriptsize}
    \includegraphics[width=0.95\textwidth]{Files/text/main/fig_countries_correlates_aerstyle.pdf}
    \caption{Frontier distance and observed policy characteristics.
    \textit{Notes:} Output improvement at the own frontier against initial
    debt, the observed liquidity share, and the observed front-loading share;
    solid lines are linear fits. The linear cross-sectional associations are
    weak---attenuated by the output-ceiling countries and the conditional
    nature of the distance measure---which is why the main text reports the
    composition--distance link as a group pattern
    (Section~\ref{subsec:results_countries}) rather than a regression.}
    \label{fig:countries_correlates}
\end{figure}


\subsection{Full Frontier Plot}


\begin{figure}[H]
    \centering
     \captionsetup{font=scriptsize}
    \includegraphics[width=0.7\textwidth]{Files/text/fig_countries_scatter_AC_aerstyle.pdf}
    \caption{Distance to the country-specific policy frontier.}
    \label{fig:countries_scatterp}
\end{figure}


\end{document}